\documentclass[12pt,a4paper,]{report}
\usepackage[T1]{fontenc}
\usepackage[utf8]{inputenc}
\usepackage[english,italian]{babel}
%-----------------------------------------------
%margini foglio
%-----------------------------------------------
\usepackage[a4paper, total={6.5in, 9.5in}]{geometry}
%-----------------------------------------------
%pacchetti
%-----------------------------------------------
\usepackage{hycolor}
\usepackage{xcolor}
\usepackage[colorlinks=true, linkcolor=black, urlcolor=blue]{hyperref}

\usepackage{graphicx}
\usepackage{float}


\usepackage{amsthm}
\usepackage{amsmath}
\usepackage{thmtools}
\usepackage{listings} 

\usepackage{wrapfig}
\usepackage{xcolor}

\usepackage{amsmath,amssymb}
\usepackage{latexsym}

\usepackage{pdfpages}

\usepackage{multirow}

\usepackage{array} 
% \usepackage{parskip} %rimuove indentazione paragraph e altro
% usare \\ per andare a capo senza spazi, altrimenti le righe vuote vengono
% interpretate come paragrafi

%permette underiline con linebreak, l'opzione serve ad evitare che rimpiazzi italic come default
\usepackage[normalem]{ulem}

%pacchetto per la tripla immagine in linea
\usepackage{subcaption}
%-----------------------------------------------

%teoremi
%-----------------------------------------------
\declaretheoremstyle[
headformat=\NAME\ \NUMBER:\NOTE,
headfont=\bfseries\itshape,
notefont=\normalfont\bfseries, %\large se si vuole piu' grande
notebraces={}{},
headpunct=,
bodyfont=\normalfont,
postheadspace=\newline,
qed=\qedsymbol
]{mystyle}
\declaretheorem[style=mystyle]{teorema}

\declaretheorem[style=mystyle]{metodo}

\declaretheoremstyle[
headformat=\NAME\ \NUMBER:\NOTE,
headfont=\small,
notefont=\normalfont,
notebraces={}{},
headpunct=,
bodyfont=\normalfont,
postheadspace=\newline,
% qed=\qedsymbol
]{lessonstyle}
\declaretheorem[style=lessonstyle]{lezione}

\declaretheoremstyle[
% headformat=\NAME\NOTE,
headfont=\bfseries\itshape,
notefont=\normalfont\bfseries,
notebraces={- }{},
headpunct=,
bodyfont=\normalfont,
postheadspace=\newline,
qed=\qedsymbol
]{defstyle}
\declaretheorem[style=defstyle]{definizione}


% \declaretheorem[style=mystyle]{algoritmo}

\declaretheorem[style=mystyle, numberwithin=part]{algoritmo}

\newcounter{utilita}
\newenvironment{utilita}[1][]{%
    \par\addvspace{\medskipamount}
    \refstepcounter{utilita}%
    \noindent{\bfseries\itshape Utility MATLAB \theutilita: #1}%
    \par\bigskip\noindent
}{%
    \par\nobreak\hfill$\square$%
    \par\addvspace{\medskipamount}
}


\declaretheoremstyle[
notebraces={: }{},
headpunct=:,
headformat=\NAME\NOTE,
bodyfont=\normalfont\itshape
]{notastyle}
\declaretheorem[style=notastyle]{nota}

\declaretheoremstyle[                                                 
headformat=\NAME\ Teorema \theteorema\NOTE,
headfont=\bfseries\itshape,
notefont=\normalfont\bfseries,
notebraces={\ (}{)},                                                  
headpunct=,         
bodyfont=\normalfont,                                                 
postheadspace=\newline,
qed=\qedsymbol                                                        
]{dimostrazionestyle}
\declaretheorem[style=dimostrazionestyle,
  numbered=no]{dimostrazione}

\declaretheoremstyle[
headformat=\NAME\NOTE,
headfont=\bfseries\itshape,
notefont=\normalfont\bfseries,
notebraces={\ (}{)},
headpunct=,
bodyfont=\normalfont,
postheadspace=\newline,
qed=\qedsymbol
]{dimostrazioneliberastyle}
\declaretheorem[style=dimostrazioneliberastyle,
  name=Dimostrazione,
  numbered=no]{dimostrazionelibera}

%----------------------------------------------

%codice
%---------------------------------------------
\lstset{language=Matlab}

\definecolor{backcolour}{RGB}{68 71 90}
\definecolor{parole}{RGB}{248 248 242}
\definecolor{commento}{RGB}{98 114 164}

\lstdefinestyle{mystyleold}{
% backgroundcolor=\color{backcolour},
basicstyle=\ttfamily\small\color{parole},
keywordstyle=\color{orange},
commentstyle=\color{commento}
}

\lstdefinestyle{mystyle}{
basicstyle=\ttfamily\small,
keywordstyle=\color{black}\bfseries,
commentstyle=\color{commento},
% numbers=left,
frame=lines,
breaklines=true,
breakatwhitespace=false,
% showspaces=false,
showstringspaces=false
}


\lstset{style=mystyle}
%---------------------------------------------

%comandi
%-----------------------------------------------
\newcommand{\s}{\hspace{1mm}}
\newcommand{\R}{\mathbb{R}}

% Imposta la distanza tra il testo e il bordo delle celle
\renewcommand{\arraystretch}{1.7}
%-----------------------------------------------


%mostra fino a ssection nell'indice
%-----------------------------------------------
\setcounter{tocdepth}{2}    
\setcounter{secnumdepth}{3}

%documento
%-----------------------------------------------
\title{Calcolo Numerico}
\author{Francesco Bonistalli}
\date{January 2025}

\begin{document}

\maketitle

\tableofcontents

% \listoftheorems[ignoreall, show={teorema}]

\chapter{Rappresentazione dei numeri e analisi dell'errore}
Quaderno 1

\begin{lezione}
    [Lezione 1]
\end{lezione}

\section{Numeri in virgola mobile}

Scelto un qualunque numero intero $\beta > 1$, ogni numero non nullo $x \in \mathbb{R}$ ammette una rappresentazione in base $\beta$:

\[
x = sign(x)\beta^e \sum_{j=1}^{\infty} \alpha_j\beta^{-j}
\]
\noindent
In cui:
\begin{itemize}
  \item $\beta$ : Base on $\beta \in \mathbb{N}$, $\beta > 1$
  \item $e$ : Esponente con $e \in \mathbb{Z}$
  \item $\alpha$ : Cifre con $\alpha_j \in \{0,\ldots,\beta - 1\}$
  \item $\displaystyle\sum_{j=1}^{\infty} \alpha_j\beta^{-j}$ = Mantissa
  \item $sign(x) = \begin{cases} 1 & x > 0 \\ -1 & x < 0 \end{cases}$
\end{itemize}
Vale il seguente

\begin{teorema}
    [Teorema di Rappresentazione]
    Data una base intera $\beta > 1$ e un qualsiasi numero reale $x$ diverso da zero, esiste un'unica rappresentazione in base $\beta$ tale che:
    
    1) Sia $\alpha_1 \neq 0$
    
    2) Non vi sia un intero $k$ per cui si abbia $\alpha_j = \beta - 1, \forall j > k$
    
    \noindent
    In simboli:\\
    Sia $\beta \in \mathbb{N}$, $\beta > 1$, $x \in \mathbb{R}\backslash\{0\}$ allora:
    
    \[
    \exists\, x = sign(x)\beta^e \sum_{j=1}^{\infty} \alpha_j\beta^{-j}
    \text{ se: } \alpha_1 \neq 0 \wedge \nexists k \in \mathbb{N} : \alpha_j = \beta - 1, \forall j > k
    \]
    \noindent
    Quindi: $\displaystyle\frac{1}{\beta} \leq \sum_{j=1}^{\infty} \alpha_j\beta^{-j} < 1$
        
\end{teorema}

\begin{definizione}
    [Numeri di macchina]
    Dati $\beta, m, L, U$ fanno parte dell'insieme dei numeri di macchina:
    \[
    x \in \mathbb{R} = sign(x) \cdot \beta^e \sum^m_{j=1} \alpha_j\beta^{-j} 
    \]
    con:
    \begin{itemize}
        \item $L\le e \le U$
        \item $\alpha_i \in\left\{0,\dots,\beta-1\right\}$
        \item $\alpha_1 \ne 0$ (escluso per lo 0)
    \end{itemize}
\end{definizione}
\noindent Osservazioni:
\begin{itemize}
    \item Cardinalità dell'insieme $=|F(\beta, m, L, U)| = 2 \cdot (U - L + 1) \cdot (\beta^m - \beta^{m-1}) + 1$
    \item Possibili segni $= 2$
    \item Possibili esponenti $= (U - L + 1)$
    \item Possibili cifre $= (\beta^m - \beta^{m-1})$
    \item Lo zero $= 1$

  \item Qual'è il \textbf{numero più grande} in $F(\ldots)$?
        \[
        \beta^U \cdot (\beta - 1) \sum_{j=1}^{m} \beta^{-j} = \boxed{\beta^U(1 - \beta^{-m})}
        \]
  
  \item Qual'è il \textbf{più piccolo numero} positivo in $F(\ldots)$?
        \[
        \beta^L \cdot \beta^{-1} = \boxed{\beta^{L-1}}
        \]
\end{itemize}
Facendo variare la mantissa e tenendo fissi esponente e segno, si ottengono tutti i numeri di macchina appartenenti all'intervallo $[\beta^{e-1}, \beta^e]$.
\\
I valori possibili per la mantissa corrispondono a tutti i modi in cui si possono scegliere le cifre $\alpha_j$. In generale si hanno $m$ cifre libere, poiché $\alpha_1 \neq 0$ (per il teorema visto sopra); per $\beta = 2$ questo vincolo fissa $\alpha_1 = 1$, lasciando $m-1$ cifre binarie libere. Ad esempio con $m = 3$ restano 2 cifre, quindi $2^2 = 4$ combinazioni possibili.
\\
Ogni intervallo $[\beta^{e-1}, \beta^e]$ contiene dunque lo stesso numero di valori (4 nell'esempio), equispaziati al suo interno; passando all'intervallo successivo lo spazio tra i valori raddoppia, poiché la distanza tra due numeri di macchina consecutivi scala proporzionalmente a $\beta^e$.
\\
Siano $x_1, x_2 \in F \cap [\beta^i, \beta^{i+1}]$ adiacenti. La loro distanza è
\[
|x_1 - x_2| = \beta^{i+1} \cdot \beta^{-m} = \beta^{i+1-m}
\]

\subsection{Rappresentazione in precisione doppia}
Con \uline{precisione doppia} (\textit{double precision}) si intende la rappresentazione su \texttt{64 bit}.\\
Il nome deriva dalla precedente nomenclatura data alla rappresentazione su \texttt{32 bit} chiamata \textit{single precision}.
\[
\begin{array}{c c c c c}
\beta & m & L & U & \\ \hline
2 & 53 & -1021 & 1024
\end{array}
\quad
64\ \text{bit} = 8\ \text{byte}
\]
\begin{nota}
    [Struttura IEEE 754 e bit implicito]
    I 64 bit sono ripartiti come segue:
    \[
    \begin{array}{l c}
    \text{Segno} & 1\ \text{bit} \\
    \text{Esponente} & 11\ \text{bit} \\
    \text{Mantissa (campo memorizzato)} & 52\ \text{bit} \\ \hline
    \text{Totale} & 64\ \text{bit}
    \end{array}
    \]
    Per i numeri \textbf{normalizzati} il primo bit della mantissa non viene memorizzato: è un \textbf{bit implicito} ($\alpha_1 = 1$ sempre, in base 2).\\
    In base 2, infatti un numero normalizzato ha sempre la forma "$1.\,qualcosa \times 2^e$". Il primo bit dopo la virgola è quindi sempre 1: non può essere 0, altrimenti potremmo ridurre l'esponente e rinormalizzare. Siccome è sempre 1, è ridondante memorizzarlo: lo sappiamo già.\\
    La \textbf{precisione effettiva} è quindi $1 + 52 = 53$ bit (da cui $m = 53$ nella tabella), ma il campo mantissa nella rappresentazione occupa solo \textbf{52 bit}.
\end{nota}

\subsection{Approssimazioni}
Dato un numero reale, ci poniamo ora il problema di trovare un numero di macchina all'interno dell'insieme F che lo approssimi.\\
Sia $x \in \mathbb{R}\backslash F$, si possono verificare 3 casi:
\begin{enumerate}
  \item $|x| > \beta^U(1 - \beta^{-m})$ $\to$ Overflow: $|x|$ è maggiore del massimo numero rappresentabile.
  \item $|x| < \beta^{L - 1}$ $\to$ Underflow: $|x|$ è minore del minimo numero rappresentabile.
  \item $\beta^{L-1} < |x| < \beta^U(1 - \beta^{-m})$ $\to$ si utilizza una funzione di arrotondamento per scegliere con quale valore approssimarlo.
\end{enumerate}

\begin{definizione}
    [Round down (Troncamento)]
    Taglia la somma alle prime $m$ cifre, scartando tutto il resto:
    \[
    Tr(x)
    = \operatorname{sign}(x)\,\beta^{e}
    \sum_{j=1}^{m} \alpha_j \beta^{-j}
    \]
\end{definizione}
\begin{definizione}
    [Round up]
    Si prende il troncamento e si aggiunge un'unità nell'ultima cifra rappresentabile:
    \[
    Up(x)
    = \operatorname{sign}(x)\,\beta^{e}
    \left[
    \left(
        \sum_{j=1}^{m} \alpha_j \beta^{-j}
    \right)
    + \beta^{-m}
    \right]
    \]
\end{definizione}
\begin{definizione}
    [Round to nearest]
    \label{def: round to nearest}
    \[
    Rn(x)=
    \begin{cases}
    Tr(x) & \text{se } \alpha_{m+1} < \dfrac{\beta}{2},\\[6pt]
    Up(x) & \text{se } \alpha_{m+1} \ge \dfrac{\beta}{2}.
    \end{cases}
    \]
\end{definizione}

\subsubsection{Overflow e Underflow}
\begin{definizione}
    [Overflow]
    Se $ |x| > \beta^{U}\bigl(1-\beta^{-m}\bigr) $
    allora si ha \textbf{overflow}, cio\`e valore maggiore del massimo rappresentabile (si utilizza una rappresentazione ad hoc).
\end{definizione}
\begin{definizione}
    [Underflow]
    Se $|x| < \beta^{L-1}$ allora si ha \textbf{underflow}, cio\`e valore minore del minimo rappresentabile
    (in genere associato a $0$).
\end{definizione}

\paragraph{Osservazione}
Il round to nearest (\autoref{def: round to nearest}) permette di scegliere ogni volta il numero di macchina che minimizza l'errore assoluto:
\[
|x - Rn(x)| \leq \min |x - z| = \frac{1}{2}\beta^{e-m} \text{ con } z \in F(\beta, m, L, U)
\]
Cioè la metà della distanza tra due numeri macchina.\\
\textbf{Spiegazione}: come abbiamo visto due numeri di macchina consecutivi differiscono solo per l'ultima cifra della mantissa, la quale è pesata per $\beta^{-m}$.
Ponendosi quindi a livello di un qualsiasi esponente $e$  questa distanza viene scalata per tale $e$: $\beta^{e-m}$.\\
Quello che concludiamo da questa situazione è che, prendendo un numero reale che, nel peggiore dei casi, si trova perfettamente a metà tra due numeri di macchina, allora questo disterà al massimo $\frac{1}{2}\beta^{e-m}$ dal numero di macchina a cui verrà approssimato.

\begin{definizione}
    [Errore assoluto nell'arrotondamento]
    \[
    \delta_x := |x - Rn(x)|
    \]
\end{definizione}
\begin{definizione}
    [Errore relativo nell'arrotondamento]
    \[
    \varepsilon_x := \frac{|x - Rn(x)|}{|x|} = \frac{\delta x}{|x|}, \,x \neq 0
    \]  
\end{definizione}
\noindent
Questa quantità può essere maggiorata grazie al fatto che abbiamo appena visto il limite superiore del numeratore e che, al denominatore, abbiamo 
\[
    |x| = \beta^e (\alpha_1\beta^{-1}+\alpha_2\beta^{-2}+\dots) \ge \beta^e (\alpha_1\beta^{-1}) \underbrace{=}_{\alpha_1 = 1}\beta^e \beta^{-1}
\]
L'errore relativo dell'approssimazione di un numero reale quindi è:
\[
\boxed{|\varepsilon_x| \leq \frac{\frac{1}{2}\beta^{e-m}}{\beta^e \cdot \beta^{-1}} = \frac{1}{2}\beta^{-m+1}} \xrightarrow{\beta=2} 2^{-m}
\]
da cui notiamo che \uline{l'errore relativo non dipende da $e$}, dunque non dipende dall'intervallo in cui ci troviamo. Dunque \uline{$\varepsilon_x$ è uniformemente limitato dall'alto da $U := \frac{1}{2}\beta^{-m+1}$, quantità chiamata Precisione di Macchina}.

\begin{definizione}
    [Precisione di macchina]
    \label{def: prec di macchina}
    La quantità $U= \frac{1}{2}\beta^{1-m}$ è detta precisione di macchina.
\end{definizione}
\noindent
Si ha:
\begin{itemize}
    \item nel caso di precisione doppia: $U \approx 10^{-16}$
    \item nel caso di precisione singola: $U \approx 10^{-8}$
\end{itemize}

\begin{table}[H]
\centering
\renewcommand{\arraystretch}{1.5} % Aumenta lo spazio verticale tra le righe
\begin{tabular}{|l|c|c|c|c|}
\hline
\textbf{FORMATO} & \textbf{SEGNO} & \textbf{ESPONENTE} & \textbf{MANTISSA} & \textbf{U} \\ \hline
SINGOLA          & 1              & 8                  & 23                & $10^{-8}$  \\ \hline
QUADRUPLA        & 1              & 15                 & 112               & $10^{-32}$ \\ \hline
MEZZA            & 1              & 5                  & 10                & $10^{-4}$  \\ \hline
BFLOAT           & 1              & 8                  & 7                 & $10^{-3}$  \\ \hline
\end{tabular}
\caption{Numero di bit}
\end{table}

\section{Numeri Sottonormalizzati (Subnormal Numbers)}
Sono quei numeri più piccoli di $\beta^{L-1}$.\\
Se ho $x \in \left[ 0, \beta^{L-1}\right]$ possiamo rinunciare all'assunzione $\alpha_1=1$ e usare un valore speciale dell'esponente che indica $\beta^L$ e $\alpha_1 = 0$.
Così facendo riesco ad ottenere una rappresentazione del tipo $\beta^L\cdot 0,0101001$ che quindi riguarda valori compresi in $\left[ \beta^{L-m}, \beta^{L-1}\right)$.
\begin{figure}[H]
    \centering
    \includegraphics[width=0.5\linewidth]{img/subnormal numbers.png}
    \caption{I numeri sottonormali sono quelli in blu.}
\end{figure}
\noindent
Quello che si fa è infatti fissare $L$ e di conseguenza, variando $m$ da $1$ al suo valore massimo, ottenere tutti i numeri sottonormalizzati per quei valori di esponente minimo e mantissa.\\
Dunque nella rappresentazione numerica si guadagnano i numeri compresi tra $\left[ 0, \beta^{L-1}\right]$ ma essi \textbf{non} avranno la \textbf{stessa accuratezza relativa}, in particolare \textbf{più mi avvicino a 0 più perdo cifre significative}.\\
\paragraph{Osservazione} 
\uline{Nei numeri sottonormali la spaziatura tra numeri consecutivi è costante}
e pari a $\beta^{L-m}$, indipendentemente da $|x|$. Di conseguenza
l'errore relativo $|\varepsilon_x| \leq \frac{\frac{1}{2}\beta^{L-m}}{|x|}$
cresce al diminuire di $|x|$: non è più uniformemente limitato da $U$ (precisione di macchina).
All'estremo destro dell'intervallo $[\beta^{L-m}, \beta^{L-1})$ si ha
$|\varepsilon_x| \approx U$, all'estremo sinistro $|\varepsilon_x| \approx 1$.

\begin{definizione}[Numeri di macchina estesi]
    $F(\beta,m,L,U) \cup \{\text{numeri sottonormali}\}$
    è detto insieme dei \textbf{numeri di macchina estesi}.
    Essi consentono un \emph{underflow graduale}, ma i numeri sottonormali
    hanno in generale errore relativo superiore a $U$.
\end{definizione}

\begin{lezione}
    [Lezione 2 01/03]
\end{lezione}

\section{Operazioni di macchina}

Un'operazione di macchina (o operazione elementare) è un'operazione aritmetica elementare $(+, -, \times, /)$ eseguita in aritmetica floating point con precisione finita, anziché in aritmetica esatta.\\
Si indica con il simbolo $\oplus, \ominus, \otimes, \oslash$ per distinguerla dall'operazione matematica esatta:
\begin{itemize}
    \item $x \oplus y$ = operazione di macchina (somma in floating point)
    \item $x + y$ = operazione esatta (somma matematica)
\end{itemize}

\begin{definizione}
    [operazione di macchina]
    \label{def: op di macchina}
    Data un'operazione aritmetica $\mathrm{op} \in \{+,-,\times,/\}$, la corrispondente operazione di macchina è definita come:
    \[
        x \;\widetilde{\mathrm{op}}\; y := RN(x \;\mathrm{op}\; y)
    \]
    ovvero il risultato esatto arrotondato al numero di macchina più vicino.
    Si indica con i simboli $\oplus, \ominus, \otimes, \oslash$.
\end{definizione}
\begin{nota}
    L'addizione "floating point" \uline{non} gode della proprietà associativa, mentre la moltiplicazione non gode della proprietà distributiva rispetto all'addizione.
\end{nota}
\paragraph{Problema}
Quando si sommano o sottraggono due numeri di macchina dello stesso segno che hanno lo stesso esponente e con le mantisse che differiscono di poco (quindi con le prime cifre uguali), si incorre in una perdita di cifre significative nel risultato.\\
\uline{Tale fenomeno produce una notevole amplificazione degli errori relativi.}
\begin{nota}
    (Non si verifica nella moltiplicazione e nella divisione).
\end{nota}

\begin{definizione}
    [Cancellazione]
    Fenomeno che si verifica quando si sottraggono due numeri di macchina dello stesso segno con lo stesso esponente e mantisse che differiscono di poco, producendo una notevole amplificazione degli errori relativi.
\end{definizione}
Le cause di tale inconveniente \textbf{non} sono da attribuire all'operazione di \uline{sottrazione} effettuata dalla macchina, \textbf{ma} dipendono dagli errori già presenti nei due operandi e \uline{da quelli introdotti dal processo di approssimazione effettuato dal calcolatore per poter memorizzare i numeri in memoria}.

\paragraph{Esempio in base 10}
Lavoriamo in un sistema a \textbf{5 cifre significative} decimali per vedere il problema chiaramente.
\[
x = 1.2348, \quad y = 1.2346
\]
Entrambi rappresentati esattamente. Il risultato esatto è:
\[
x - y = 0.0002
\]
Fin qui nessun problema. Ma supponiamo che i valori reali da cui derivano $x$ e $y$ fossero:
\[
x_{\text{vero}} = 1.23481, \quad y_{\text{vero}} = 1.23457
\]
Troncati a 5 cifre diventano $x = 1.2348$ e $y = 1.2346$: errori assoluti entrambi dell'ordine di $5 \times 10^{-5}$.
\\
Il risultato vero sarebbe $x_{\text{vero}} - y_{\text{vero}} = 0.00024$, ma noi calcoliamo $0.0002$.
\[
\varepsilon_{\text{risultato}} = \frac{|0.00024 - 0.0002|}{0.00024} \approx 17\%
\]
Gli errori originali erano dell'ordine di $\approx 10^{-5}$ --- il risultato ha un errore del $17\%$, ovvero abbiamo avuto una amplificazione di circa $10^4$.


\section{Errore nel calcolo di una funzione}
Data $f:\mathbb{R}^n\to\mathbb{R}$ definita come:
\[
f(P_0)\quad , \quad P_0=
\begin{pmatrix}
    x_1\\
    \vdots\\
    x_n
\end{pmatrix}
\]
3 fonti di errore:
\begin{enumerate}
    \item Dobbiamo approssimare $f$ con una funzione che usa solo $[+, -, \times, /]$. Tali funzioni si dicono \textbf{razionali}, ovvero si possono scrivere o come polinomi o come rapporto tra polinomi: $\frac{p(x)}{q(x)}$.\\
    Si indica con $\tilde{f}$ la \textbf{rappresentazione razionale} di $f$ (\textbf{errore analitico}).
    \item $\tilde{f}(x)$ è tradotta in un algoritmo che usa le operazioni di macchina (\autoref{def: op di macchina}), per cui quello che si valuta veramente è $\tilde{f}_a(x)$ (\textbf{errore algoritmico}).
    \item Al posto di $P_0$ si valuta $P_1 = RN(P_0)$ (\textbf{errore inerente}).
\end{enumerate}
La \textbf{conclusione} quindi è che vorremmo $f(P_0)$ ma in realtà calcoliamo $\tilde{f}_a(P_1)$.
\begin{definizione}
    [Errore totale nel calcolo di funzione]
    $ $
    \begin{itemize}
        \item Errore totale assoluto:
        $$\delta_f = f(P_0) - \tilde{f}_a(P_1)$$
        \item Errore totale relativo:
        \[
        \varepsilon_f = \frac{\delta_f}{f(P_0)}
        \]
    \end{itemize}
\end{definizione}
\section{Funzioni razionali}
Se $f$ è una funzione razionale, ovvero calcolabile esattamente tramite sole operazioni di $[+,\,-,\,\times,\,/]$, allora $\tilde{f} = f$ e l'\textbf{errore analitico è nullo}. La decomposizione dell'\textbf{errore totale} si riduce \textbf{perciò} a \textbf{due sole componenti}:
\[
    f(P_0) - \tilde{f}_a(P_1)
    = \underbrace{f(P_0) - f(P_1)}_{\delta_d\;\textbf{(inerente)}}
    + \underbrace{f(P_1) - \tilde{f}_a(P_1)}_{\delta_a\;\textbf{(algoritmico)}}
\]
dove $P_0 = [x_1^{(0)},\dots,x_n^{(0)}]$ sono i dati esatti e
$P_1 = \mathrm{RN}(P_0)$ i loro arrotondamenti in macchina.
Le sezioni seguenti stimano ciascuna delle due componenti.

\subsection{Stimare l'errore inerente}
\begin{definizione}
    [Errore inerente assoluto]
    Errore dovuto al fatto che al posto di $f(P_0)$ valutiamo $\tilde{f}(P_1)$:
    \[
    \delta_d = f(P_0) - f(P_1)
    \]
    Dallo sviluppo di Taylor $f(P_1)$ attorno a $P_0$ abbiamo:
    \[
    f(P_1) \approx f(P_0) + \sum_j \frac{\partial f}{\partial x_j}(P_0)\cdot(x_j^{(1)}-x_j^{(0)})
    \]
    Da cui, sostituendo nella definizione di errore inerente abbiamo:
    \[
    \boxed{\delta_d =} f(P_0) - f(P_1) = \boxed{- \sum_{j=1}^n \frac{\partial f}{\partial x_j} (P_0) \cdot  (x_j^{(1)} - x_j^{(0)})} = +\sum \rho_j \delta_j
    \]
    dove $\rho_j = \frac{\partial f(P_0)}{\partial x_j} $ è il coefficiente di amplificazione dell'errore inerente e $\delta_j = (x_j^{(1)} - x_j^{(0)})$ è l'errore di arrotondamento nella componente $j$.
\end{definizione}
\begin{definizione}
    [Errore inerente relativo]
    \[
    \boxed{\varepsilon_d = \sum_{j=1}^n \frac{x_j^{(0)}}{f(P_o)}\cdot \frac{\partial f}{\partial x_j} (P_0) \cdot \frac{(x_j^{(1)} - x_j^{(0)})}{x_j^{(0)}}}  = +\sum \gamma_j \varepsilon_j
    \]
    dove $\gamma_j = \frac{x_j^{(0)}}{f(P_o)}\cdot \frac{\partial f}{\partial x_j} (P_0) $ è il coefficiente di amplificazione dell'errore relativo e $ \varepsilon_j = \frac{(x_j^{(1)} - x_j^{(0)})}{x_j^{(0)}}$ è l'errore relativo di arrotondamento nella componente $j$.
\end{definizione}
%aggiungere nota sul come si arriva a questa equazione
\subsection{Sistema mal condizionato (e ben)}
\begin{definizione}
    [Problema malcondizionato]
    Un problema (calcolo funzione) si dice \textbf{malcondizionato} se i $\gamma_j$ sono tali per cui:
    \[
    \varepsilon_d \gg \max_{j=1,\ldots,n} |\varepsilon_j|
    \]
\end{definizione}
\begin{definizione}
    [Problema bencondizionato]
    Al contrario il problema si dice \textbf{bencondizionato} se $|\varepsilon_d| \approx \max_{j} |\varepsilon_j|$.
\end{definizione}

\subsection{Errori inerenti nelle operazioni aritmetiche}

\begin{table}[H]
    \centering
    \begin{tabular}{|c|c|c|}
    \hline
    Op & \textbf{$\delta_d$} & \textbf{$\varepsilon_d$} \\
    \hline
    \textbf{$x \oplus y$} & $\delta_x + \delta_y$ & $\frac{x}{x+y}\varepsilon_x + \frac{y}{x+y} \varepsilon_y$ \\
    \hline
    \textbf{$x \ominus y$} & $\delta_x - \delta_y$ & $\frac{x}{x-y}\varepsilon_x - \frac{y}{x-y} \varepsilon_y$ \\
    \hline
    $y \otimes x $ & $y\delta_x + x\delta_y$ & $\varepsilon_x + \varepsilon_y$ \\
    \hline
    \textbf{$x \oslash y$} & $\frac{1}{y}\delta_x - \frac{x}{y^2}\delta_y$ & $\varepsilon_x - \varepsilon_y$  \\
    \hline
    \end{tabular}
\end{table}

\subsection{Stimare l'errore algoritmico}
Una volta fissato un algoritmo $\tilde{f}_a$, è possibile stimare $\delta_a$ e
$\varepsilon_a$. Si assume che l'\textbf{errore inerente sui dati iniziali sia
nullo} ($P_1 \equiv P_0$) e si studia come si propagano gli errori di
arrotondamento sui dati intermedi, introdotti dalle operazioni
$\oplus, \ominus, \otimes, \oslash$.

\subsection{Limitare dall'alto}

Poiché $\delta_f \doteq \delta_d + \delta_a$ si ha:
\[
|\delta_f| \leq |\delta_d| + |\delta_a| \leq E_d + E_a
\]
dove si definisce il \textbf{limite superiore all'errore inerente}:
\[
E_d = \sum_{j=1}^{n} A_{x_j}\,|\delta_{x_j}|,
\qquad
A_{x_j} = \sup_{x \in D}\left|\frac{\partial f}{\partial x_j}(x)\right|,
\quad
\delta_{x_j} = x_j^{(1)} - x_j^{(0)}
\]
$D$ è detto \textbf{insieme di indeterminazione}: un insieme che contiene i
punti di interesse e tale che $\mathrm{RN}(P_0) = P_1 \in D,\; \forall P_0 \in D$.
\\
Analogamente per l'errore relativo:
\[
|\varepsilon_f| \leq |\varepsilon_d| + |\varepsilon_a|
\]

\begin{definizione}
    [Algoritmo stabile]
    Si definisce come tale un algoritmo con
    $
    |\varepsilon_a| \approx |\varepsilon_d|
    $
    .
\end{definizione}
\begin{definizione}
    [Algoritmo instabile]
    Si definisce come tale un algoritmo con
    $
    |\varepsilon_a| >> |\varepsilon_d|
    $
    .
\end{definizione}


\begin{lezione}
    [Lezione 3 04/03]
\end{lezione}

\section{Problema diretto e inverso con gli errori}

\textbf{Esercizio (tipico d'esame).}
Data la funzione $f(x,y) = (y-x)(x^2+y^2)$, si scriva un algoritmo per il calcolo della funzione e si trovi l'espressione dell'errore algoritmico relativo.
\noindent
\textbf{Risoluzione}: Algoritmo per il Calcolo di $(y-x)(x^2+y^2)$:
\[
\begin{array}{ll}
r_1 \leftarrow x \otimes x & \quad (= x^2)\\
r_2 \leftarrow y \otimes y & \quad (= y^2)\\
r_3 \leftarrow y \ominus x & \quad (= y - x)\\
r_4 \leftarrow r_1 \oplus r_2 & \quad (= x^2 + y^2)\\
r_5 \leftarrow r_3 \otimes r_4 & \quad (= (y-x)(x^2+y^2))
\end{array}
\]

\begin{nota}
Questo non è l'unico algoritmo (= ordine di operazioni) possibile: non esistono risposte univoche per questa parte della domanda.
\end{nota}

\noindent\textbf{Calcolo dell'errore algoritmico relativo (tecnica del grafo).}\\
Si assume $\varepsilon_x = \varepsilon_y = 0$ (nessun errore inerente sui dati iniziali) e si denota con $\varepsilon_i$ l'errore di arrotondamento introdotto dall'operazione $i$-esima.
I coefficienti di amplificazione si ricavano dai nodi del grafo di calcolo:

\begin{figure}[H]
    \centering
    \includegraphics[width=0.6\textwidth]{img/grafoerrore.png}
\end{figure}
\begin{itemize}
  \item $r_1 = x \otimes x$: moltiplicazione, $\varepsilon_{r_1} = \varepsilon_1$.
  \item $r_2 = y \otimes y$: moltiplicazione, $\varepsilon_{r_2} = \varepsilon_2$.
  \item $r_3 = y \ominus x$: sottrazione, $\varepsilon_{r_3} = \varepsilon_3$
    (i coefficienti $\tfrac{y}{y-x}$ e $-\tfrac{x}{y-x}$ amplificano $\varepsilon_y=0$ e $\varepsilon_x=0$).
  \item $r_4 = r_1 \oplus r_2$: addizione con coefficienti $\tfrac{x^2}{x^2+y^2}$ e $\tfrac{y^2}{x^2+y^2}$:
    \[
      \varepsilon_{r_4} = \varepsilon_4 + \frac{x^2}{x^2+y^2}\,\varepsilon_1 + \frac{y^2}{x^2+y^2}\,\varepsilon_2.
    \]
  \item $r_5 = r_3 \otimes r_4$: moltiplicazione, entrambi i coefficienti valgono $1$:
    \[
      \varepsilon_{r_5} = \varepsilon_5 + 1 \cdot \varepsilon_{r_3} + 1\cdot \varepsilon_{r_4}.
    \]
\end{itemize}

Sostituendo $\varepsilon_{r_3}$ e $\varepsilon_{r_4}$ nell'espressione di $\varepsilon_{r_5}$ si ottiene l'errore algoritmico relativo:
\[
\boxed{\varepsilon_a = \varepsilon_5 + \varepsilon_3 + \varepsilon_4 + \frac{x^2}{x^2+y^2}\,\varepsilon_1 + \frac{y^2}{x^2+y^2}\,\varepsilon_2 }
\]
\noindent
\textbf{Problema diretto}: si maggiori l'errore algoritmico relativo in termini della precisione di macchina $U$ (se siamo in doppia precisione $U \approx 10^{-16}$).

\noindent
\textbf{Risoluzione}: Poiché ogni arrotondamento soddisfa $|\varepsilon_i| \leq U$ (ricordarsi che la precisione di macchina limita superiormente l'errore relativo nell'arrotondamento (\autoref{def: prec di macchina})) e i coefficienti di amplificazione sono:
\[
  \frac{x^2}{x^2+y^2} \leq 1, \qquad \frac{y^2}{x^2+y^2} \leq 1,
\]
prendendo il valore assoluto nell'espressione di $\varepsilon_a$:
\[
  |\varepsilon_a| \leq |\varepsilon_5| + |\varepsilon_4| + |\varepsilon_3|
  + \frac{y^2}{x^2+y^2}\,|\varepsilon_2| + \frac{x^2}{x^2+y^2}\,|\varepsilon_1|
  \leq U + U + U + U + U
\]
\[
  \boxed{|\varepsilon_a| \leq 5U}
\]
L'algoritmo è dunque stabile: l'errore algoritmico è limitato da un piccolo multiplo della precisione di macchina, indipendentemente dai valori di $x$ e $y$.



\subsection{Esempi}

\section{Funzioni non razionali}
Se la funzione non può essere valutata solo con $+,-,\times, /$, ad esempio funzioni esponenziali o sin/cos, quello che si fa è considerare $\tilde{f}(x)$ funzione razionale tale che $\tilde{f} \approx f$ almeno in un dominio $D \subseteq \mathbb{R}^n$ che ci interessa.
\\
In tal caso possiamo decomporre l'errore in 3 parti:
\begin{align*}
f(P_0) - \tilde{f}_a(P_1) &= f(P_0) - f(P_1) + f(P_1) - \tilde{f}(P_1) + \tilde{f}(P_1) - \tilde{f}_a(P_1) \\[0.5em]
\delta_f &= f(P_0) - \tilde{f}_a(P_1) \\
\delta_d &= f(P_0) - f(P_1) \\
\delta_{an} &= f(P_1) - \tilde{f}(P_1) \to\text{ errore analitico assoluto} \\
\delta_a &= \tilde{f}(P_1) - \tilde{f}_a(P_1)
\end{align*}
Analogamente: $\varepsilon_{an} = \dfrac{f(P_1) - \tilde{f}(P_1)}{f(P_1)}\to$ errore analitico relativo.\\
Vale che: $\delta_f = \delta_d + \delta_{an} + \delta_a$ e $\varepsilon_f = \varepsilon_d + \varepsilon_{an} + \varepsilon_a$.\\
Nel caso non razionale, l'errore algoritmico $(\delta_a, \varepsilon_a)$ è calcolato per l'approssimazione razionale $\tilde{f}\to$ tecnica grafo.\\
L'errore inerente $(\delta_d, \varepsilon_d)$ è valutato sulla $f$ originale, ma questo non è un problema poiché la formula dell'errore inerente coinvolge le derivate di $f$ e non richiede che si usino solo operazioni elementari.

\subsection{Errore analitico assoluto e relativo}
\begin{definizione}
    [Errore analitico (assoluto)]
    Presente solo per \textbf{funzioni non razionali}: è l'errore introdotto sostituendo $f$ con un'approssimazione razionale $\tilde{f} \approx f$ (ad es.\ troncamento di una serie). Si misura con i dati arrotondati $P_1$, prima di applicare l'algoritmo:
    \[
    \delta_{an} = f(P_1) - \tilde{f}(P_1), \qquad
    \varepsilon_{an} = \frac{f(P_1) - \tilde{f}(P_1)}{f(P_1)}
    \]
    L'errore totale per funzioni non razionali si decompone in tre contributi:
    \[
    \delta_f = \delta_d + \delta_{an} + \delta_a
    \]
\end{definizione}

\subsection{Stimare l'errore analitico}
Vediamo solo il caso $f: \mathbb{R} \to \mathbb{R}$, con $f$ derivabile $k+1$ volte, $k \in \mathbb{N}$. In tal caso si può considerare lo sviluppo di Taylor con resto di Lagrange.\\
Reminder: $f\in C^{n+1}(a,b)$, $x_0\in(a,b)$ allora $\forall x\in(a,b)$ vale
\[
f(x) = f(x_0) + f'(x_0)\cdot(x-x_0) + \frac{f''(x_0)}{2}(x-x_0)^2 + \dots + \frac{f^{(n)}(x_0)}{n!}(x-x_0)^n
+ \underbrace{R(x)}_{\text{resto di Lagrange}}
\]

\[
R(x) = \frac{f^{(n+1)}(\eta)}{(n+1)!}(x-x_0)^{n+1}, \qquad \eta\in(x,x_0)
\]

\[
T_n(x) = \sum_{j=0}^n \frac{f^{(j)}(x_0)}{j!}(x-x_0)^j
\]

\begin{definizione}
    [Stima dell'errore analitico]
    L'errore analitico si stima approssimando $f$ con lo sviluppo di Taylor troncato $T_n$: l'errore commesso è esattamente il resto di Lagrange trascurato,
    \[
    \boxed{f(x) - T_n(x) = \frac{f^{(n+1)}(\eta)}{(n+1)!}(x-x_0)^{n+1}}
    \]
\end{definizione}
\paragraph{Esempio: approssimazione di $e$}
Approssimiamo $e$ con Taylor in $x_0 = 0$ troncato al grado 2. Vogliamo trovare l'errore analitico quando valutiamo $e^x$ per $x \in [-1,1]$.
\[
T_2(x) = 1 + x + \frac{x^2}{2}, \qquad
\delta_{an} = e^x - \left(1 + x + \frac{x^2}{2}\right) = \frac{e^n}{3!}(x-0)^3 = \frac{e^n \cdot x^3}{6}
\]
Maggiorando in modulo l'errore analitico relativo:
\[
|\varepsilon_{an}| \le \max_{n,x\in[-1,1]} \frac{\left|\dfrac{e^n \cdot x^3}{6}\right|}{|e^x|}
\le \frac{\max_{n\in[-1,1]}|e^n| \cdot \max_{x\in[-1,1]}|x^3|}{\min_{x\in[-1,1]} e^x} \cdot \frac{1}{6}
= \frac{e \cdot 1}{e^{-1}} \cdot \frac{1}{6} = \frac{e^2}{6} \approx 1{,}23
\]
Se rifacciamo i conti con un generico $k$ (numero di termini dello sviluppo di Taylor), abbiamo che:
\[
|\varepsilon_{an}| = \frac{|e^x - T_k(x)|}{|e^x|}
\le \frac{\max_{n\in[-1,1]}|e^n| \cdot \max_{x\in[-1,1]}|x^{k+1}|}{\min_{x\in[-1,1]} e^x} \cdot \frac{1}{(k+1)!}
= \frac{e^2}{(k+1)!}
\]
Quantità piccola abbastanza velocemente all'aumentare di $k$.


\chapter{Richiami di algebra lineare}

\section{Basi}

\begin{nota}
    Per convenzione, i vettori sono sempre \textbf{colonna}: $x \in \mathbb{C}^n$ è un vettore $n \times 1$. Questo è implicito ovunque negli appunti: il prodotto $Ax$ è definito solo se $A \in \mathbb{C}^{m \times n}$ e $x$ è $n \times 1$; ogni volta che si scrive $Av = \lambda v$ si sta assumendo $v$ colonna.
\end{nota}

\begin{definizione}[Matrice quadrata] 
    $A \in \mathbb{C}^{n \times n}$, stesso numero di righe e colonne.
\end{definizione}
\begin{definizione}[Prodotto scalare Euclideo] 
    $\langle x, y \rangle = y^H x = \sum_j \overline{y_j} x_j$ per $x,y \in \mathbb{C}^n$.
\end{definizione}
\begin{definizione}[Matrice trasposta] 
    $(A^T)_{ij} = a_{ji}$.
\end{definizione}
\begin{definizione}[Matrice coniugata trasposta] 
    Per matrici complesse, $A^H = (\overline{A})^T$.
\end{definizione}
\begin{nota}
    [Trasposta coniugata di un prodotto]
    \label{note: tras con prod}
    Per la trasposta coniugata di un prodotto di matrici (o matrice per vettore),
    l'ordine dei fattori si inverte:
    \[
    (AB)^H = B^HA^H
    \]
    (analoga alla proprietà della trasposta semplice, $(AB)^T=B^TA^T$, ma con coniugio
    aggiuntivo su ciascun fattore).
\end{nota}
\begin{definizione}[Matrice simmetrica (e anti simmetrica)]
    Anti simmetrica quando $A^T = -A$.
\end{definizione}
\begin{definizione}[Matrice diagonale] Elementi non nulli solo sulla diagonale principale: $a_{ij}=0$ se $i \neq j$.\end{definizione}
\begin{definizione}[Matrice triangolare] \textit{Superiore}: $a_{ij}=0$ se $i>j$. \textit{Inferiore}: $a_{ij}=0$ se $i<j$.\end{definizione}
\begin{definizione}
    [Traccia di una matrice]
    Si definisce come traccia di una matrice la sommatoria degli elementi sulla diagonale:
    \[
    tr(A) = \sum a_{ii}
    \]
\end{definizione}
\begin{nota}
    La traccia è pari alla somma degli autovalori (contati con molteplicità algebrica) per \textbf{ogni} matrice quadrata. Il ragionamento usa due fatti distinti:
    \begin{enumerate}
        \item \textbf{Invarianza per similitudine}: dalla proprietà ciclica della traccia, $tr(S^{-1}AS) = tr(SAS^{-1}) = tr(A)$, quindi matrici simili hanno la stessa traccia.
        \item \textbf{Forma di Jordan}: ogni matrice è simile alla sua forma di Jordan $J$, che è triangolare superiore con gli autovalori sulla diagonale, quindi $tr(J) = \sum_j \lambda_j$.
    \end{enumerate}
    Combinando: $tr(A) = tr(J) = \sum_{j=1}^n \lambda_j$.
\end{nota}

\subsection{Proprietà delle matrici triangolari}
Di seguito alcune delle principali proprietà di questo tipo di matrici.
\begin{itemize}
    \item Gli \textbf{autovalori} sono gli elementi sulla diagonale.
    \item La \textbf{traccia} coincide con la somma degli autovalori (questo vale per ogni matrice quadrata, sempre contando le molteplicità).
    \item Il \textbf{determinante} è dato dal prodotto degli elementi sulla diagonale.
    \item La matrice risulta \textbf{invertibile} se non presenta elementi nulli sulla diagonale.
    \item Il \textbf{prodotto} tra matrici triangolari dello stesso tipo produce una matrice triangolare di quel tipo.
    \item L'\textbf{inversa} di una matrice triangolare di un tipo è anch'essa di quel tipo.
\end{itemize}


\begin{lezione}
    [Lezione 4 08/03]
\end{lezione}

\begin{definizione}
    [Vettori linearmente indipendenti]
    Dei vettori sono linearmente indipendenti se lo $0$ è l'unico fattore a moltiplicarli possibile per ottenere come somma di essi un vettore nullo.
\end{definizione}
\paragraph{Proprietà}
Se il numero di vettori linearmente indipendenti corrisponde alla dimensione di uno spazio vettoriale allora questi vettori si dicono base di quello spazio.

\begin{definizione}
    [Prodotto scalare Euclideo in $\mathbb{R}^n$]
    $\langle x, y \rangle = x^T y = \sum_{j=1}^n x_j y_j$ per $x,y \in \mathbb{R}^n$. Induce la norma $\|x\|_2 = \sqrt{x^Tx}$.
\end{definizione}

\begin{definizione}
    [Vettori ortogonali]
    Due vettori $x$ e $y$ sono ortogonali se:
    \[
    y^Hx=0
    \]
\end{definizione}
\subsection{Operazioni tra matrici}
$ $
\begin{figure}[H]
    \centering
    \includegraphics[width=0.6\linewidth]{img/prodotto matrix.png}
    \caption{Prodotto}
\end{figure}
\noindent
Il \textbf{prodotto} tra matrici è \textbf{possibile se il numero di colonne della prima è uguale al numero di righe della seconda}.\\
La matrice prodotto $P = A\cdot B$ avrà numero di righe pari alla prima  ($A$) e numero di colonne pari alla seconda ($B$).
\paragraph{Proprietà del prodotto tra matrici}
Il prodotto gode della proprietà associativa:
\[
(AB)C=A(BC)
\]
In caso quindi di prodotti con più operandi non importa l'ordine.
\begin{nota}
    È importante notare che, nel caso di più operandi, il costo totale dell'operazione può variare di molto e quindi risulta cruciale scegliere le giuste associazione per ridurre i costi.
\end{nota}

\section{Determinante, rango e matrici inverse}

\subsection{Determinante (anche casi particolari e proprietà)}
\begin{definizione}
    [Determinante]
    Si definisce \textit{determinante} di A la quantità:
    \[
    det(A) = a_{11} \qquad n = 1
    \]
    \[
    det(A) = a_{11}a_{22} - a_{12}a_{21} \qquad n=2
    \]
    \[
    det(A) = \sum^n_{j=1} (-1)^{i+j}a_{ij}\cdot det(A_{ij}) \qquad n\geq 3
    \]
    e rappresenta il fattore di scala del "volume" quando si applica la trasformazione $A$.
\end{definizione}
\paragraph{Casi particolari}
In caso di matrici diagonali o triangolari si ha:
\[
det(A) = \prod^n_{j=1} a_{jj}
\]
Si ha quindi costo $O(n)$.
\begin{definizione}
    [Matrice singolare]
    Matrice A tale che 
    \[
    det(A) = 0
    \]
\end{definizione}
\begin{nota}
    [Importante]
    Ci saranno quindi molto utili le matrici \textbf{non singolari} in quanto queste non "perdono le dimensioni" (più rozzamente "non si schiacciano").
\end{nota}
\paragraph{Determinante e autovalori}
    Per ogni matrice quadrata vale 
    \[det(A) = \prod_{j=1}^n \lambda_j\]
    (prodotto degli autovalori con molteplicità algebrica). La dimostrazione usa due fatti:
    \begin{enumerate}
        \item \textbf{Invarianza per similitudine}: per la moltiplicatività del determinante, $det(S^{-1}AS) = det(S^{-1})det(A)det(S) = det(A)$.
        \item \textbf{Forma di Jordan}: ogni matrice è simile alla sua forma di Jordan $J$ (triangolare superiore), quindi $det(J) = \prod_j \lambda_j$ (ricordarsi anche che matrici simili condividono gli autovalori).
    \end{enumerate}
    Conseguenze dirette: $det(A) = 0 \Longleftrightarrow$ esiste almeno un autovalore nullo $\Longleftrightarrow A$ è singolare. 
    Inoltre \uline{$rk(A)$ coincide con il numero di autovalori non nulli contati con molteplicità}.
    \label{th: rk autovalori non nulli}

    \subsection{Minore di ordine k, principali e principali di testa}
\begin{definizione}
    [Minore di ordine k]
    Determinante di una sottomatrice ottenuta selezionando $k$ righe e $k$ colonne.
\end{definizione}
\begin{definizione}
    [Minore principale]
    Minore in cui gli indici di righe e colonne coincidono.
\end{definizione}
\begin{definizione}
    [Minore principale di testa]
    Gli indici sono $1,...,k$.
\end{definizione}
\begin{definizione}
    [Rango]
    Massimo numero di righe o colonne linearmente indipendenti.
    Può essere trovato anche come:
    Massimo ordine dei k minori $\neq 0$ in A.
    \[
    rk(A) = dim(Im(A))
    \]
    dove l'immagine di A è data dalle x tali che esiste y tale che $Ay=x$.
\end{definizione}
\begin{nota}
    Si può valutare indifferentemente il numero di righe o colonne, il risultato non cambia.
\end{nota}
\begin{teorema}
    [Teorema di Binet-Cauchy]
    \label{th: Binet-Cauchy}
    Dati $A\in \mathbb{C}^{m\times n}$, $B\in \mathbb{C}^{n\times m}$ e $C=A\cdot B \in \mathbb{C}^{m\times m}$ vale che:
    \[
    det(C) = 
    \begin{cases}
        det(A)\cdot det(B) \quad m=n \\
        0 \quad m>n \\
        \sum_jA_{[j]}B_{[j]} \quad m<n
    \end{cases}
    \]
    con $A_{[j]}B_{[j]}$ sottomatrici ottenute con lo stesso sottoinsieme di indici j.
\end{teorema}
\begin{definizione}
    [Matrice Inversa]
    Data $A^{n\times n}$ se esiste si chiama matrice inversa e si indica con $A^{-1}$ l'\uline{unica} matrice tale che:
    \[
    A\cdot A^{-1} = A^{-1} \cdot A = I
    \]
\end{definizione}
\paragraph{Proprietà matrice inversa}
\begin{itemize}
    \item in primis:
    \[
    \boxed{A^{-1} \text{ esiste } \Longleftrightarrow det(A) \neq 0 \Longleftrightarrow rk(A) = n = max}
    \]
    (o anche se e solo se A ha n righe/colonne linearmente indipendenti).
    \item Inoltre si ha che 
    \[
    det(I) = det(A\cdot A^{-1}) = det(A)\cdot det(A^{-1}) \Longrightarrow \boxed{det(A^{-1}) = \frac{1}{det(A)}}
    \]
    \item Inoltre si ha che se scriviamo $B = A^{-1}$:
    \[
    b_{ij} = (-1)^{i+j} \cdot \frac{det(A_{ij})}{det(A)}
    \]
    \item Infine:
    \[
    (A\cdot B)^{-1} = B^{-1}A^{-1} 
    \]
    (vale anche per più di 2 termini, basta sempre invertire l'ordine).
\end{itemize}

\begin{definizione}
    [Matrice Hermitiana]
    Se $A = A^H$, ovvero coincide con la propria trasposta coniugata (o matrice aggiunta), allora si dice \textbf{Hermitiana}.
\end{definizione}
\begin{definizione}
    [Matrice anti Hermitiana]
    Se $A^H = -A$.
\end{definizione}

\paragraph{Proprietà}
Una matrice Hermitiana \uline{ha valori reali sulla diagonale}, infatti per rispettare la definizione tutti gli elementi della diagonale devono essere uguali al loro coniugato, e questo è possibile solo nel caso in cui essi sono reali.
\begin{definizione}[Matrice normale]
    Matrice tale che $A^HA = AA^H$.
\end{definizione}
\begin{definizione}[Matrice unitaria]
    Matrice tale che $A^HA = AA^H = I$, ovvero la sua trasposta coniugata è anche la sua inversa.
\end{definizione}
\begin{definizione}[Matrice ortogonale]
    Matrice tale $A^TA = AA^T = I$, ovvero unitaria nel caso reale: la sua trasposta è anche la sua inversa.
\end{definizione}
\begin{nota}
    Una matrice di questo tipo (unitaria/ortogonale) rappresenta una trasformazione che preserva lunghezze e angoli (isometrie): non deforma lo spazio ma si limita a ruotarlo e/o rifletterlo.
\end{nota}
\begin{definizione}
    [Matrice di permutazione]
    $P \in \mathbb{R}^{n\times n}$ è detta matrice di permutazione se è ottenuta dalla matrice identità permutando le sue righe o le sue colonne.
\end{definizione}
\begin{nota}
    Utilizzo: per lo scambio di righe $M = PA$, mentre per lo scambio di colonne l'ordine da seguire degli operandi diventa $M = AP$.
\end{nota}
\noindent
\paragraph{Proprietà}
Ogni matrice di permutazione è ortogonale:
\begin{enumerate}
    \item $PP^T = P^T P = I\to$ si dimostra facendo vedere che a $P^T$ è associata la permutazione inversa.
    \item Il prodotto di matrici di una sequenza di permutazione $P$ produce una permutazione delle colonne di $A$ $(AP)$ o delle righe $(PA)$.
\end{enumerate}

\section{Sistemi lineari}

\begin{definizione}
    [Sistema lineare]
    Data $A \in \mathbb{C}^{m\times n}$ e $b \in \mathbb{C}^m$ l'equazione $Ax = b$ rappresenta il sistema lineare di $m$ equazioni in $n$ incognite.
    \[
    \begin{cases}
    a_{11}x_1 + \ldots + a_{1n}x_n = b_1 \\
    \ldots + \ldots + \ldots = \ldots &  \Longrightarrow x \in \mathbb{C}^n \text{ vettore di incognite} \\
    a_{m1}x_1 + \ldots + a_{mn}x_n = b_m
    \end{cases}
    \]
\end{definizione}
\begin{teorema}
    [Teorema di Rouché-Capelli]
    \label{th: Rouche-Capelli}
    Data $A\in \mathbb{C}^{m\times n}$ e $b\in\mathbb{C}^{m\times 1}$ esistono soluzioni di $Ax=b$ se e solo se:
    \[
    rk(A) = rk(A|b) \quad con \quad (A|b)\in\mathbb{C}^{m\times n+1}
    \]
    Se vale questo, allora l'unicità della soluzione è garantita per:
    \begin{itemize}
        \item $rk(A) = n$ soluzione unica
        \item $rk(A) < n$ infinite soluzioni
        \item $rk(A) > n$ (possibile solo se $m>n$) l'insieme delle soluzioni è della forma $x_0 + V = \{x_0 + v : v\in V\}$ con $V$ spazio vettoriale di dimensione $n-rk(A)$
    \end{itemize}
    dove $n$ indica il numero di colonne di A (= numero di incognite).\\
    In particolare, per l'ultimo caso, si hanno:
    \begin{itemize}
        \item $V= \{\textit{soluzioni di} Ax=0\} = \{\textit{soluzioni del sistema omogeneo}\} = ker(A)$
        \item $x_0 : \s Ax_0=b$ ovvero $x_0$ soluzione particolare
    \end{itemize}
\end{teorema}

\begin{nota}
    Ricordarsi che il valore del rango non cambia se calcolato sulle righe o sulle colonne, si ha sempre lo stesso valore.
\end{nota}

\paragraph{Osservazione}
\uline{Se $m = n$ e $\det(A) \neq 0$ si ha automaticamente $rk(A) = n$ quindi la soluzione è unica} (vedi \autoref{th: rk autovalori non nulli}).

\begin{definizione}
    [Kernel]
    Con kernel (o nucleo) si intende l'insieme delle soluzioni del sistema omogeneo, ovvero di $Ax=0$.
    Esso rappresenta quindi l'insieme di tutti i vettori che vengono mandati nel vettore nullo da $A$.
\end{definizione}

\begin{definizione}
    [Sistema omogeneo]
    Un sistema si dice omogeneo quando $b=0$, \textbf{ha sempre almeno una soluzione}, ovvero la soluzione banale $x=0$.
\end{definizione}
\paragraph{Soluzione unica di un sistema}
Un sistema ha soluzione unica se
\[
m = n \wedge det(A)\neq 0 \Longrightarrow rk(A) = n \Longrightarrow dim(Ker(A)) = 0 \Longrightarrow Ker(A) = \{0\}
\]
e si individua come $x = A^{-1}b$.

\begin{lezione}
    [Lezione 5 11/03]
\end{lezione}

\section{Autovalori, autovettori e diagonalizzabilità}

\begin{definizione}
    [Autovalore e autovettore]
    Un autovalore, ovvero un valore scalare, e un autovettore, ovvero un vettore ad esso associato, sono una coppia $\lambda, v$ di elementi che soddisfa:
    $$Av = \lambda v$$
    Gli autovettori sono quelle direzioni che rimangono invariate durante la trasformazione $A$, gli autovalori indicano di quanto queste direzioni sono dilatate.
\end{definizione}
\begin{definizione}
    [Polinomio caratteristico]
    \[P_A(\lambda) = det(A - \lambda I)\]
\end{definizione}
\begin{nota}
    Il polinomio caratteristico posto uguale a 0 ci dice per quali valori di $\lambda$ la matrice $A - \lambda I$
    \footnote{Questa matrice è derivata direttamente dalla definizione di autocoppia trasformata in sistema lineare omogeneo: $Av = \lambda v \to (A- \lambda I)v=0$.}
    diventa singolare, ovvero per quali $\lambda$ esiste un vettore (non nullo) che la matrice $A - \lambda I$ manda nel vettore nullo.\\
    Non stiamo facendo altro che cercare dei lambda tali per cui la definizione di autocoppia sia valida:
    \[
    (A-\lambda I)v = Av - \lambda I v = Av - \lambda v = \lambda v - \lambda v = 0
    \]
\end{nota}
\begin{definizione}
    [Molteplicità algebrica]
    Molteplicità di un autovalore come radice del polinomio caratteristico, si indica con $\alpha(\lambda)$.
\end{definizione}

\begin{nota}
    Il rango di una matrice ($rk(A)$) è uguale al numero di autovalori non nulli della matrice contati con la loro molteplicità algebrica, in quanto $det(A) = 0$ se e solo se esiste almeno un autovalore nullo.
\end{nota}
\begin{definizione}
    [Molteplicità geometrica]
    \[
    \gamma (\lambda) = dim(ker(A-I\lambda))
    \]
    ovvero la dimensione dello spazio delle soluzioni del sistema omogeneo, quindi la dimensione dello spazio degli autovettori di A associati a $\lambda$.
    \begin{nota}
        Ricordarsi che gli autovettori associati a un certo autovalore formano un sottospazio.
    \end{nota}
    \noindent
    Per Rouché-Capelli si ha:
    \[
    \gamma(\lambda) = n - rk(A-I\lambda)
    \]
\end{definizione}
\subsection{Matrici simili}
\begin{definizione}
    [Matrici simili]
    \label{def: matrici simili}
    A e B (quadrate) sono simili se esiste S tale che
    \[
    B = S^{-1}AS
    \]
\end{definizione}
\begin{teorema}
    [Proprietà di base delle matrici simili]
    \label{th: prop base matr simili}
    Si hanno 2 proprietà di base:
    \begin{itemize}
        \item Se due matrici sono simili queste \textbf{condividono gli autovalori} (con rispettive molteplicità).
        \item gli autovettori della matrice B sono dati da $S^{-1}v$.
    \end{itemize}
\end{teorema}
\paragraph{Proprietà}
Per effetto della definizione e della condivisione di autovalori si ha che due matrici simili condividono anche stessi:
\begin{itemize}
    \item determinante: $det(A) = \prod_j \lambda_j$ (vedi nota nella sezione Determinante)
    \item raggio spettrale
    \item polinomio caratteristico
\end{itemize}
\begin{dimostrazionelibera}
    [Stesso polinomio caratteristico]
    \[
    \det(\lambda I - B) = \det(\lambda I - S^{-1}AS) = \det\!\left(S^{-1}(\lambda I - A)S\right)=
    \]
    \[
    = \underbrace{\det(S^{-1})}_{1/\det S}\,\det(\lambda I - A)\underbrace{\det(S)}_{\det S}
    = \det(\lambda I - A)
    \]
\end{dimostrazionelibera}
\subsection{Matrice diagonalizzabile}
\begin{definizione}
    [Matrice diagonalizzabile]
    Una matrice si dice diagonalizzabile se simile a una matrice diagonale:
    \[
    S^{-1}AS = diag(\lambda_1, \dots, \lambda_n)
    \]
    Questo implica, per il \autoref{th: prop base matr simili}, che tale matrice diagonale contenga gli autovalori di A, mentre la matrice $S$ conterrà gli autovettori (questo si ricava facilmente dalla definizione di autovettori e autovalori).\\
\end{definizione}
\begin{teorema}
    [Diagonalizzabilità di una matrice]
    Per determinare se una matrice sia diagonalizzabile facciamo affidamento alla seguente relazione:
    \[
    A \in \mathbb{C}^{n\times n} \text{ è diagonalizzabile } \Longleftrightarrow \alpha(\lambda) = \gamma(\lambda) \quad \forall \lambda \textit{ autovalore di }A
    \]
    dove ricordiamo essere $\alpha$ la molteplicità algebrica e $\lambda$ la molteplicità geometrica.\\
    Concludiamo quindi anche che \textbf{una matrice è diagonalizzabile se ha tutti gli autovalori distinti}
    \footnote{Questo tornerà utile nelle conseguenze del secondo teorema di Gershgorin.}.
\end{teorema}
\paragraph{Osservazione}
Dire che una matrice è diagonalizzabile equivale a dire che esiste una base di $\mathbb{C}^n$ formata da autovettori della matrice (esiste infatti una matrice simile, con cui quindi condivido gli autov).
\begin{nota}
    Una matrice diagonalizzabile rappresenta concettualmente una trasformazione che, scelto il giusto sistema di riferimento (la base degli autovettori), si riduce a una serie di \textbf{scalature indipendenti} lungo $n$ direzioni fisse - nessuna rotazione, nessun "mescolamento" tra le direzioni, solo stiramento/compressione (ed eventualmente inversione, se $\lambda$ < 0) lungo ciascun asse.
\end{nota}
\begin{nota}
    Se una matrice quadrata ha tutti gli autovalori distinti \textbf{allora} ognuno di essi ha le molteplicità uguali a 1.
\end{nota}
\begin{teorema}
    [Teorema spettrale]
    \label{th: teorema spettrale}
    Data $A\in \mathbb{C}^{n\times n}$ Hermitiana allora A è diagonalizzabile con autovalori reali, esiste pertanto una base unitaria di $\mathbb{C}^n$ fatta di autovettori di A.
\end{teorema}
\paragraph{Proprietà autovalori e autovettori}
\begin{enumerate}
    \item \textbf{Importante} per esercizi esame:\\
    Sia $(\lambda, v)$ autocoppia di $A \in \mathbb{C}^{n\times n}$ e $\delta \in \mathbb{C}$ $\implies$ $(\lambda + \delta, v)$ è autocoppia di $A + \delta I$.
    
    \item Dato un polinomio $p(x) = \displaystyle\sum_{j=0}^{s} p_j x^j$ possiamo definire la matrice
    \[
    P(A) = \sum_{j=0}^{s} p_j A^j = p_0 I + p_1 A + \ldots + p_s A^s
    \]
    
    se $(\lambda, v)$ è autocoppia di $A$ $\implies$ $(p(\lambda), v)$ è autocoppia di $P(A)$
    
    \item \label{item: proprieta autovalori inversa} \textbf{Importante} per esercizi esame:\\
    Se $A \in \mathbb{C}^{n\times n}$ invertibile e $(\lambda, v)$ è autocoppia di $A$ $\implies$ $(\lambda^{-1}, v)$ è autocoppia di $A^{-1}$.
    \begin{dimostrazionelibera}
        \[
        Av= \lambda v \implies A^{-1}Av = \lambda A^{-1}v \implies v = \lambda A^{-1}v \implies \lambda^{-1} v = A^{-1}v
        \]
    \end{dimostrazionelibera}
    \item Se $A \in \mathbb{R}^{n\times n}$ $\implies$ $P_A(\lambda)$ è un polinomio a coefficienti reali. In particolare, se $\lambda$ verifica $P_A(\lambda) = 0$ allora
    \[
    \overline{P_A(\lambda)} = (-1)(\overline{\lambda})^n + (-1)^{n-1}\sigma_1(\overline{\lambda})^{n-1} + \ldots + \sigma^n = P_A(\overline{\lambda})
    \]
    
    se $\lambda$ è autovalore $\implies$ $\overline{\lambda}$ è autovalore.
\end{enumerate}
\begin{definizione}
    [Raggio spettrale]
    \[
    \rho(A) = \underset{j= 1, ..., n}{max} |\lambda_j|
    \]
\end{definizione}
\begin{definizione}
    [Matrice convergente]
    Una matrice si dice tale se:
    $$\rho(A)<1$$
\end{definizione}

\begin{teorema}
    [Teorema matrice convergente]
    \label{th: matrice convergente}
    \[
    A \text{ è convergente } \Longleftrightarrow \lim_{k\to\infty} A^K = 0 \quad \textit{(con 0 = matrice nulla)}
    \]
\end{teorema}


\begin{lezione}
    [Lezione 6 15/03]
\end{lezione}

\section{Norme}

\subsection{Norme vettoriali}
Esempi di norme
\begin{itemize}
    \item $||x||_1 = \sum^n_j |x_j|$ (spostamento a griglia)
    \item $||x||_2 = \sqrt{\sum^n_j |x_j|^2}$
    \item $||x||_\infty = max|x_j|$
\end{itemize}

\begin{definizione}
    [Norme equivalenti]
    Due norme vettoriali $||\cdot||_p, ||\cdot||_q$ si dicono \textbf{equivalenti} se $\exists \, \alpha, \beta \in \mathbb{R}^+$ tali per cui $\forall x \in \mathbb{C}^n$:
    \[
    \alpha \, ||x||_p \leq ||x||_q \leq \beta \, ||x||_p
    \]
    Se vale ciò allora $\beta^{-1} ||x||_q \leq ||x||_p \leq \alpha^{-1} ||x||_q$.
\end{definizione}
\noindent
Quindi tutte le norme su $\mathbb{C}^n$ sono equivalenti
\begin{itemize}
    \item $||x||_2 \leq ||x||_1 \leq \sqrt{n} \, ||x||_2$
    \item $||x||_\infty \leq ||x||_1 \leq \sqrt{n} \, ||x||_\infty$
    \item $||x||_\infty \leq ||x||_2 \leq \sqrt{n} \, ||x||_\infty$
\end{itemize}
\begin{nota}
    il calcolo delle norme vettoriali costa $O(n)$.
\end{nota}

\subsection{Norme matriciali}
\begin{definizione}
    [Norma indotta]
    Una norma matriciale $||\cdot||_m$ si dice \textbf{indotta} da una norma vettoriale $||\cdot||_v$ se è definita come:
    \[
    ||A||_m = \sup_{x \in \mathbb{C}^n} \frac{||Ax||_v}{||x||_v}
    \]
    Esempi di norme indotte da $||\cdot||_1, ||\cdot||_2, ||\cdot||_\infty$ :
    \begin{itemize}
        \item $||A||_1 $        
        \item $||A||_2 $        
        \item $||A||_\infty $ 
    \end{itemize}
    Esempio di norma non indotta (norma di Frobenius):
    \[
    ||A||_F := \sqrt{\sum_{i,j=1,\ldots,n} |a_{ij}|^2} \to\text{ non è indotta perché } ||I||_F = \sqrt{n} \neq 1
    \]
    \uline{Per tutte le norme indotte \textbf{deve} valere} $\boxed{||I|| = 1}$, infatti:
    \[
    ||I|| = \sup_{x \neq 0} \frac{||I \cdot x||}{||x||} = \sup_{x \neq 0} \frac{||x||}{||x||} = 1
    \]
\end{definizione}

\begin{definizione}
    [Norma 1]
    \[
    ||A||_1 = \underset{j}{max} \sum_i^n |a_{ij}|
    \]
    Data dalla massima sommatoria dei moduli degli elementi di una colonna.
\end{definizione}
\begin{definizione}
    [Norma 2]
    \[
    ||A||_2 = \sqrt{\rho (A^HA)}
    \]
    Data dalla radice del raggio spettrale della matrice $A^HA$ (ovvero radice del massimo modulo degli autovalori di $A^HA$).
\end{definizione}
\begin{definizione}
    [Norma di Frobenius]
    \[
    ||A||_F = \sqrt{\sum |a_{ij}|^2}
    \]
    Data dalla radice della sommatoria del modulo degli elementi al quadrato.
    \begin{nota}
        Questa norma non è indotta: $||I||_F = \sqrt{n} \neq 1$.
    \end{nota}
\end{definizione}
\begin{definizione}
    [Norma infinita]
    \[
    ||A||_\infty = \underset{i}{max} \sum_j^n |a_{ij}|
    \]
    Data dalla massima sommatoria dei moduli degli elementi di una riga.
\end{definizione}

\begin{definizione}
    [Norma compatibile]
    Una norma matriciale $||\cdot||_m$ si dice \textbf{compatibile} con una norma vettoriale $||\cdot||_v$ se $\forall A \in \mathbb{C}^{n\times n}$ e $\forall x \in \mathbb{C}^n$ vale: $||Ax||_v \leq ||A||_m \cdot ||x||_v$.
\end{definizione}
\noindent
\textbf{Proprietà}:
Se $||\cdot||_m$ è la norma matriciale indotta da una vettoriale $||\cdot||_v$ allora $||\cdot||_m$ e $||\cdot||_v$ sono compatibili.
\begin{dimostrazionelibera}
    $\forall x \neq 0$ $\quad \dfrac{||Ax||_v}{||x||_v} \leq ||A||_m$
     vale perché la norma è definita come il massimo della quantità a sinistra.
\end{dimostrazionelibera}

\begin{teorema}
    [Teorema di Hirsh]
    \label{th: Hirsh}
    Data $A\in \mathbb{C}^{n\times n}$ allora per ogni norma matriciale vale:
    \[
    \rho(A) = max\{|\lambda|:\lambda \text{ autovalore di A }\} \leq ||A||
    \]
\end{teorema}

\begin{dimostrazione}
    Sia $(\lambda, x)$ autocoppia di $A$ con $x \neq 0$, ovvero $Ax = \lambda x$.
    Costruiamo la matrice
    \[
        B = \bigl[\, x \mid 0 \mid \cdots \mid 0 \,\bigr] \in \mathbb{C}^{n \times n}
    \]
    avente $x$ come prima colonna e le restanti colonne nulle. Allora:
    \[
        A \cdot B
        = \bigl[\, Ax \mid A \cdot 0 \mid \cdots \,\bigr]
        = \bigl[\, Ax \mid 0 \mid \cdots \,\bigr]
        = \bigl[\, \lambda x \mid 0 \mid \cdots \,\bigr]
        = \lambda \bigl[\, x \mid 0 \mid \cdots \,\bigr]
        = \lambda B
    \]
    Per la sottomoltiplicatività della norma matriciale:
    \[
        |\lambda| \cdot \|B\| = \|\lambda B\| = \|A \cdot B\| \leq \|A\| \cdot \|B\|
    \]
    Poiché $x \neq 0$ si ha $\|B\| \neq 0$, quindi si può dividere per $\|B\|$:
    \[
        |\lambda| \leq \|A\|
    \]
    Poiché $\lambda$ è un autovalore arbitrario di $A$, prendendo il massimo tra tutti gli autovalori:
    \[
        \rho(A) = \max_j |\lambda_j| \leq \|A\|
    \]
\end{dimostrazione}
\begin{nota}
    Hirsch vale anche per le norme \textbf{non indotte}. Geometricamente, gli autovalori di $A$
    si trovano tutti all'interno del disco di raggio $\|A\|$ centrato nell'origine del piano complesso.
\end{nota}

\section{Localizzazione degli autovalori}

\subsection{Teoremi di Gershgorin}
\begin{definizione}
[Cerchi di Gershgorin]
Data $A\in \mathbb{C}^{n\times n}$ definiamo gli insiemi
\[
\mathcal{F}_i(A) = \{z \in \mathbb{C} : |z - a_{ii}| \leq \rho_i \}
\]
con raggio del cerchio $\rho_i = \sum^n_{j\neq i}|a_{ij}|$ e con $z$ l'insieme dei punti nello spazio che rispettano tale condizione.\\
$\mathcal{F}_i(A)$ si dice cerchio di Gershgorin associato alla riga \textit{i} di \textit{A}.
\end{definizione}
\begin{teorema}
    [Primo teorema di Gershgorin]
    Dato $\lambda$ autovalore di A allora \[\lambda \in \bigcup  \mathcal{F}_i(A)\]
\end{teorema}

\begin{dimostrazione}
    Sia $(\lambda, x)$ autocoppia di $A$ con $x \neq 0$, ovvero $Ax = \lambda x$.
    Scriviamo $x = [x_1, \dots, x_n]^T$ e scegliamo l'indice
    \[
        k \in \{1, \dots, n\} \quad \text{tale che} \quad
        |x_k| \geq |x_j| \quad \forall j \in \{1, \dots, n\}
    \]
    ovvero $k$ è l'indice della componente di modulo massimo. In particolare,
    poiché $x \neq 0$, si ha $|x_k| > 0$.
    \\
    Consideriamo la $k$-esima riga dell'equazione $Ax = \lambda x$:
    \[
        \sum_{j=1}^{n} a_{kj} x_j = \lambda x_k
    \]
    Separiamo il termine $j = k$:
    \[
        \sum_{\substack{j=1 \\ j \neq k}}^{n} a_{kj} x_j = \lambda x_k - a_{kk} x_k
        = (\lambda - a_{kk})\, x_k
    \]
    Prendendo il modulo di entrambi i membri e applicando la disuguaglianza triangolare:
    \[
        |\lambda - a_{kk}| \cdot |x_k|
        = \left| \sum_{\substack{j=1 \\ j \neq k}}^{n} a_{kj} x_j \right|
        \leq \sum_{\substack{j=1 \\ j \neq k}}^{n} |a_{kj}| \cdot |x_j|
    \]
    Dividiamo per $|x_k| > 0$:
    \[
        |\lambda - a_{kk}|
        \leq \sum_{\substack{j=1 \\ j \neq k}}^{n} |a_{kj}| \cdot \frac{|x_j|}{|x_k|}
    \]
    Per scelta di $k$ si ha $|x_j| \leq |x_k|$, dunque $\dfrac{|x_j|}{|x_k|} \leq 1$ per ogni $j$, e quindi:
    \[
        |\lambda - a_{kk}|
        \leq \sum_{\substack{j=1 \\ j \neq k}}^{n} |a_{kj}|
        = \rho_k
        \quad \Longrightarrow \quad \lambda \in \mathcal{F}_k(A)
    \]
    Poiché $\mathcal{F}_k(A) \subseteq \displaystyle\bigcup_{i=1}^{n} \mathcal{F}_i(A)$, si conclude:
    \[
        \lambda \in \bigcup_{i=1}^{n} \mathcal{F}_i(A)
    \]
    sta in un cerchio quindi sta nell'unione di tutti.
\end{dimostrazione}

\begin{definizione}
    [Matrice a predominanza diagonale forte]
    \[
    |a_{jj}| > \sum^n_{h\neq j}|a_{jh}| \quad \forall j = 1,...,n
    \]
\end{definizione}

\begin{teorema}
    [Primo corollario di Gershgorin]
    Se A è a predominanza diagonale forte allora A è non singolare (ovvero $det(A)\neq 0$).\\
\end{teorema}
\begin{dimostrazione}
    I cerchi di Gershgorin non contengono lo 0 perché il raggio del cerchio è più piccolo del modulo del centro. Quindi il disco non può contenere lo 0.
\end{dimostrazione}

\begin{lezione}
    [Lezione 7 18/03]
\end{lezione}

\begin{teorema}
    [Secondo teorema di Gershgorin]
    Sia $A \in \mathbb{C}^{n \times n}$, $1 \le s < n$ (con $s$ intero) e sia $M_1$ dato dall'unione di $s$ cerchi di Gershgorin associati ad $A$, e sia $M_2$ l'unione dei restanti $n - s$ cerchi di Gershgorin associati ad $A$.
    \\
    Se
    \[
    M_1 \cap M_2 = \varnothing
    \]
    allora $M_1$ contiene esattamente $s$ autovalori e $M_2$ contiene esattamente $n - s$ autovalori (contati rispetto alla molteplicità algebrica).
\end{teorema}

\paragraph{Conseguenze dirette del teorema}
\begin{enumerate}
    \item \uline{Se $A$ ha $n$ cerchi disgiunti, allora $A$ ha $n$ autovalori distinti e quindi $A$ è diagonalizzabile}.
    
    \item Se $A \in \mathbb{R}^{n \times n}$ e $F_i(A)$ è un cerchio isolato (cioè non tocca gli altri), ovvero
    \[
    F_i(A) \cap \left( \bigcup_{\substack{j=1 \\ j \ne i}}^{n} F_j(A) \right) = \varnothing,
    \]
    allora $F_i(A)$ contiene un autovalore \textbf{reale}.
\end{enumerate}

\begin{dimostrazione}
    [Conseguenza 2]
    Sia $\lambda$ autovalore di $A$ con $\lambda \in \mathcal{F}_i(A)$. Per il secondo teorema di
    Gershgorin, poiché $\mathcal{F}_i(A)$ è isolato ($s = 1$, $M_1 = \mathcal{F}_i(A)$,
    $M_1 \cap M_2 = \varnothing$), $\mathcal{F}_i(A)$ contiene \textbf{esattamente un autovalore}.
    
    \medskip\noindent
    Supponiamo per assurdo che $\lambda \notin \mathbb{R}$, ovvero $\lambda \in \mathbb{C} \setminus \mathbb{R}$.\\
    Poiché $A \in \mathbb{R}^{n \times n}$, per la proprietà 4 degli autovalori, $\bar{\lambda}$ è
    anch'esso autovalore di $A$, con $\bar{\lambda} \neq \lambda$.

    \medskip\noindent
    Osserviamo ora che $\mathcal{F}_i(A)$ è simmetrico rispetto all'asse reale: il suo centro
    $a_{ii} \in \mathbb{R}$ (in quanto $A$ reale) e il raggio $\rho_i \in \mathbb{R}$, quindi
    \[
        \lambda \in \mathcal{F}_i(A)
        \;\Longrightarrow\;
        |\lambda - a_{ii}| \leq \rho_i
        \;\Longrightarrow\;
        |\bar{\lambda} - a_{ii}| = \overline{|\lambda - a_{ii}|} = |\lambda - a_{ii}| \leq \rho_i
        \;\Longrightarrow\;
        \bar{\lambda} \in \mathcal{F}_i(A)
    \]
    Dunque $\mathcal{F}_i(A)$ contiene due autovalori distinti ($\lambda$ e $\bar{\lambda}$),
    contraddicendo il secondo teorema di Gershgorin.

    \medskip\noindent
    Per assurdo, deve essere $\lambda \in \mathbb{R}$.
\end{dimostrazione}

\begin{definizione}
    [Grafo associato]
    Un grafo associato a una matrice è identificato come
    \[
    G(A) = (V, E)
    \]
    dove $V$ sono i vertici che vanno da $1$ a $n$ ($n$ dimensione della matrice quadrata) e $E$ sono gli archi definiti come $E = \{ (i,j) : a_{ij} \neq 0 \}$.
\end{definizione}
\begin{definizione}
    [Grafo fortemente connesso]
    Un grafo si dice fortemente connesso se per ogni coppia di vertici $(i, j)\s\exists$ un cammino orientato che parte da $i$ e arriva a $j$.
\end{definizione}
\begin{definizione}
    [Matrice irriducibile]
    Una matrice A il cui grafo associato G(A) è fortemente connesso si dice irriducibile.
\end{definizione}
\begin{nota}
    questo significa che non si possono riordinare le incognite in modo da separare il sistema in due sottosistemi indipendenti: tutte le equazioni "dipendono" tra loro, direttamente o indirettamente.
\end{nota}
\begin{teorema}
    [Terzo teorema di Gershgorin]
    Sia $A \in \mathbb{C}^{n \times n}$ una matrice \textbf{irriducibile} e con $\lambda$ autovalore di $A$ allora si ha:
    \[
    \lambda \in 
    \left\{
    \begin{array}{c}
    \text{bordo dell'unione dei} \\
    \text{cerchi di Gershgorin di A}
    \end{array}
    \right\}
    \Longrightarrow
    \lambda \in 
    \text{bordo di ogni cerchio } \mathcal{F}_i(A)
    \]
\end{teorema}
\begin{nota}
    Senza irriducibilità, l'affermazione sarebbe falsa. Una matrice riducibile si può permutare in forma triangolare a blocchi: 
    in quel caso gli autovalori di una riga e dell'altra sono indipendenti, e un autovalore di una riga "ignora" completamente i cerchi di Gershgorin associati alle altre righe.
    Il primo può quindi stare sul bordo dell'unione senza avere nulla a che fare con i cerchi degli altri blocchi.
\end{nota}
\begin{definizione}
    [Matrice a predominanza diagonale debole]
    \[
    |a_{jj}| \geq \sum^n_{h\neq j}|a_{jh}| \quad \forall j = 1,...,n
    \]
\end{definizione}
\begin{teorema}
    [Terzo corollario di Gershgorin]
    \label{th: terzo cor gersh}
    Sia $A$ a predominanza diagonale debole, irriducibile, \uline{e} tale che \uline{esista almeno un indice}
    $j_0\in\{1,\dots,n\}$ per cui la disuguaglianza sia \textbf{stretta}, ovvero
    \[
    |a_{j_0j_0}| > \sum^n_{h\neq j_0}|a_{j_0h}|
    \]
    Allora $A$ è non singolare ($\det(A)\neq 0$) e quindi invertibile.
\end{teorema}
\begin{dimostrazione}
    Se $0$ dovesse essere autovalore di $A$ allora dovrebbe stare sul bordo di tutti i cerchi. Dato che però c'è un cerchio che non tocca lo $0$ (quello relativo alla riga per la quale vale la predominanza forte) allora lo $0$ non può appartenere alla lista degli autovalori.
\end{dimostrazione}

\section{Matrici riducibili e a blocchi}

\subsection{Matrici a blocchi}
Scrivere una matrice in termine delle sue sottomatrici al posto delle entrare scalari.
\begin{nota}
    Nei casi rilevanti si assume che i blocchi sulla diagonale sono quadrati e che la matrice A sia quadrata.
\end{nota}
\paragraph{Determinante matrici diagonali a blocchi}
In questo caso il determinante è dato dal prodotto dei determinanti delle matrici a blocco che formano la diagonale.
\begin{definizione}
    [Matrice Riducibile]
    Data $A\in \mathbb{C}^{n\times n}$ con $n\geq2$ si dice riducibile se $\exists \pi \in \mathbb{R}$ tale che
    \[
    \pi A\pi^T = 
    \left[
    \begin{matrix}
        A_{11} & A_{12} \\
        0 & A_{22}
    \end{matrix}
    \right]
    \]
    ovvero se esiste una permutazione che porta A in forma triangolare a blocchi superiore.
\end{definizione}
\begin{nota}
    Se una matrice è riducibile allora esisterà anche una permutazione che la trasforma in forma diagoanle a blocchi inferiore.
\end{nota}
\subsection{Risolvere un sistema lineare con matrice riducibile}
Si consideri il sistema lineare
$
Ax = b, \; A \in \mathbb{C}^{n \times n}, \; b \in \mathbb{C}^n, \; x \in \mathbb{C}^n.
$
Supponiamo che la matrice \(A\) sia riducibile, cioè che esista una matrice di permutazione \(\pi\) tale che
\[
\pi A \pi^T =
\begin{pmatrix}
A_{11} & A_{12} \\
0      & A_{22}
\end{pmatrix}
= B.
\]
Poiché \(\pi^T \pi = I\), vale la seguente equivalenza:
\[
Ax = b 
\;\Longleftrightarrow\;
\pi A x = \pi b
\;\Longleftrightarrow\;
\pi A \pi^T (\pi x) = \pi b.
\]
Definendo:
$y = \pi x$ (permutazione delle incognite), e 
$c = \pi b$ (permutazione dei termini noti),
il sistema diventa:
\[
By = c.
\]
Scrivendo esplicitamente il sistema a blocchi:
\[
\begin{pmatrix}
A_{11} & A_{12} \\
0      & A_{22}
\end{pmatrix}
\begin{pmatrix}
y_1 \\
y_2
\end{pmatrix}
=
\begin{pmatrix}
c_1 \\
c_2
\end{pmatrix},
\]
si ottiene il sistema equivalente:
\[
\begin{cases}
A_{11} y_1 + A_{12} y_2 = c_1 \\
A_{22} y_2 = c_2
\end{cases}
\]
La procedura risolutiva è quindi:
\begin{enumerate}
  \item Risolvo il sistema \(A_{22} y_2 = c_2\), di dimensione \((n-k) \times (n-k)\)
  \item Sostituisco nella prima equazione:
  \[
  A_{11} y_1 = c_1 - A_{12} y_2;
  \]
  \item Risolvo il sistema \(A_{11} y_1 = c_1 - A_{12} y_2\), di dimensione \(k \times k\).
\end{enumerate}

\subsubsection{Costo computazionale}

Assumendo che la risoluzione diretta di un sistema lineare di dimensione \(n\) costi \(O(n^3)\), si ha:
\begin{itemize}
  \item Procedura diretta su \(Ax=b\): \(O(n^3)\);
  \item Procedura a blocchi:
  $O\bigl((n-k)^3 + k^3\bigr) + \text{termini di ordine inferiore}.$
\end{itemize}

Nel caso migliore \(k = \frac{n}{2}\), si ottiene:
\[
O\left(\left(n - \frac{n}{2}\right)^3 + \left(\frac{n}{2}\right)^3\right)
= O\left(\frac{n^3}{4}\right),
\]
con un guadagno di un fattore \(4\) nel numero di operazioni.
\begin{nota}
    Questa strategia può essere applicata ricorsivamente ai blocchi \(A_{11}\) e \(A_{22}\) qualora essi risultino a loro volta riducibili.
\end{nota}

\begin{lezione}
    [Lezione 8 22/03]
\end{lezione}

\begin{teorema}
    [Teorema matrice riducibile (con grafo)]
    $A\in \mathbb{C}^{n\times n}$ è riducibile $\Longleftrightarrow$ G(A) \textbf{non} è fortemente connesso
\end{teorema}

\begin{dimostrazione}
    %controllare da quaderno
    \underline{($\Rightarrow$)} $A$ è riducibile $\implies \exists\,\pi$ tale che
    \[
    B = \pi A \pi^T = \begin{bmatrix} A_{11} & A_{12} \\ 0 & A_{22}\end{bmatrix}, \qquad
    \begin{array}{l} k \text{ righe} \\ n-k \text{ righe}\end{array}
    \]
    $\implies G(B)$ non è fortemente connesso perché da un indice in $\{k+1,k+2,\dots,n\}$ non si può arrivare a nessun indice nell'insieme $\{1,\dots,k\}$ $\implies G(A)$ non è fortemente connesso.

    \bigskip
    \noindent
    \underline{($\Leftarrow$)} $G(A)$ non fortemente connesso $\implies \exists\, i,h\in\{1,\dots,n\}$, $i\neq h$, tali per cui partendo da $i$ non posso arrivare ad $h$.
    Allora consideriamo:
    \[
    P = \{x\in\{1,\dots,n\} : x \text{ raggiungibile da } i\}, \qquad
    Q = \{x\in\{1,\dots,n\} : x \text{ irraggiungibile da } i\}
    \]
    $P\neq\varnothing$ (contiene $i$), $Q\neq\varnothing$ (contiene almeno $h$), $P\cup Q = \{1,\dots,n\}$.

    \begin{nota}
        Non ci possono essere archi che partono da un elemento di $P$ e raggiungono un elemento di $Q$.
    \end{nota}

    Allora consideriamo una permutazione $\pi$ tale che sposta in testa tutti gli elementi di $Q$ e sposta in coda tutti gli elementi di $P$:
    \[
    \pi A \pi^T =
    \begin{array}{cc} \scriptstyle Q & \scriptstyle P \end{array}
    \begin{bmatrix} A_{11} & A_{12} \\ 0 & A_{22} \end{bmatrix}
    \begin{array}{l} \scriptstyle Q \\ \scriptstyle P \end{array}
    \implies A \text{ riducibile}
    \]
    \begin{nota}
        La dimostrazione fornisce anche un metodo costruttivo per trovare $\pi$.
    \end{nota}
\end{dimostrazione}

\subsection{Esempi (skippable)}

Sia
\[
A = \begin{bmatrix} 3 & 0 & 6 \\ 2 & 4 & 7 \\ 3 & 0 & 6\end{bmatrix}
\]
con grafo associato $G(A)$: archi $1\to3$, $3\to1$, $2\to1$, $2\to3$ (nessun arco entra nel nodo 2).

\begin{center}
\begin{tabular}{c|c|c}
    nodo & nodi raggiungibili & nodi irraggiungibili \\
    \hline
    1 & 1, 3 & 2 \\
    2 & 1, 2, 3 & $\varnothing$ \\
    3 & 1, 3 & 2
\end{tabular}
\end{center}

Tagliamo $P=\{1,3\}$, $Q=\{2\}$ e prendiamo la permutazione $2\to1$, $1\to2$, $3\to3$:
\[
\pi = \begin{bmatrix} 0&1&0 \\ 1&0&0 \\ 0&0&1\end{bmatrix}
\qquad\implies\qquad
\pi A \pi^T =
\begin{bmatrix}
4 & 2 & 7 \\
0 & 3 & 6 \\
0 & 3 & 6
\end{bmatrix}
% NOTA DI TRASCRIZIONE: l'ultimo elemento in basso a destra, nella foto, sembra
% leggersi "0" invece di "6" -- ma ricalcolando esplicitamente pi*A*pi^T a partire
% dalla A data (con A_{33}=6), l'unico valore algebricamente coerente è 6.
% Verifica sul quaderno se la matrice A originale avesse un valore diverso in A_{33}.
\]
che è in forma triangolare a blocchi (blocco $A_{11}$ di taglia $1\times1$, blocco $A_{22}$ di taglia $2\times2$) $\implies A$ riducibile.

\begin{nota}
    Si può provare ad applicare il procedimento di riduzione a forma triangolare a blocchi in maniera \textbf{ricorsiva}. Supponiamo di aver trovato $\pi$ tale che
    \[
    \pi A \pi^T = \begin{bmatrix} A_{11} & A_{12} \\ 0 & A_{22}\end{bmatrix}
    \]
    Si prova ad applicare lo stesso metodo ad $A_{11}$ e $A_{22}$, se sono a loro volta riducibili. Se lo sono entrambi, si trovano $\pi_1\in\mathbb{R}^{k\times k}$, $\pi_2\in\mathbb{R}^{(n-k)\times(n-k)}$ matrici di permutazione tali che
    \[
    \pi_1 A_{11}\pi_1^T = \begin{bmatrix} A_{11}^{(1)} & A_{12}^{(1)} \\ 0 & A_{22}^{(1)}\end{bmatrix}, \qquad
    \pi_2 A_{22}\pi_2^T = \begin{bmatrix} A_{33}^{(1)} & A_{34}^{(1)} \\ 0 & A_{44}^{(1)}\end{bmatrix}
    % NOTA DI TRASCRIZIONE: i pedici del secondo blocco (33,34,44 vs. quanto letto
    % come "23,34,44" nella foto) sono stati resi coerenti con la numerazione
    % progressiva dei 4 sotto-blocchi risultanti dallo split di A11 e A22.
    \]
    Combinando le due permutazioni:
    \[
    \begin{bmatrix}\pi_1 & 0 \\ 0 & \pi_2\end{bmatrix}(\pi A \pi^T)\begin{bmatrix}\pi_1^T & 0 \\ 0 & \pi_2^T\end{bmatrix}
    = \begin{bmatrix} A_{11}^{(1)} & A_{12}^{(1)} & * \\ 0 & A_{22}^{(1)} & * \\ \hline & 0 & A_{33}^{(1)}\end{bmatrix}
    \]
    Continuando ricorsivamente sui blocchi diagonali si arriva a una matrice $t\times t$ triangolare a blocchi:
    \[
    \begin{bmatrix} A_{11} & \cdots & A_{1t} \\ & \ddots & \vdots \\ & & A_{tt}\end{bmatrix}
    \]
\end{nota}

\subsection{Sistemi con matrice triangolare a blocchi}

Il sistema lineare associato a una matrice $t\times t$ triangolare a blocchi è della forma
\[
\begin{cases}
A_{11}x_1 = c_1 - A_{12}x_2 - \cdots - A_{1t}x_t \\
\quad\vdots \\
A_{t-1,t-1}x_{t-1} = c_{t-1} - A_{t-1,t}x_t \\
A_{tt}x_t = c_t \quad\to \text{sostituisco nell'equazione sopra}
\end{cases}
\]
Questo è il \textbf{metodo di sostituzione all'indietro (a blocchi)}.

\subsection{Sistemi con matrice triangolare}
Supponiamo $A = \begin{bmatrix} a_{11} & \cdots & a_{1n} \\ & \ddots & \vdots \\ 0 & & a_{nn}\end{bmatrix}$ (triangolare superiore) e di voler risolvere $Ax=b$.
\noindent
In formula le componenti della soluzione sono:
\[
x_j =
\begin{cases}
\dfrac{b_n}{a_{nn}} & \text{se } n=j \\[10pt]
\dfrac{\displaystyle b_j - \sum_{i=j+1}^{n}a_{ji}x_i}{a_{jj}} & \text{altrimenti}
\end{cases}
\]
con un costo di $2(n-j)+1$ operazioni per l'entrata $j$. Sommando su tutte le entrate:
\[
\sum_{j=1}^{n}\big[2(n-j)+1\big] = n + 2\sum_{j=1}^{n}(n-j) = n+n(n-1) = n^2
\]
\[
\implies \text{risolvere un sistema triangolare costa } O(n^2) \text{ operazioni}
\]
(molto meglio del caso $A$ generica, in cui si ha $O(n^3)$).


\chapter{Metodi per sistemi lineari}

Sistemi del tipo $Ax=b$ con $A\in \mathbb{C}^{n\times n}$, $b\in \mathbb{C}^n$ e A generica.

\section{Intro}
\begin{itemize}
    \item \textbf{Metodi diretti}: metodi che, se operassimo in matematica esatta, calcolerebbero in un numero finito di passi la soluzione.
    \item \textbf{Metodi iterativi}: metodi che generano una successione $\{x_k\}_{k\in\mathbb{N}}$ tale che \[\lim_{k\to +\infty} x_k = x\]
\end{itemize}

\section{Metodi diretti}
\paragraph{Metodo di Cramer}
Prendiamolo come esempio.
Costa come calcolarsi $n+1$ determinanti, è quindi troppo costoso, si parla infatti di $O(n^4)$ o di $O(n!)$.
\subsection{Metodo di eliminazione di Gauss}
Si applica la riduzione per righe alla matrice aumentata $[A|b]$ per ottenere un sistema triangolare superiore equivalente.
Si risolve poi il sistema triangolare con la sostituzione all'indietro.
Riduzione per righe $O(\frac{2}{3}n^3)$ operazioni, + $O(n^2)$  per la sostituzione all'indietro, quindi il costo totale del metodo è $O(n^3)$.

\subsection{Fattorizzazione LU}
Fattorizzazione della matrice di partenza $A$ come prodotto delle due matrici trinagolari $L$ e $U$ (\textit{Lower} e \textit{Upper}):
\[
A = \begin{bmatrix}
a_{11} & a_{12} & \cdots & a_{1n} \\
a_{21} & a_{22} & \cdots & a_{2n} \\
\vdots & \vdots & \ddots & \vdots \\
a_{n1} & a_{n2} & \cdots & a_{nn}
\end{bmatrix}
=
\underbrace{\begin{bmatrix}
1 & & & \\
l_{21} & 1 & & \\
\vdots & & \ddots & \\
l_{n1} & l_{n2} & \cdots & 1
\end{bmatrix}}_{L}
\underbrace{\begin{bmatrix}
u_{11} & u_{12} & \cdots & u_{1n} \\
 & u_{22} & \cdots & u_{2n} \\
 & & \ddots & \vdots \\
 & & & u_{nn}
\end{bmatrix}}_{U}
\]
Se uno conosce L e U può risolvere il sistema, che diventa $LUx=b$, in 2 step:
\begin{itemize}
    \item $Lc=b$ costo $O(n^2)$
    \item $Ux=c$ costo $O(n^2)$
\end{itemize}
mentre portare inizialmente la matrice in forma triangolare costa $O(\frac{2}{3}n^3)$.


\begin{lezione}
    [Lezione 9 25/03]
\end{lezione}

\subsection{Strategie di pivoting}
\subsubsection{Pivoting Parziale}
Tramite questo tipo di strategie si mira a evitare l'uso negli algoritmi di componenti nulle o molto piccole che, rispettivamente, interrompono l'esecuzione (es. in eliminazione di Gauss un elemento nullo blocca il proseguimento) o amplificano gli errori di arrotondamento.\\
\begin{definizione}
    [Pivoting Parziale]
    Prima di ogni passo $j$ dell'algoritmo si valuta il pivot attuale ed eventualmente si scambia la sua riga, con una (indichiamola con $r$) di quelle rimanenti (= di passi successivi), se si verifica la seguente condizione:
    \[
    |a^{(j)}_{rj}| \geq \underset{i = j,...,n}{max} |a^{(j)}_{ij}|
    \]
\end{definizione}
\subsection{Fattorizzazione LU con pivoting parziale}
Si parla di pivoting parziale poiché andiamo a ricercare il pivot nella colonna corrente, andando quindi a effettuare, eventualmente, uno scambio solo tra righe (se avessimo operato anche scambi tra colonne, per arrivare al pivot migliore, si sarebbe parlato di pivoting totale).\\
Ogni scambio di riga lo rappresentiamo con la moltiplicazione a sinistra per una matrice di permutazione.
\[
\begin{array}{c}
A \to \pi_1 A \to H_1\pi_1 A \to \pi_2 H_1\pi_1 A \to \ldots \to H_{n-1}\pi_{n-1} \ldots H_1\pi_1 A \\
H_{n-1}\pi_{n-1} \ldots H_1\pi_1 A = H_{n-1}\tilde{\pi}_{n-1} \ldots H_1\pi_1 B \\
H_{n-1}\pi_{n-1} \ldots H_1\pi_1 A = U \to \text{ triangolare superiore}
\end{array}
\]
Proprietà:
\[
\pi_j H_i = \pi_j(I - N_i) = (I - \pi_j N_i)\pi_j = \tilde{H}_i \cdot \pi_j
\]
Dove $\tilde{H}_i$ è la matrice simile ad $H_i$ in cui si scambiano le righe secondo $H_i$\\
Esempio sul quaderno\\
Alla fine vogliamo scriverci una fattorizzazione del tipo $\pi A = LU$ con L triangolare inferiore e U triangolare superiore. U viene trovata alla fine dell'algoritmo, mentre L si trova applicando gli scambi di righe agli elementi sottodiagonali di L, che abbiamo trovato fino a quel punto.\\
Esempio sul quaderno
\subsection{Calcolo del determinante}
Una conseguenza del teorema Binet-Cauchy (\ref{th: Binet-Cauchy}) è che se $A=LU$ allora:
\[
det(A) = det(L)\cdot det(U) = det(U) = \prod_{j=1}^n U_{jj}
\]
dato che $det(L) = 1$.

\paragraph{Conseguenza} possiamo quindi calcolare $det(A)$ con costo pari all'eliminazione di Gauss, ovvero $O(\frac{2}{3}n^3)$.

\medskip\noindent
Quaderno 2

\section{Condizionamento di un sistema lineare}
\begin{nota}
    Il condizionamento misura quanto la soluzione x è sensibile a piccole perturbazioni nei dati in ingresso.
    \\
    Da un punto di vista più intuitivo possiamo considerare il problema $Ax=b$: risolvere questo problema significa trovare l'intersezione di iperpiani. Se gli iperpiani si intersecano con angoli "ben separati" allora il problema può essere classificato come ben condizionato, se invece i piani sono quasi paralleli allora il problema è mal condizionato: una piccola perturbazione di b può spostare moltissimo il punto di intersezione.
\end{nota}
\noindent
Con condizionamento di un sistema lineare si intende il rapporto tra errore commesso sul risultato di un calcolo e incertezza sui dati in ingresso.
In poche parole il condizionamento misura quanto la soluzione è sensibile a piccole perturbazioni dei dati.
\begin{definizione}
    [Numero di condizionamento]
    \label{def: numero di condizionamento}
    La quantità
    \[
    \mu(A) = ||A||\cdot||A^{-1}||
    \]
    si dice numero di condizionamento della matrice A rispetto alla norma $||\cdot||$.
\end{definizione}
\begin{nota}
    Le due componenti di questo prodotto indicano:
    \begin{itemize}
        \item $||A^{-1}||$ coefficiente di amplificazione degli effetti sui dati in ingresso sulla soluzione $x$ (pensare al fatto che la soluzione del sistema è definita come $x = A^{-1}b$).
        \item $||A||$ normalizza questa amplificazione, rende relativo l'errore che altrimenti sarebbe assoluto (si arriva a questa definizione dopo alcune sostituzioni ecc. effettuate a partire dalla definizione di errore relativo su $x$).
    \end{itemize}
\end{nota}
\begin{teorema}
    [Limite inferiore del numero di condizionamento]
    \label{th: mu maggiore uguale 1}
    Per ogni norma compatibile $||\cdot||$ e per ogni matrice $A$ invertibile vale
    \[
        \mu(A) \geq 1
    \]
\end{teorema}
\begin{dimostrazione}
    Per la sottomoltiplicatività della norma matriciale:
    \[
        1 = ||I|| = ||A A^{-1}|| \leq ||A|| \cdot ||A^{-1}|| = \mu(A)
    \]
\end{dimostrazione}
\begin{nota}
    Il numero di condizionamento dice: di quanto possono essere amplificati gli errori nei dati quando li trasferisci nella soluzione.
\end{nota}
\subsection{Sistema lineare mal condizionato}
\begin{definizione}
    [Sistema lineare Mal condizionato]
    Se $$\mu(A)>>1$$.
\end{definizione}
\begin{definizione}
    [Sistema lineare Ben condizionato]
    Se $$\mu(A)\approx 1$$ si dice ben condizionato.
\end{definizione}

\begin{lezione}
    [Lezione 10 05/04]
\end{lezione}
\begin{definizione}
    [Errore algoritmico nell'eliminazione di Gauss]
    Dipende dal valore assoluto dei moltiplicatori
    $
    l_{ij} = \frac{a_{ij}^{(i)}}{a_{ji}^{(i)}}
    %forse la versione corretta e':  $l_{ij} = a_{ij}^{(j)}/a_{jj}^{(j)}$
    $:
    se $|l_{ij}|$ è grande (in senso relativo) allora l'errore algoritmico è grande.
\end{definizione}
\begin{nota}
    La strategia di pivoting parziale serve a tenere sotto controllo l'errore algoritmico.
\end{nota}
\begin{teorema}
    [Regola fondamentale dell'algebra computazionale]
    Per risolvere $Ax=b$ mai fare $A^{-1}\cdot b$ perché è più instabile e costa di più di fare l'eliminazione di Gauss.
\end{teorema}
\begin{nota}
    Trovare $A^{-1}$ costa circa il doppio di realizzare la fattorizzazione LU ($O(\frac{2}{3}n^3$)), su MatLab si utilizza quindi $A \backslash b$.
\end{nota}
\section{Metodi iterativi per Ax=b}
I metodi visti fino ad ora hanno costo cubico che va bene per dimensioni fino a 10k, si passa ora all'idea di generare una successione del tipo
\begin{equation}
    \label{eq: metodo iterativo}
    x^{(k+1)} = Hx^{(k)}+c
\end{equation}
dove moltiplicare per H è conveniente (meglio di O($n^3$)) e si arriva vicino alla soluzione vera ($x^{(k)}\approx x$ per un $k$ piccolo).
\begin{definizione}
    [Metodo iterativo convergente]
    Un metodo iterativo (di punto fisso) per $Ax=b$  definito come \autoref{eq: metodo iterativo} è convergente se:
    \[
    \lim_{k\to +\infty} x^{(k)} = x \quad \forall x^{(0)}\in \mathbb{C}
    \]
\end{definizione}
\begin{teorema}
    [Convergenza di un Metodo Iterativo]
    \label{th: metod it conv}
    Un metodo iterativo è convergente se e solo se ($\Leftrightarrow$)
    \[\rho(H) < 1\]
    ovvero se $H$ è una matrice convergente.
\end{teorema}
\noindent
Osservazioni (con * indichiamo un metodo iterativo):
\begin{itemize}
    \item Se $\|H\| < 1$ per una qualsiasi norma matriciale allora $(*)$ è convergente, poiché per il teorema di Hirsh (\ref{th: Hirsh}) $\rho(H) \leq \|H\|$ .
    \item Se $|\det(H)| > 1$ allora $(*)$ non è convergente.
    \begin{dimostrazionelibera}
        Dato che il prodotto degli autovalori = determinante, allora se, in modulo, il determinante è maggiore di uno significa che c'è almeno un autovalore maggiore di uno e quindi il raggio spettrale non potrà che essere $>1$.
    \end{dimostrazionelibera}
    \item $|\det(H)| < 1$ è condizione necessaria alla convergenza ma non sufficiente.
\end{itemize}

\begin{dimostrazione}
Poiché $x$ è soluzione vera del sistema $Ax=b$, allora soddisfa il punto fisso: 
\[
x = Hx + c
\]
\begin{nota}
    Questa premessa è \uline{fondamentale}, il teorema non vale per una qualunque $H$ convergente ma solo per quelle che formano con $c$ una \textbf{coppia consistente}, ovvero è tale che il punto fisso dell'iterazione coincide con la soluzione vera (quella di $Ax=b$).
\end{nota}
\noindent
Definiamo l'\textbf{errore al passo} $k$ come $e^{(k)} = x^{(k)} - x$.
Calcoliamo come evolve l'errore da un passo al successivo:
\[
    e^{(k+1)}
    = x^{(k+1)} - x
    = \bigl(Hx^{(k)} + c\bigr) - \bigl(Hx + c\bigr)
    = H\bigl(x^{(k)} - x\bigr)
    = H e^{(k)}
\]
Applicando la relazione ricorsivamente:
\[
    e^{(k+1)} = He^{(k)} = H^2 e^{(k-1)} = \cdots = H^{k+1} e^{(0)}
\]
Quindi la convergenza del metodo equivale a:
\[
    \lim_{k \to +\infty} e^{(k)} = 0 \quad \forall\, e^{(0)}
    \;\Longleftrightarrow\;
    \lim_{k \to +\infty} H^k e^{(0)} = 0 \quad \forall\, e^{(0)}
    \;\Longleftrightarrow\;
    \lim_{k \to +\infty} H^k = 0
\]
L'ultimo passaggio vale per ogni $e^{(0)}$ se e solo se $H^k \to 0$ come matrice,
che per il \textbf{teorema della matrice convergente} (\ref{th: matrice convergente}) è equivalente a $\rho(H) < 1$.
\\
Pertanto:
\[
    \lim_{k \to +\infty} x^{(k)} = 0 \quad \forall\, x^{(0)}
    \;\Longleftrightarrow\;
    \rho(H) < 1
\]
\end{dimostrazione}

\subsection{Velocità di convergenza}
Dal teorema \ref{th: metod it conv} sappiamo che $\dfrac{||e^{(k)}||}{||e^{(0)}||} \leq ||H^k|| \quad (e^{(k)} = H^k e^{(0)})$\\
Inoltre per un risultato di algebra lineare sappiamo che:
\[
\lim_{k \to \infty} \sqrt[k]{||H^k||} = \rho(H) \implies ||H^k|| \approx \rho(H)^k
\]
Quindi se abbiamo due metodi iterativi con matrici di iterazione $H_1$ e $H_2$ tali che $$\rho(H_1) < \rho(H_2)$$ allora ci aspettiamo che il metodo più veloce a convergere sia quello con $H_1$.\\
Inoltre come stima (non rigorosa) del numero di iterazioni che servono a ottenere $||e^{(k)}|| \leq \delta$ (es. $\delta = 10^{-8}$), possiamo imporre la condizione
\[
\rho(H)^k \leq \delta \implies k \log(\rho(H)) \leq \log(\delta) \implies k \geq \frac{\log(\delta)}{\log(\rho(H))} \text{ e poi sostituire}
\]
\subsection{Criteri d'arresto}
\begin{itemize}
    \item Scegliere le iterazioni a priori se si conosce $\rho(H)$
    \item Controllare il residuo $\frac{||x^{(k+1)}-x^{(k)}||}{||b||}=\frac{||b-Ax^{(k)}||}{||b||}\leq tolleranza$
    \item $||x^{(k+1)}-x^{(k)}||< \delta$ 
\end{itemize}
\subsection{Metodi di Jacobi e Gauss-Seidel}

Alla base della costruzione di questi due metodi c'è la scomposizione del sistema generico $Ax = b$ in un sistema di punto fisso (argomento trattato più approfonditamente nel capitolo \autoref{sec: punto fisso}), ovvero quel sistema che nasce da una riscrittura di $A$ come:
\[
A = M - N 
\]
da cui
\[
Ax = b \longrightarrow (M - N)x = b \longrightarrow Mx = Nx -b
\]
Se consideriamo la $x$ a sinistra come quella del nuovo passo ($x^{(k+1)}$) allora possiamo definire dei metodi basati su questo tipo di sistemi.
\subsubsection{Definizione dei metodi di Jacobi e Gauss-Seidel}
Per la definizione dei metodi di Jacobi e Gauss-Seidel si considera A come:
\[
A = D - E - F
\]
dove:
\begin{itemize}
    \item D è la matrice diagonale con gli elementi diagonali di A
    \item E è la L invertita di segno e con diagonale nulla
    \item F è la U invertita di segno e con diagonale nulla
\end{itemize}
\begin{metodo}
    [Metodo di Jacobi]
    \[
    \begin{cases}
    x^{(k+1)} = D^{-1} (E+F)x^{(k)}+D^{-1}b \\
    x^{(0)} = \text{ random/scelto}
    \end{cases}
    \]
    che possiamo scrivere come 
    \[
    \begin{cases}
    x^{(k+1)} = H_J\cdot x^{(k)}+C_J \\
    x^{(0)} = \text{ random/scelto}
    \end{cases}
    \]
    Se lo riconduciamo alla forma generica $Mx = Nx -b$ notiamo come a sinistra abbiamo la matrice diagonale, mentre a destra una matrice quasi piena.
    Quello che ne concludiamo è che le equazioni del passo di questo metodo non faranno altro che formare un sistema diretto (incognite moltiplicate per matrice diagoanle), in cui ogni componente della nuova $x$ incognita è calcolata dal prodotto dei vari coefficienti della matrice $A$ ogni volta per l'intero vecchio vettore delle valutazioni:
    \[
    \begin{cases}
        M_{11}x_1^{(k+1)} = N_{1i} x^{(k)}+C_{J1}\\
        M_{22}x_2^{(k+1)} = N_{2i} x^{(k)}+C_{J2}\\
        \;\vdots
    \end{cases}
    \]
    Metodo che quindi, ad ogni iterazione, prende le $n$ equazioni del sistema, e da ciascuna estrae una nuova stima dell'incognita corrispondente, tutte in parallelo, tutte usando le stime vecchie.
    \begin{nota}
        Condizione \textbf{necessaria} per Jacobi è che A non abbia 0 sulla diagonale, ovvero che D sia invertibile.
    \end{nota}
\end{metodo}
\paragraph{Osservazione}
Per via del prodotto tra una matrice diagonale e una matrice avente solo 0 sulla diagonale la risultante, ovvero la matrice del metodo di Jacobi, avrà 0 sulla diagonale (caratteristica utile alla dimostrazione della convergenza).

\paragraph{Equazione del passo di Jacobi}
\[
x_i^{(k+1)} = \frac{1}{a_{ii}} \left( b_i - \sum_{\substack{j=1 \\ j \neq i}}^{n} a_{ij} x_j^{(k)} \right), \quad \text{per } i = 1, \dots, n
\]
\begin{nota}
    Il segno $"-"$ qui è dovuto al fatto che $E$ ed $F$ sono definite con segno invertito rispetto ad $A$.
\end{nota}
\begin{nota}
    Non si può sovrascrivere/calcellare nessuna componente di $x^{(k)}$ finchè non ho calcolato tutte le entrate di $x^{(k+1)}$ (\textbf{notare come nella sommatoria al variare di $j$ compaia il termine $x_j$, servono dunque tutte le componenti del vettore $x$ per il calcolo di ogni componente del nuovo vettore $x$}).
    Questo è da tenere in conto nei casi di implementazioni con matrici molto grandi in quanto sarà necessaria molta memoria per poter contenere contemporaneamente tutte le entrate di 2 passi dell'algoritmo.
\end{nota}
\begin{metodo}
    [Metodo di Gauss-Seidel]
    \[
    \begin{cases}
    x^{(k+1)} = (D-E)^{-1}F x^{(k)}+(D-E)^{-1}b \\
    x^{(0)} = \text{ random/scelto}
    \end{cases}
    \]
    che possiamo scrivere come 
    \[
    \begin{cases}
    x^{(k+1)} = H_{GS}\cdot x^{(k)}+C_{GS} \\
    x^{(0)} = \text{ random/scelto}
    \end{cases}
    \]
    Questa volta l'utilizzo delle valutazioni precedenti (passo $k$) è diverso: come $M$ abbiamo una matrice triangolare inferiore e come $N$ una triangolare superiore mancante di diagonale.\\
    Quello che ne risulta è che ad ogni iterazione il sistema che si viene a formare è composto da equazioni che, una dopo l'altra, vedono aumentare i termini di $x^{(k+1)}$ (stiamo "scendendo" le righe di $M$, triangolare inferiore, ma solo dopo aver calcolato la stima della componente precedente) e diminuire quelli di $x^{(k)}$ ("scendendo" le righe di $N$, triangolare superiore).
    
    La differenza con Jacobi è che non appena viene calcolata una nuova stima $x_i^{(k+1)}$, questa viene utilizzata \textbf{immediatamente} per calcolare $x_{i+1}^{(k+1)}$, senza attendere il completamento dell'intero passo \textit{k+1esimo}.\\
    Invece di aggiornare tutte le componenti in parallelo, l'aggiornamento
    avviene \textbf{in serie}, sfruttando le informazioni più fresche già
    disponibili. Le stime risultano quindi più accurate già durante lo stesso
    passo, e il metodo converge tipicamente più velocemente di Jacobi.
    \\
    Il rovescio della medaglia è che il metodo non è parallelizzabile e richiede
    di processare le incognite in ordine.
\end{metodo}

\paragraph{Equazione del passo di Gauss-Seidel}
\[
x_i^{(k+1)} = \frac{1}{a_{ii}} \left( b_i - \sum_{j=1}^{i-1} a_{ij} x_j^{(k+1)} - \sum_{j=i+1}^{n} a_{ij} x_j^{(k)} \right), \quad \text{per } i = 1, \dots, n
\]
\begin{nota}
    A differenza di Jacobi posso sovrascrivere in memoria $x_j^{(k)}$ una volta che ho calcolato $x_j^{(k+1)}$ in quanto il calcolo del passo successivo della componente i-esima richiede del passo precedente solo le componenti successive a i. Questo è utile per la gestione della memoria.
\end{nota}

\begin{teorema}
    [Teorema convergenza Jacobi/Gauss-Seidel]
    Se $A\in \mathbb{C}^{n\times n}$ soddisfa una di queste due condizioni:
    \begin{enumerate}
        \item A a predominanza diagonale forte
        \item A irriducibile $\land$ a predominanza diagonale debole
    \end{enumerate}
    allora sia Jacobi che Gauss-Seidel per un sistema lineare risultano \textbf{convergenti}.
\end{teorema}

\begin{lezione}
    [Lezione 11 08/04]  
\end{lezione}

\begin{dimostrazione}   
    Dobbiamo far vedere che $\rho(H_J)<1$ (nel caso 1 del teorema).
    \[
    H_J = 
    \begin{cases}
    \dfrac{a_{ij}}{a_{ii}} & i\neq j\\[6pt]
    0 & i=j
    \end{cases}
    \]
    \[
    \|H_J\|_\infty = \max_{i=1,\dots,n}\sum_{\substack{j=1\\ j\neq i}}^n \left|\frac{a_{ij}}{a_{ii}}\right| = \max_{i=1,\dots,n} \frac{\sum_{j\neq i}|a_{ij}|}{|a_{ii}|} < 1
    \]
    \[
    \implies \quad 1 > \|H_J\|_\infty \geq \rho(H_J) \quad \implies \quad H_J \text{ convergente}
    \]
    Se siamo nel caso 2. $\|H_J\|_\infty\leq1$, però sappiamo che $\exists$ un cerchio di Gershgorin in $H_J$ che ha raggio $<1$ (il \autoref{th: terzo cor gersh} richiede almeno un indice con disuguaglianza stretta).\\
    Quindi se guardiamo i cerchi di Gershgorin di $H_J$, questi hanno tutti centro in $0$ (ricordarsi della forma della matrice di J.) e ce n'è almeno $1$ con raggio $<1$ e tutti gli altri con raggio $\leq1$ (per Terzo Teorema di Gershgorin).

    \medskip
    \noindent
    $A$ irriducibile $\Rightarrow H_J$ è irriducibile $\overset{3^\circ\text{ Gersh.}}{\implies}$ se ho un autovalore sul bordo dell'unione dei cerchi, deve stare sul bordo di ogni cerchio $\Rightarrow$ non ci sono autovalori di modulo $1$, ed in particolare sono tutti minori di $1$ $\Rightarrow \rho(H_J)<1$.

    \begin{nota}
        [Osservazione]
        Le assunzioni 1) e 2) implicano che $a_{ii}\neq0\ \forall i$.
    \end{nota}

    \medskip
    \noindent
    \textbf{Dim G.S.} supponiamo per assurdo che $H_{GS} = (D - E)^{-1} F$ abbia un autovalore $\lambda \in \mathbb{C}$ tale che $|\lambda| \geq 1$, allora, dalla definizione di autovalore:
    \begin{align*}
        0 &= \det(\lambda I - H_{GS}) = \det\left(\lambda I - (D-E)^{-1}F\right) \\
        &= \det\left((D-E)^{-1}\left[\lambda(D-E) - F\right]\right) \\
        &= \det(D-E)^{-1} \cdot \det\left(\lambda(D-E) - F\right) \\
        &= \det(D-E)^{-1} \cdot \lambda^n \det\left(D - E - \lambda^{-1}F\right) \\
        &\underbrace{=}_{\neq 0} \quad \underbrace{\lambda^n}_{\neq 0} \det\left(D - E - \lambda^{-1}F\right)
    \end{align*}
    \noindent
    Dimostriamo che questo (*) è assurdo:
    \[
    D - E - \lambda^{-1}F = \begin{bmatrix}
        a_{11} & a_{12} & \cdots & a_{1n} \\
        \vdots & & & \vdots \\
        a_{n1} & & \cdots & a_{nn}
    \end{bmatrix}
    \]
    se guardiamo ai cerchi di Gershgorin di $M_\lambda$.\\
    Questi hanno gli stessi centri di quelli di A e raggi minori o uguali. In particolare $A$ è a dominanza diagonale forte  $\Rightarrow M_\lambda$ dominanza diagonale debole.\\
    Ricapitolando:
    \begin{itemize}
        \item A dominanza diagonale debole $\to M_\lambda$ dom diag debole
        \item A irriducibile $\to M_\lambda$ irriducibile (ha zero nelle stesse posizioni delle fuori diag)
    \end{itemize}
    Quindi $M_\lambda$ non è singolare per i teoremi di Gershgorin e questo contraddice *.\\
    Da questa contraddizione concludiamo che $H_{GS}$ non ha autovalore di modulo $\geq 1$.
\end{dimostrazione}
    
\begin{definizione}
    [Matrice Tridiagonale]
    Matrice che ha elementi diversi da 0 solo sulla diagonale principale, sulla prima sottodiagonale e sulla prima sopradiagonale.
\end{definizione}

\begin{teorema}
    [Matrici tridiagonali]
    Se $A\in \mathbb{C}^{n\times n}$ è tridiagonale allora 
    \[
    \rho(H_{GS}) = \rho(H_J)^2
    \]
    in particolare convergono entrambi o divergono entrambi.
\end{teorema}
\begin{nota}
    Nel caso tridiagonale, se convergono, G-S più veloce.
\end{nota}

\section{Sistemi lineari sovradeterminati e approssimazione ai minimi quadrati}
Sistemi $Ax=b$ dove $A\in\mathbb{C}^{m\times n}$ e $b\in\mathbb{C}^m$ con:
\[
m>n
\]
Dato che spesso la soluzione non esiste, infatti è normale che $x$ soddisfi contemporaneamente più condizioni incompatibili tra loro, allora è di nostro interesse minimizzare il vettore residuo definito come $r(x) = b - Ax$; si sostituisce quindi il problema $Ax=b$ con:
\begin{equation}
\label{eq: min quad}
\mathcal{P}: \underset{x\in \mathbb{C}}{min}||r(x)||_2 = \underset{x\in \mathbb{C}}{min}||b-Ax||_2
\end{equation}
Indichiamo con $x^*$ il punto minimo di $\mathcal{P}$.
Per trovare questo punto, ovvero la soluzione ottima del problema, ci restringiamo al caso $A \in \mathbb{R}^{m \times n}$ e $b \in \mathbb{R}^m$.
Per cui il nostro problema risulta essere:
\[
\min_{x \in \mathbb{R}^n} \|Ax - b\|_2
\quad
\textit{dove:}
\s\s
\|b - Ax\|_2 = \sqrt{\sum_{j=1}^{m} r_j(x)^2}
\]
con $r_j$ vettore residuo.\\
Prima semplificazione: cerco il minimo di $\|b - Ax\|_2^2$:
\[
\psi(x) = \|b - Ax\|_2^2 = (b - Ax)^T (b - Ax) :
\mathbb{R}^n \to \mathbb{R}.
\]
Funzione convessa $\Rightarrow$ il suo punto minimo $x^\star$ verifica $\nabla \psi(x^\star) = 0 \in \mathbb{R}^n$ con $0$ vettore.
\[
\boxed{r(x) = b - Ax} = [r_1(x), \ldots, r_m(x)],
\qquad
\psi(x) = \sum_{i=1}^{m} [r_i(x)]^2
\]

\[
[r_i(x)]^2 =
\left[
b_i - \sum_{j=1}^{n} a_{ij} x_j
\right]^2
\]

\[
\frac{\partial r_i^2(x)}{\partial x_s}
= 2 \left[
b_i - \sum_{j=1}^{n} a_{ij} x_j
\right] (-a_{is})
= 2 \left[
\sum_{j=1}^{n} a_{is} a_{ij} x_j - b_i a_{is}
\right]
\]

\[
\frac{\partial \psi(x)}{\partial x_s}
= 2 \left[
\sum_{i=1}^{m} \sum_{j=1}^{n} a_{is} a_{ij} x_j
- \sum_{i=1}^{m} b_i a_{is}
\right]
\]

\[
\textit{Osservazione: } \sum_{i=1}^{m} b_i a_{is} = (A^T b)_s
\]
\[
\Rightarrow \sum_{i=1}^{m} a_{is} \sum_{j=1}^{n} a_{ij} x_j
= \sum_{i=1}^{m} a_{is} (A x)_i
= (A^T A x)_s
\]

\[
\frac{\partial \psi(x)}{\partial x_s}
= 2 \left( (A^T A x)_s - (A^T b)_s \right)
\]

\[
\Rightarrow \nabla \psi(x) = 2 (A^T A x - A^T b)
\]
In definitiva abbiamo trovato che:
\[
\nabla \psi(x) = 0
\;\Longleftrightarrow\;
A^T A x - A^T b = 0
\;\Longleftrightarrow\;
A^T A x = A^T b
\]
dove l'ultima equazione descrive un sistema lineare quadrato di dimensione $n \times n$.\\
Per cui, \textbf{per trovare $x^\star$ \`e sufficiente risolvere
un sistema lineare quadrato}.

\begin{definizione}
    [Sistema delle equazioni normali]
    \label{def: eq normali}
    Viene definito come tale il sistema:
    \[
    A^TAx=A^Tb
    \]
\end{definizione}
\begin{definizione}
    [Problema lineare ai minimi quadrati]
    Il problema $\mathcal{P}$ (\ref{eq: min quad}) viene detto \textit{Problema Lineare ai Minimi Quadrati}.
    Con questo nome si intende quindi una classe di problemi in cui si vuole trovare la soluzione che meglio approssima un sistema lineare incompatibile, minimizzando la somma dei quadrati degli errori.
\end{definizione}

\begin{teorema}
    [Teorema Sistema delle equazioni normali]
    Se $A\in \mathbb{R}^{m\times n}$ con $m>n$, allora il \textit{Sistema delle equazioni normali} \textbf{ammette} soluzione \[\forall b \in \mathbb{R}^n\] 
    Questa soluzione è unica \textbf{se e solo se} 
    \[
    rk(A)=n=max
    \]
\end{teorema}
Pertanto il sistema delle equazioni normali ha \textbf{sempre} almeno una soluzione e se $rank(A) = n$ allora questa è \textbf{unica}.

\begin{dimostrazione}
\textbf{Esistenza.}
Per il teorema di Rouché-Capelli, il sistema $A^TAx = A^Tb$ ammette soluzione
se e solo se $A^Tb \in \operatorname{Im}(A^TA)$. Mostriamo che questa condizione
è sempre verificata.
\\
Poiché $\mathbb{R}^m = \operatorname{Im}(A) \oplus \operatorname{Im}(A)^\perp$,
ogni vettore $b \in \mathbb{R}^m$ si decompone univocamente come
\[
    b = b_1 + b_2,
    \qquad b_1 \in \operatorname{Im}(A),\quad b_2 \in \operatorname{Im}(A)^\perp
\]
dove $\operatorname{Im}(A)^\perp = \{y \in \mathbb{R}^m : y^T z = 0\ \forall z \in \operatorname{Im}(A)\}$
e
$\operatorname{Im}(A) = \{y = Az\} $
Quindi:
\[
    A^T b = A^T b_1 + A^T b_2
\]
Analizziamo i due termini separatamente.
\\
\textit{Primo termine.} Poiché $b_1 \in \operatorname{Im}(A)$, esiste $z \in
\mathbb{R}^n$ tale che $b_1 = Az$, quindi
\[
    A^T b_1 = A^T A z \in \operatorname{Im}(A^T A)
\]
\\
\textit{Secondo termine.} Poiché $b_2 \in \operatorname{Im}(A)^\perp$, risulta
$b_2 \perp a_j$ per ogni colonna $a_j$ di $A$, ovvero
$\langle a_j, b_2 \rangle = 0$ per ogni $j = 1,\dots,n$. Quindi:
\[
    A^T b_2 =
    \begin{bmatrix} \langle a_1, b_2 \rangle \\ \vdots \\ \langle a_n, b_2 \rangle
    \end{bmatrix} = 0
\]
In conclusione $A^T b = A^T A z + 0 = A^T A z \in \operatorname{Im}(A^TA)$,
e l'esistenza è dimostrata.

\bigskip
\noindent\textbf{Unicità.}
La soluzione è unica se e solo se $\det(A^T A) \neq 0$.
Studiamo il segno di $\det(A^T A)$ tramite la formula di Binet-Cauchy (versione
generalizzata): per $A \in \mathbb{R}^{m \times n}$ con $m > n$,
\[
    \det(A^T A)
    = \sum_{\mathcal{I}} \det\!\bigl(A^T_{\mathcal{I}}\bigr)\cdot\det(A_{\mathcal{I}})
\]
dove la somma è estesa a tutti i sottoinsiemi $\mathcal{I} \subseteq \{1,\dots,m\}$
con $|\mathcal{I}|=n$, e $A_{\mathcal{I}} \in \mathbb{R}^{n\times n}$ è il minore
di ordine $n$ ottenuto selezionando le righe di indice $\mathcal{I}$.
\\
Poiché $A^T_{\mathcal{I}} = (A_{\mathcal{I}})^T$, si ha
$\det(A^T_{\mathcal{I}}) = \det(A_{\mathcal{I}})$, quindi:
\[
    \det(A^T A) = \sum_{\mathcal{I}} \det(A_{\mathcal{I}})^2 \geq 0
\]
Questa somma è strettamente positiva se e solo se esiste almeno un minore
$A_{\mathcal{I}}$ non singolare, cioè:
\[
    \det(A^T A) \neq 0
    \;\Longleftrightarrow\;
    \exists\, \mathcal{I} : \det(A_{\mathcal{I}}) \neq 0
    \;\Longleftrightarrow\;
    \operatorname{rk}(A) = n
\]
\end{dimostrazione}

\begin{nota}
    Quanto visto si traduce nel caso complesso con $A^H$ al posto di $A^T$
\end{nota}

\begin{lezione}
    [Lezione 12 12/04]
\end{lezione}

\begin{definizione}
    [Sistema delle equazioni normali nel caso complesso]
    \[
    A^HAx=A^Hb
    \]
\end{definizione}
\subsection{Mal condizionamento delle equazioni normali}
In molti casi la matrice $A^TA$ risulta mal condizionata. Nel caso quadrato $m = n$ si ha:
\[
\mu(A^TA) = \mu(A^2) = \mu(A)^2
\]
Questo significa che il condizionamento si \textbf{eleva al quadrato}: se $\mu(A) = 10^6$, risolvere le equazioni normali comporta un condizionamento di $10^{12}$, rendendo la soluzione numericamente inaffidabile.

\section{Fattorizzazione QR}
Metodo più sicuro dal punto di vista del condizionamento rispetto a quello appena trattato (\textit{Equazioni Normali}): non si moltiplica la matrice $A$, ovvero i coefficienti del problema che vogliamo risolvere, per se stessa e quindi portando al quadrato eventuali errori; ma si trasforma essa in una forma comoda da risolvere tramite operazioni che non alterano la sensibilità del problema agli errori (moltiplichiamo per matrici $H$ unitarie, che quindi non modificano la norma di $A$ e quindi non hanno effetto sul numero di condizionamento (\autoref{def: numero di condizionamento})).\\
In particolare avere $Q$ unitaria ci garantisce di non cambiare la norma di $A$.
\begin{definizione}
    [Fattorizzazione QR]
    Data $A\in \mathbb{C}^{n\times n}$ una (mai unica) fattorizzazione QR di A è data da 
    \begin{itemize}
        \item $Q\in \mathbb{C}^{m\times m}$ unitaria
        \item $R\in \mathbb{C}^{m\times n}$ triangolare superiore
    \end{itemize}
    In particolare R non è quadrata e può essere vista come due blocchi 
    $
    \left[
    \begin{matrix}
        R_1 \\ 
        0
    \end{matrix}
    \right]
    $
    con $R_1$ triangolare superiore.
\end{definizione}
\begin{teorema}
    [Eistenza della fattorizzazione QR]
    Data $A\in \mathbb{C}^{n\times n}$ \textbf{esiste sempre} una fattorizzazione QR di A, in particolare non è mai unica.
\end{teorema}
\begin{nota}
    [Proprietà algebrica]
    moltiplicare un vettore per una matrice unitaria non cambia la sua norma euclieda (norma due).
\end{nota}
\subsection{Proprietà della fattorizzazione QR}
\begin{enumerate}
    \item \textbf{Importante}: $x \in \mathbb{C}^m$, $Q \in \mathbb{C}^{m\times n}$, Q unitaria
    \[
    \boxed{||x||_2 = ||Qx||_2} \xrightarrow{dimostr.} ||x||_2^2 = x^H x \iff ||Qx||_2^2 = (Qx)^H(Qx) = x^H Q^H Qx = x^H x
    \]
    quindi moltiplicare un vettore per una matrice unitaria non ne cambia la norma.
    \begin{nota}
        Il penultimo passaggio sfrutta \autoref{note: tras con prod}.
        Vale anche per la norma della matrice inversa, stessa dimostrazione.
    \end{nota}
    \item $A \in \mathbb{C}^{m\times n}$, $Q \in \mathbb{C}^{m\times n}$
    \[
    ||QA||_F = ||A||_F \text{ dim su quad}
    \]
    \item $A \in \mathbb{C}^{m\times m}$
    \[
    ||QA||_2 = ||A||_2 \text{ dim su quad}
    \]
\end{enumerate}
\subsection{Metodo QR per problemi ai minimi quadrati}
\begin{metodo}[Metodo QR per minimi quadrati]
Dati $A \in \mathbb{C}^{m\times n}$, $b \in \mathbb{C}^m$ risolvere
\[
\underset{x\in \mathbb{C}}{min}||Ax-b||_2 
\]
è equivalente a risolvere
\[
\underset{x\in \mathbb{C}}{min}||R_1x-c_1||_2^2+||c_2||_2^2 
\]
come dimostrato dai seguenti passaggi, dove si sfrutta $Q^HQ= QQ^H = I$ (=unitaria):
\[
||Ax - b||_2^2 = ||\underbrace{QRx - b}_{\text{vettore in } \mathbb{C}^m}||_2^2 = ||Q^H(QRx - b)||_2^2 = ||Rx - \underbrace{Q^H b}_{c\,\in\mathbb{C}^m}||_2^2 = ||Rx - c||_2^2
\]
\begin{nota}
    Moltiplico per $Q^H$ in quanto essa è unitaria e non cambia il valore della norma.\\
    Lo facciamo perché è nostro interesse arrivare a qualcosa che coinvolga soltanto $R$ in quanto triangolare superiore.
\end{nota}
\noindent
Scrivendolo a blocchi:
\[
\left\|\begin{bmatrix} R_1 \\ 0 \end{bmatrix} x - \begin{bmatrix} c_1 \\ c_2 \end{bmatrix}\right\|_2^2 = \left\|\begin{matrix} R_1 x - c_1 \\ c_2 \end{matrix}\right\|_2^2 = ||R_1 x - c_1||_2^2 + ||c_2||_2^2
\]
Dunque il nostro problema $\mathcal{P}$ è equivalente a:
\[
\underset{x\in\mathbf{C}^n}{min}\;||R_1 x - c_1||_2^2 + ||c_2||_2^2
\]
Dove $c_1$ è formato dalle prime $n$ componenti di $c$ e $c_2$ dalle rimanenti $m-n$.
\begin{nota}
    Nel valore ottimo del residuo, $c_2$ non dipende da x quindi non abbiamo il controllo sul secondo termine.
\end{nota}
Nel caso di rango massimo la soluzione del problema classico è quella del sistema $R_1x-c_1$: infatti se $A$ ha rango massimo allora anche $R$ ha rango massimo perché $R=Q^HA$ e quindi $R_1$ è invertibile (colonne di $R_1$ linearmente indipendenti).
\subsubsection{Recap metodo QR}
\begin{enumerate}
    \item calcolo $A=QR$ (costo $O(mn^2)$)
    \item calcolo $c=Q^Hb$ (costo $O(mn)$)
    \item risolvo $R_1x=c_1$ (costo $O(n^2)$)
    \item calcolo $||c_2||_2$ (costo $O(m-n)$)
\end{enumerate}
\begin{nota}
    $||c_2||_2$ corrisponde al valore ottimo del problema ai minimi quadrati.
\end{nota}
\end{metodo}
\subsection{Come si calcola la fattorizzazione QR}
Osservazione di partenza: il prodotto di 2 matrici unitarie è unitario.
\medbreak
\noindent
L'idea è quella di moltiplicare A a sx con delle matrici in modo da portarla in forma R, per farlo si usano le \textbf{matrici di Householder}.\\
Il ruolo di queste trasformazioni è azzerare gli elementi sotto la diagonale senza alterare (e quindi peggiorare) gli errori numerici.
La loro \textbf{peculiarità} è il riuscire ad \textbf{azzerare molteplici elementi "in un colpo solo"}.
\begin{nota}
    Nella fattorizzazione LU, per esempio, l'azzeramento degli elementi di una colonna invece avveniva in $j-n$ operazioni.
\end{nota}
Per ottenere un qualcosa che azzeri gli elementi senza toccare le norme ci serve uno strumento che "schiaccia" tutta la norma di un vettore sulla prima componente, introduciamo quindi le:
\begin{definizione}
    [Trasformazione di Householder]
    Dato $v\in\mathbb{C}^m$ si definisce la trasformazione di Householder associata a $v$ la matrice
    \[
    H_v = I -  \frac{2}{||v||^2_2}\cdot vv^H
    \]
    \begin{nota}
        In particolare H è della forma "identità + rango 1".
    \end{nota}
    \noindent
    I componenti della matrice sono così pensati:
    \begin{itemize}
        \item $I-...$ : dato che la matrice $H$ moltiplica un vettore colonna $x$ per annullare le sue componenti sotto la prima, strutturiamo la trasformazione in modo da sottrarre allo stesso (da qui la $I-...$) una quantità ad hoc.
        \item $vv^H$: serve a costruire (una volta moltiplicato per $x$) la proiezione di $x$ lungo $v$.
        In particolare l'espressione finale $vv^Hx$  è un vettore sempre parallelo a $v$, scalato proporzionalmente a quanto $x$ "punta" nella direzione di $v$. 
        \item $\frac{2}{||v||^2_2}$: è uno scalare volto a normalizzare la proiezione che altrimenti dipenderebbe dalla lunghezza di $v$, arbitraria per definizione (si veda dopo per come scegliere $v$).
        In particolare la moltiplicazione per due, che si traduce con una doppia sottrazione, serve a riprodurre una riflessione (invece di una semplice proiezione)
    \end{itemize}
\end{definizione}
\noindent
Le proprietà di questa matrice sono quindi:
\begin{itemize}
    \item $H_v$ unitaria $\forall v$
    \item dato $x\in\mathbb{C}^m \quad \exists v\in\mathbb{C}^m : H_v\cdot x = 
    \left[
    \begin{array}{c}
         *\\[-10pt]
         0\\[-10pt]
         \vdots\\[-10pt]
         0\\
    \end{array} 
    \right]
    = \alpha \cdot e_1
    $ 
    dove $e_1$ è un vettore unitario, in particolare è il  primo vettore della base canonica.\\
    Osservazione: dato che moltiplicare per matrici unitarie preserva la norma dei vettori, necessariamente $|\alpha| = ||x||_2$
    \item \uline{\textbf{Come scegliere} $v$}: Per ottenere la trasformazione di Householder che manda $x$ in un multiplo di
    $
    e_1 =
    \begin{bmatrix}
    1 \\[-10pt] 0 \\[-10pt] \vdots \\[-10pt] 0
    \end{bmatrix},
    $
    con $x$ che nelle iterazioni del metodo QR è il vettore colonna parziale che parte dall'elemento $k$ (= futuro pivot) e arriva all'ultimo della colonna.\\
    Si pu\`o scegliere 
    $
    \begin{cases}
    v = x + \theta \, \|x\|_2 \, e_1\\
    |\theta| = 1
    \end{cases}
     $
    ad esempio:
    \[
    v = x \pm \|x\|_2 \, e_1
    =
    \begin{bmatrix}
    x_1 \pm \|x\|_2 \\[-10pt]
    x_2 \\[-10pt]
    \vdots \\[-10pt]
    x_m
    \end{bmatrix}
    \]
    \textbf{Esempio}: preso $x = \begin{bmatrix} 3 \\[-10pt] 4 \end{bmatrix}$, si ha $\|x\|_2 = \sqrt{3^2+4^2} = 5$, quindi scegliendo il segno $+$:
    \[
    v = \begin{bmatrix} 3 \\[-10pt] 4 \end{bmatrix} + 5 \begin{bmatrix} 1 \\[-10pt] 0 \end{bmatrix} = \begin{bmatrix} 8 \\[-10pt] 4 \end{bmatrix}
    \]
    Si verifica: $v^Tv = 80$, $\;v^Tx = 40$, quindi
    \[
    Hx = x - 2\frac{v^Tx}{v^Tv}\,v = \begin{bmatrix}3\\[-10pt]4\end{bmatrix} - \begin{bmatrix}8\\[-10pt]4\end{bmatrix} = \begin{bmatrix}-5\\[-10pt]0\end{bmatrix} = -\|x\|_2\,e_1 \checkmark
    \]
    Quindi si può calcolare $H_v$ "facilmente",
    basta memorizzare $v$,
    che si trova con $O(m)$ operazioni.
    \begin{nota}
        L'idea di questa ultima proprietà descritta è di trovare v come vettore normale al piano di riflessione, dato che la matrice di Householder rappresenta una riflessione rispetto a un iperpiano. 
    \end{nota}
\end{itemize}
\begin{nota}
    [Importante]
    Notare come $v$ sia uguale a $x$ tranne che per il primo elemento, che differisce per il modulo $2$ del valore di $x$.
\end{nota}
\begin{nota}
    [Importante]
    Occhio a tenere sempre a mente che la matrice $H$ azzera la parte voluta del vettore voluto ma ha anche effetto sulle altre colonne della matrice a cui viene moltiplicata, per esempio se consideriamo una matrice $A$ che ha come prima colonna il vettore dell'esempio precedente:
    \[
    A = \begin{bmatrix} 3 & 1 \\[-10pt] 4 & 2 \end{bmatrix}
    \]
    vediamo come la moltiplicazione per la stessa $H$ trovata prima ha effetti anche sulla seconda colonna:
    \[
    H = I - \frac{1}{40}\begin{bmatrix}8\\[-10pt]4\end{bmatrix}\begin{bmatrix}8&4\end{bmatrix} = \begin{bmatrix}-0.6&-0.8\\[-10pt]-0.8&0.6\end{bmatrix}
    \]
    \[
    H\begin{bmatrix}1\\[-10pt]2\end{bmatrix} = \begin{bmatrix}-0.6&-0.8\\[-10pt]-0.8&0.6\end{bmatrix}\begin{bmatrix}1\\[-10pt]2\end{bmatrix} = \begin{bmatrix}-0.6-1.6\\[-10pt]-0.8+1.2\end{bmatrix} = \begin{bmatrix}-2.2\\[-10pt]0.4\end{bmatrix}
    \]

    Quindi:
    \[
    HA = \begin{bmatrix}-5 & -2.2 \\[-10pt] 0 & 0.4\end{bmatrix}
    \]
\end{nota}
Una volta ottenuta la forma 
\[
H_n...H_1 A=R
\]
se si moltiplica a sx per entrambi i membri, in modo da annullare tutte le matrici H di A, si ottiene una forma del tipo
\[
A = H^H_1...H^H_n R
\]
per via della proprietà di essere unitarie delle matrici $H$.
Da questa forma ricaviamo che, dalla definizione di QR, Q è:
\[
Q = H^H_1...H^H_n
\]
\begin{nota}
    Q risulta unitaria come da definizione perché prodotto di unitarie.
\end{nota}
\begin{nota}
    [Costo]
    il costo computazionale del calcolo di QR è $O(mn^2)$.
\end{nota}
\begin{nota}
    Nel caso $m=n$ si può considerare la fattorizzazione QR anche per risolvere $Ax=b$.
\end{nota}

\begin{lezione}
    [Lezione 13 15/04]
\end{lezione}

\chapter{Approssimazione di funzioni}

\section{Interpolazione}

\begin{definizione}
    [Problema di Interpolazione]
    Trovare una funzione polinomiale $P_k(x)$ tale che $P_j(x) = y_j$.
\end{definizione}
\noindent
Pertanto costruire un interpolante significa trovare una funzione che passa \textbf{esattamente} per i punti dati. 
\begin{figure}[H]
    \centering
    \includegraphics[width=0.35\linewidth]{img/funz interpol.png}
\end{figure}
\begin{definizione}
    [Interpolazione Polinomiale (o parabolica)]
    Dati $k$ punti distinti e $k$ valori di $f(x)$ vogliamo trovare $P_k(x)$ polinomio di grado al più $k$ tale che $P_j(x) = y_j \quad \forall j=1,...,k$\\
    In particolare il polinomio è definito come
    \[
    P_k(x) = \sum^k_{j=0} a_j x^j
    \]
    dove i coefficienti $a_i$ sono trovati risolvendo il sistema:
    \[
    \begin{cases}
        P(x_0) = a_0+a_1 x_0 + ... + a_k x_0^k = y_0 \\
        .\\
        .\\
        .\\
        P(x_k) = a_0+a_1 x_k + ... + a_k x_k^k = y_k 
    \end{cases}
    \]
\end{definizione}
\begin{definizione}
    [Matrice di Vandermonde]
    Il sistema sopra descritto può essere espresso tramite la matrice di Vandermonde definita come
    \[
    V = 
    \left[
    \renewcommand{\arraystretch}{0.8}
    \begin{matrix}
        1 & x_0 & ... & x_0^k\\
        1 & . & .& .\\
        1 & . & .& .\\
        1 & . & .& .\\
        1 & x_k & ... & x_k^k
    \end{matrix}
    \right]
    \]
    Il problema quindi diventerà 
    \[
    Va=y
    \]
\end{definizione}
\begin{teorema}
    [Esistenza e unicità del polinomio di interpolazione]
    \label{th: esistenza unicita interpolazione}
    Il \underline{polinomio di interpolazione} sui punti $(x_0, y_0)...(x_k, y_k)$, di grado al più $k$, \underline{esiste ed è unico} se e solo se il sistema $Va=y$ ha un'unica soluzione.
\end{teorema}
\noindent
Questo è un risultato diretto del Teorema di Rouché-Capelli (\autoref{th: Rouche-Capelli}) applicato al sistema $Va=y$: essendo $V$ quadrata, \uline{la soluzione esiste ed è unica} se e solo se $det(V)\neq 0$, ovvero $V$ invertibile (come dimostrato dal prossimo teorema (\autoref{th: determinante Vandermonde})).\\
L'esistenza e unicità del polinomio di interpolazione è il \textbf{concetto su cui si basa} tutta \textbf{la teoria} che ne consegue.
\begin{teorema}
    [Determinante della matrice di Vandermonde]
    \label{th: determinante Vandermonde}
    Il determinante della matrice di Vandermonde è:
    \[
    det(V) = \underset{k\ge \,j\,>\,i\,\ge0}{\prod}\left( x_j - x_i\right)
    \]
    Se $x_i \neq x_j \;\forall i \neq j \Longrightarrow det(V)\neq 0 \Longrightarrow$ la soluzione di $Va=y$ esiste ed è unica.
\end{teorema}
\begin{nota}
    Quanto appena affermato risulta vero perché se prendiamo $x$ distinte non importa quanto siano astruse le $y$, trovo sempre e comunque un polinomio di quel grado che le interpola.
\end{nota}
\paragraph{Commento}
La matrice di Vandermonde ha un condizionamento, ovvero una sensibilità ai piccoli cambiamenti, che cresce molto rapidamente con l'aumentare dei nodi di interpolazione.
\paragraph{Conclusioni}
Non conviene risolvere $Va=y$ per trovare il polinomio quando $k$ è moderato ($k\ge 10$), significa che se si hanno piccole variazioni sui valori di $y_i$ si ottengono relativemente grandi variazioni sui coefficienti $a_j$ calcolati.
\paragraph{Osservazione}
\uline{Non è detto che il polinomio abbia grado esattamente k}, questo dipende infatti dai valori di $y_j$.
Avere grado $k$ significa che $a_k\ne 0$ e questo dipende dal valore della soluzione di $Va=y$.

\section{Rappresentazione del polinomio di interpolazione}
Quello fatto prima è stato cercare di scrivere $P_k(x)$ come combinazione lineare delle funzioni $1,x,x^2,...,x^k$, ovvero stiamo considerando lo spazio dei polinomi di grado al più $k$ con la base $\{1,x,x^2,...,x^k\}$.\\
Quello di trovare le coordinate di $P_k(x)$ rispetto alla base $B$, dove $B$ è la base dei monomi $(\{1,x,x^2,...,x^k\}$), ovvero quella utilizzata nella matrice di Vandermonde, è però un problema mal condizionato.\\
Il problema ben condizionato che possiamo ritrovare per il nostro scopo è quello ottenuto potendo scegliere $B$, cioè vogliamo trovare $B= \{q_0(x), \dots, q_k(x)\}$ base dello spazio dei polinomi di grado $\le k$ tale per cui si trova in maniera conveniente la rappresentazione
\[
P_k(x) = \sum^k_{j=0} c_j\, q_j(x)
\]
\begin{nota}
    [\textbf{Importante}]
    Il polinomio di interpolazione è \textbf{unico}, sono diversi invece i modi in cui lo posso scrivere: stiamo infatti considerando basi diverse dello stesso spazio dei polinomi.
\end{nota}
\noindent
Alla luce dell'unicità e della possibilità di utilizzare base diverse procediamo con l'introdurre delle basi funzionali alla risoluzione del problema.
\subsection{Interpolazione nella base di Lagrange}
Idea di base: trovare un polinomio con coefficienti $c_j$ facili da trovare, ad esempio proprio gli $y_j$.
Per ottenere questo risultato nella sommatoria che descrive il polinomio abbiamo bisogno di funzioni $q$ che valgano 1 in corrispondenza di $j$ e che siano nulle in tutti gli altri punti.\\
Si sceglie quindi come base dello spazio vettoriale di polinomi l'insieme:
\[
\{l_0(x),...,l_k(x)\}
\]
dove 
\[
l_j(x) = 
\frac{(x-x_0)\cdot \s ... \s \cdot (x-x_{j-1})\cdot (x-x_{j+1})\cdot \s ... \s \cdot(x-x_k)}{(x_j-x_0)\cdot \s ... \s \cdot (x_j-x_{j-1})\cdot (x_j-x_{j+1})\cdot \s ... \s \cdot(x_j-x_k)} = \frac{\prod^k_{i\neq j} (x-x_i)}{\prod^k_{i\neq j} (x_j-x_i)}
\]

\paragraph{Proprietà}
\[
l_j(x_i) =
\begin{cases}
1 & \text{se } i = j, \\
0 & \text{se } i \neq j
\end{cases}
\quad \Leftarrow \quad
l_j(x_j)
= \frac{\prod_{i \neq j} (x_j - x_i)}{\prod_{i \neq j} (x_j - x_i)} = 1,
\qquad
l_j(x_h)
= \frac{\prod_{i \neq j} (x_h - x_h)}{\prod_{i \neq j} (x_j - x_i)} = 0
\]
\begin{nota}
    Gli $l_j(x)$ sono quindi polinomi che passano da 1 nel punto $x_j$ e da 0 in tutti gli altri punti (non ci interessa il comportamento di questi polinomi in punti intermedi).
\end{nota}
\noindent
In particolare, se scriviamo $P_k(x)$ come
\[
P_{(k)} = \sum^k c_j l_j(x)
\]
si ha che i coefficienti $c_j$ si trovano molto facilmente.
Infatti,
\[
\begin{aligned}
y_h = P_k(x_h)
= \sum_{j=0}^{k} c_j \, l_j(x_h)
= c_0 l_0(x_h) + c_1 l_1(x_h) + \cdots + c_h l_h(x_h) + \cdots + c_k l_k(x_h) =\\
= 0 + 0 + \cdots + c_h l_h(x_h) + \cdots + 0
= c_h \cdot 1 = c_h
\end{aligned}
\]
Quindi
$
c_h = y_h \; 
\forall h = 0, \ldots, k
$ per cui il polinomio interpolante nella base di Lagrange si scrive come:
\begin{definizione}
    [Polinomio interpolante di Lagrange]
    \label{def: polinomio interpolante di Lagrange}
    \[
    P_k(x) = \sum_{j=0}^{k} y_j \, l_j(x)
    \]
    Nel polinomio ogni valore di $f(x_j)$ pesa solo nel suo nodo.
\end{definizione}
\begin{nota}
    Quando andiamo a costruire/calcolare il polinomio di Lagrange lo si vorrà fare in certi punti, infatti $l_j$ è definito come funzione di x, ovvero i punti in cui si vuole valutare il polinomio di interpolazione.
    Ricordare sempre che questi sono i punti di valutazione, non quelli di rappresentazione.
\end{nota}

\subsection{Interpolazione nella base di Newton}
Se aggiungiamo un dato in più, ovvero se conosciamo un punto in più, allora l'interpolazione di Lagrange appena calcolata perde di validità e dobbiamo ricalcolare tutta l'interpolazione, quindi anche tutti i polinomi, da 0.
In poche parole \textbf{l'interpolazione di Lagrange non è incrementale}.\\
Nell'interpolazione di Newton si sceglie la base $n_j(x)$ con $j=0,...,k$ definita come:
\[
n_j(x) = \left\{ 1, (x - x_0), (x - x_0)(x - x_1),\dots ,(x - x_0)\cdot \s ... \s \cdot (x - x_k) \right\}
\]
che corrisponde a
\[
n_j(x) = \left\{ n_0(x), n_1(x), n_2(x), ..., n_k(x) \right\}
\]
In particolare $n_j(x)$ è un polinomio di grado $j$ che dipende solo da $x_0,..,x_{j-1}$. 
Quindi se si aggiunge un nodo di interpolazione si deve solamente calcolare un nuovo polinomio di base.\\
L'obbiettivo finale è ottenere come al solito un polinomio del tipo:
\[
P_k(x) = \sum_{j=0}^k c_j n_j(x)
\]
Per Newton i coefficienti $c_j$ si esprimono con la formula
\[
c_j = f\left[x_0,..,x_j\right]
\]
dove $f\left[...\right]$ si definisce per ricorrenza come la differenza divisa di ordine $k$ (vedi sotto).


\begin{lezione}
    [Lezione 14 19/04]
\end{lezione}

\subsection{Differenze Divise}
Abbiamo appena visto come, per costruire il polinomio di Newton, ci serva calcolare i coefficienti moltiplicativi dei polinomi che formano la base tramite le differenze divise.\\
Concettualmente, una differenza divisa è la naturale generalizzazione del rapporto incrementale
$
\Bigl(
f[x_0,x] = \frac{f(x)-f(x_0)}{x-x_0}
\Bigr)
$
a più punti: così come il rapporto incrementale coglie la variazione media di $f$ tra due soli valori, la differenza divisa di ordine $k$ estende questa idea a $k+1$ punti, usando esclusivamente le valutazioni di $f$ nei nodi (non serve conoscere $f$ in forma esplicita, né calcolarne le derivate).\\
Quello che fa non è altro che permettere di aggiornare il valore del rapporto incrementale (=frazione che misura la velocità media di variazione della funzione stessa) precedente in funzione dell'aggiunta di un punto di valutazione (dalla definizione si vede come si sottrae alla vecchia valutazione la differenza con il nuovo punto, permettendo quindi di fare il rapporto con la variazione sulle ordinate a partire dal nuovo punto stesso).

Per calcolarle nella pratica quello che si fa è organizzare i dati in una struttura tabellare (il \textbf{quadro delle differenze divise}), costruita per ricorsione a partire dai soli valori $f(x_i)$, secondo la definizione seguente:
\begin{definizione}
    [Definizione Differenza Divisa di ordine k]
    Dati $k$ punti distinti in un intervallo si definisce la differenza divisa di ordine $k$ la funzione:
    \[
    \begin{array}{l}
         f[x]=f(x)\\
        f_{[x_0,x]} = \frac{f(x) - f(x_0)}{x-x_0}\\
        f_{[x_0,...,x_{k-1},x]} = \frac{f_{[x_0,...,x_{k-2},x]} - f_{[x_0,...,x_{k-1}]}}{x-x_{k-1}}
    \end{array}
    \]
\end{definizione}
\noindent
Matrice delle differenze divise:
\[
A =
\begin{bmatrix}
f(x_0) & 0 & 0 & 0 & \cdots & 0 \\
f(x_1) & f[x_0, x_1] & 0 & 0 & \cdots & 0 \\
f(x_2) & f[x_0, x_2] & f[x_0, x_1, x_2] & 0 & \cdots & 0 \\
\vdots & \vdots & \vdots & \ddots & \ddots & \vdots \\
f(x_k) & f[x_0, x_k] & f[x_1, x_k] & \cdots & f[x_0, \dots, x_k]
\end{bmatrix}
\]
\begin{nota}
    Gli elementi sulla diagonale saranno i coefficienti che andranno a formare il polinomio.
\end{nota}
\noindent
Per calcolare un elemento in posizione $(i,j)$ nella matrice A si può usare la formula:
\[
\boxed{a_{ij}=\frac{a_{i(j-1)}-a_{(j-1)(j-1)}}{x_i - x_{j-1}}}
\]
dove le $x_i$ sono gli indici delle righe ($i$ per la riga corrente, $j$ per la colonna corrente) dati dal valore di $x$ in quella riga; in particolare si calcolano le entrate di $A$ procedendo per righe dall'alto verso il basso e per colonne da sinistra verso destra.

\paragraph{Osservazione}
Nel caso in cui si trovino tutti elementi uguali in una colonna del quadro delle differenze divise significa che le differenze divise di ordine superiore sono tutte nulle.\\
In particolare $P_k$ \textbf{ha grado pari all'ordine dell'ultima differenza divisa} $\ne 0$.
\paragraph{Esempi differenze divise}

\paragraph{Osservazione}
\uline{L'ordine dei nodi/punti di interpolazione non è essenziale per determinare $P_k(x)$.}
Questo può risultare utile in alcune circostanze per semplificare i calcoli: ad esempio quando realizziamo il quadro delle differenze divise non dobbiamo preoccuparci dell'ordine delle x degli elementi che andiamo a scrivere nella prima colonna.

\begin{teorema}
    [Teorema Errore di interpolazione polinomiale]
    \label{th: errore interpolazione polinomiale}
    Sia $f \in C^{k+1}[a,b]$ e $P_k$ il polinomio interpolante su $k+1$ nodi distinti. Allora $\forall x \in [a,b]$ esiste $\epsilon_x$ tra $x$ e i nodi tale che:
    \[
    f(x) - P_k(x) = \frac{f^{(k+1)}(\epsilon_x)}{(k+1)!}\cdot \pi(x)
    \]
    con $\pi(x) = \prod^k_j (x- x_j)$.
\end{teorema}

\subsection{Recap basi di Newton/Lagrange}

\begin{center}
\begin{tabular}{|p{0.46\textwidth}|p{0.46\textwidth}|}
\hline
\textbf{Newton} & \textbf{Lagrange} \\ \hline

$m_j(x)=(x-x_0)\cdots(x-x_{j-1})$ &
$\ell_j(x)=\displaystyle\prod_{i\neq j}\frac{x-x_i}{x_j-x_i}$ \\

$f[x_0,\ldots,x_j]$ &
$f(x_j)=y_j$ (costa tanto trovarli) \\ \hline

Il grado di $P_k(x)$ si vede facilmente. &
Per ottenere il grado di $P_k(x)$
devo espandere (cambiare) in un'altra base (es.\ monomi). \\ \hline

Da $P_k$ a $P_{k+1}$:
calcolo solo $m_{k+1}(x)$ e
$f[x_0,\ldots,x_{k+1}]$. &
Da $P_k$ a $P_{k+1}$:
devo ricalcolare tutti gli $\ell_j(x)$. \\ \hline
\end{tabular}
\end{center}
Formula dell'errore (indipendente dalla base):
\[
f(x)-P_k(x)
=
\Pi_k(x)\,
\frac{f^{(k+1)}(\xi_x)}{(k+1)!}
\]


\subsection{Interpolazione di Hermite}
\begin{nota}
    Questa sezione non fa parte del programma 2025/26
\end{nota}
\noindent
Supponiamo di avere $x_0,...,x_k\in\mathbb{R}$ e di conoscere $f(x_0),...,f(x_k)$ e $f'(x_0),...,f'(x_k)$, vogliamo trovare un polinomio $P_{2k+1}(x)$ di grado al più $2k+1$\footnote{si hanno infatti 2 condizioni per ogni nodo: una su f e una sulla derivata.} tale che:
\[
P_{2k+1}(x_j) = f(x_j)\quad ,\quad P'_{2k+1}(x_j) = f'(x_j)
\]
\uline{ovvero un polinomio che passa per i punti dati (come per Lagrange e Newton) con anche la stessa pendenza con cui ci arriva $f$.}\\
Le condizioni appena sopra citate viste come condizioni sui coefficienti del polinomio sono $2k+2$ equazioni lineari. 
Ci aspettiamo quindi che servono $2k+1$ gradi di libertà in $P_{2k+1}(x)$ per risolvere il problema. 
Per questo motivo $P_{2k+1}(x)$ viene preso di grado $\leq 2k+1$ (questo è infatti definito da $2k+2$ coefficienti).
\begin{teorema}
    [Polinomio di Hermite]
    se $x_i\neq x_j \quad \forall i\neq j$ allora \textbf{esiste ed è unico} $H_{2k+1}(x)$ che verifica
    \[
    P_{2k+1}(x_j) = f(x_j)\quad ,\quad P'_{2k+1}(x_j) = f'(x_j)
    \]
    Questo viene detto \textbf{Polinomio di Hermite}.
\end{teorema}
\subsubsection{Come si costruisce il polinomio di Hermite}
Si adotta un approccio simile all'interpolazione di Lagrange: si cercano di famiglie di polinomi $h_{0j}(x)$ e $h_{1j}(x)$ con $j=0,...,k$ che abbiano le seguenti proprietà:

\[
\begin{array}{cc}
     h_{0j}(x_i) = 
\begin{cases} 
1 & i = j \\ 
0 & i \neq j 
\end{cases} &  
h'_{0j}(x_i) = 0 \quad \forall i\\
   h_{1j}(x_i) = 0  \quad \forall i&  h'_{1j}(x_i) = 
\begin{cases}
1 & i = j \\ 
0 & i \neq j 
\end{cases}
\end{array}
\]
Se conosciamo $h_{0j}(x)$ e $h_{1j}(x)$ con queste proprietà allora scriverci $H_{2k+1}(x)$ è facile:
\[
H_{2k+1} = \sum_{j=0}^{k} f(x_j) \cdot h_{0j}(x) + \sum_{i=0}^{k} f'(x_j) \cdot h_{1i}(x)
=\sum_{j=0}^{k} f(x_j) \cdot \underbrace{h_{0j}(x_i)}_{*} + \sum_{i=0}^{k} f'(x_j) \cdot \underbrace{h_{1i}(x)}_{=0}=
\]
\begin{nota}
    *: si annullano tutti tranne il j-esimo
\end{nota}
\[
=\sum_{j=0}^{k} f(x_j) \cdot \underbrace{h_{0j}(x_i)}_{=1} = f(x_j)
\]
$H'_{2k+1} = \dots = f'(x_j)$\\
\textbf{Come si trovano $h_{0j}(x)$, $h_{1j}(x)$}\\
Idea: partire dai polinomi della base di $y$ e li cerchiamo nella funzione.
\begin{align*}
h_{0j}(x) &= (A_j x + B_j)l_j^2(x) \\
h_{1j}(x) &= (C_j x + D_j)l_j^2(x)
\end{align*}
Si determinano $A_j$, $B_j$, $C_j$, $D_j$ in modo che le condizioni $(*)$ siano soddisfatte.
 
\begin{lezione}
    [Lezione 15 22/04]
\end{lezione}
\section{Fenomeno di Runge}
Fenomeno che mostra come nell'interpolazione, aumentando il numero di nodi, invece di migliorare, l'approssimazione peggiora vicino agli estremi dell'intervallo, con forti oscillazioni.\\
Questo nasce dal fatto che, quando si considera l'interpolante polinomiale $P_k(x)$ di una certa $f$, l'unico strumento che abbiamo per diminuire l'errore di approssimazione:
\[\underset{x \in [a,b]}{max}\, |P_k(x) - f(x)|\]
è il grado di $P_k$, cioè ci aspettiamo che aumentando $k$, quindi aumentando i punti di interpolazione, l'errore vada a $0$:
\[
\lim_{k \to \infty} |P_k(x) - f(x)| = 0
\]
Tuttavia, questa aspettativa può rivelarsi falsa, ed è proprio questo il cuore del fenomeno di Runge. Carl Runge osservò che, anche per una funzione analitica e ben comportata (come $f(x) = \frac{1}{1 + x^2}$), se si effettua l'interpolazione polinomiale utilizzando nodi equidistanti su un intervallo ampio, si ottengono oscillazioni sempre più marcate del polinomio interpolante, soprattutto verso i bordi dell'intervallo (si veda grafico di esempio in \autoref{fig:runge-esempio}). Paradossalmente, aggiungere più punti non migliora l'approssimazione, ma anzi può peggiorarla nelle zone estreme.\\
\uline{Questo fenomeno si manifesta} in particolare \uline{quando}:
\begin{itemize}
    \item si usano \textbf{punti equidistanti su un intervallo lungo};
    \item si costruisce un \textbf{polinomio unico di grado elevato}.
\end{itemize}
Per evitare o ridurre il fenomeno di Runge, si possono adottare diverse strategie pratiche:
\begin{itemize}
    \item Utilizzare nodi di Chebyshev, che sono più densi vicino ai bordi e aiutano a contenere le oscillazioni (vedi lezione 7 di laboratorio a \autoref{sec: lab nodi chebyshev}).
    \item Utilizzare splines (es. spline cubiche), cioè più polinomi di basso grado incollati tra loro in modo continuo, ovvero limitare il grado del polinomio e suddividere l'intervallo in sottosezioni.
\end{itemize}
\paragraph{Curiosità}
C'è anche un limite teorico più profondo. Esiste infatti un teorema che afferma: per qualunque strategia di scelta dei nodi di interpolazione su $[a, b]$, esiste sempre una funzione continua $f \in C([a, b])$ tale che: $\lim_{k \to \infty} |P_k(x) - f(x)| \neq 0$.
\paragraph{Conclusione}
Questo significa che non esiste un insieme "perfetto" di nodi che garantisca convergenza dell'interpolante per tutte le funzioni continue. In altre parole, non esiste una scelta universale di punti che vada bene per ogni funzione e ogni intervallo. Tuttavia, alcune scelte sono molto migliori di altre in pratica. I nodi di Chebyshev, ad esempio, sono una strategia particolarmente efficace per ridurre l'errore massimo e controllare le oscillazioni.\\
In sintesi, il fenomeno di Runge e il relativo limite teorico mostrano che l'interpolazione polinomiale va usata con cautela: non sempre più punti migliorano l'approssimazione, e non esiste una strategia che funzioni sempre per ogni funzione. È fondamentale scegliere con intelligenza i nodi e, in molti casi, preferire approcci alternativi come i metodi a tratti (splines) o l'interpolazione con nodi ottimizzati.

\section{Interpolazione polinomiale a tratti}
Idea: si divide l'insieme dei nodi di interpolazione in sottoinsiemi di $s + 1$ nodi consecutivi
\[
x_0 < x_1 < \ldots < x_k: \quad \underbrace{x_0, \ldots, x_s}_{\text{gruppo 1}} = \underbrace{x_{s+1}, \ldots, x_k}_{\text{gruppo 2}}
\]
\begin{itemize}
    \item su ogni sottoinsieme costruisco il polinomio interpolante di grado $s$ $(q_1(x), \ldots, q_{k/s}(x))$
    \item Definisco la funzione interpolante $F(x)$ come: $F(x) = q_j(x)$ se $x$ appartiene al gruppo $j$ ovvero è in $[x_{s(j-1)}, x_{s(j)}]$
\end{itemize}

\subsection{Funzioni Spline}

Dati $x_0,<...<x_k\in [a,b]$ si prendono i sottointervalli:
\[
[x_0,x_1],\s...\s, [x_{k-1}, x_k]
\]
ed in ciascuno di essi si prende un polinomio interpolante di grado 3 in modo tale che la funzione definita a tratti sia derivabile 2 volte in $[a,b]$.
Pi\`u formalmente:

\begin{definizione}
    [Spline cubica]
    Si dice tale \textbf{spline cubica} per $f$ rispetto ai punti di interpolazione $(x_0, f(x_0)), \dots , (x_k, f(x_k))$ la funzione $S_3(x) : [a,b] \rightarrow \mathbb{R}$ tale che:
    \begin{enumerate}
        \item $S_3(x)|_{[x_{i-1},x_i]}$ è un polinomio di grado 3 $\forall \;i = 1, \dots, k+1$
        \item $S_3(x_i) = f(x_i) \quad \forall i = 1,...,k$
        \item $S_3(x) \in \mathcal{C}^2([a,b])$
    \end{enumerate}
    In particolare il punto 3 ci indica che:
    \begin{itemize}
        \item $p'_i (x_i) = p'_{i+1}(x_i)$
        \item $p''_i (x_i) = p''_{i+1}(x_i)$
    \end{itemize}
    \begin{nota}
        Il punto 3, unito al punto 2, ci indicano che la funzione spline deve, nei nodi, corrispondere perfettamente con la funzione di riferimento (stesso valore, pendenza, curvatura)
    \end{nota}
\end{definizione}
\begin{teorema}
    [Esistenza e unicità della spline]
    Siano $x_0,\ldots,x_k$ nodi equispaziati in $[a,b]$, in particolare
    \[
    x_i-x_{i-1}=\frac{b-a}{k}=h,
    \qquad
    f:[a,b]\to\mathbb{R}.
    \]
    Allora la spline naturale $S_3(x)$ per i nodi $x_j$ e la funzione $f$ \textbf{esiste ed \`e unica}.
    Inoltre i polinomi $p_i(x)$, $i=1,\ldots,k$, che la costituiscono sono della forma $(y_i=f(x_i))$:
    \[
    \begin{aligned}
    p_i(x)
    &=
    \Bigg[
    y_{i-1}
    +\left(m_{i-1}+\frac{2y_{i-1}}{h}\right)(x-x_{i-1})
    \Bigg]
    \left(\frac{x-x_i}{h}\right)^2 +\\
    &\quad +
    \Bigg[
    y_i
    +\left(m_i-\frac{2y_i}{h}\right)(x-x_i)
    \Bigg]
    \left(\frac{x-x_{i-1}}{h}\right)^2
    \end{aligned}
    \]
    Dove $m_0,\ldots,m_k$ sono la soluzione del sistema lineare associato:
    \[
    \begin{bmatrix}
    2 & 1 &   &        & 0 \\[-10pt]
    1 & 4 & 1 &        &   \\[-10pt]
      & 1 & \ddots & \ddots &   \\[-10pt]
      &   & \ddots & 4 & 1 \\[-10pt]
    0 &   &        & 1 & 2
    \end{bmatrix}
    \begin{bmatrix}
    m_0 \\[-10pt] m_1 \\[-10pt] \vdots \\[-10pt] m_{k-1} \\[-10pt] m_k
    \end{bmatrix}
    =
    \frac{3}{h}
    \begin{bmatrix}
    y_1-y_0 \\[-10pt]
    y_2-y_0 \\[-10pt]
    \vdots \\[-10pt]
    y_{k-1}-y_{k-2} \\[-10pt]
    y_k-y_{k-1}
    \end{bmatrix}
    \]
\end{teorema}
\begin{dimostrazione}
    \textbf{Idea:} in ogni intervallo $[x_{i-1}, x_i]$ si prende $p_i(x)$ come il polinomio di Hermite rispetto ai nodi d'interpolazione
    \[
        (x_{i-1}, y_{i-1}), \quad (x_i, y_i), \quad (x_{i-1}, m_{i-1}), \quad (x_i, m_i)
    \]
    dove i valori $m_i$ sono, per il momento, da determinare.
    \\
    Fatta questa scelta su ogni intervallo, si ottiene una funzione definita a tratti continua e $\mathcal{C}^1([a,b])$: infatti, per costruzione, ogni $p_i$ interpola sia i valori $y_{i-1}, y_i$ sia le pendenze $m_{i-1}, m_i$ agli estremi dell'intervallo, dunque
    \[
        p_i(x_i) = p_{i+1}(x_i) = y_i,
        \qquad
        p'_i(x_i) = p'_{i+1}(x_i) = m_i.
    \]
    Imponendo infine la condizione
    \[
        p''_i(x_i) = p''_{i+1}(x_i), \qquad i = 1, \dots, k-1,
    \]
    si determinano i coefficienti $m_i$ come soluzione del sistema lineare riportato nell'enunciato.
\end{dimostrazione}

\subsubsection{Precisione delle spline}
L'andamento dell'errore in questo tipo di funzioni è descritto nella sezione \autoref{sec: lab spline cubiche naturali}.\\
In particolare si fa notare come l'errore delle spline cubiche decresce, ordine di potenza $4$, rispetto alla funzione di ampiezza dell'intervallo $h$, contro un minore ordine di potenza $2$ che contraddistingue la più semplice interpolazione a tratti lineare.

\section{Approssimazione di funzioni ai minimi quadrati}
Si affronta ora il problema, che prima avevamo visto dal punto di vista matriciale e algebrico, dal punto di vista funzionale: si cerca infatti una funzione che minimizzi gli errori in corrispondenza dei nodi.
\\
Stiamo dunque rilassando la condizione del "\textit{passare esattamente dai nodi}" introdotta con l'interpolazione.
\begin{figure}[H]
    \centering
    \includegraphics[width=0.5\linewidth]{img/appr f min q.png}
\end{figure}
\noindent
Si parte da:
$f:\mathbb{R}\to\mathbb{R}$, $(x_0,y_0),\s...\s, (x_k, y_k)$ e $y_i= f(x_i)$ con $k$ molto grande.\\
Cerchiamo un approccio che generalizza l'interpolazione nelle direzioni:
\begin{enumerate}
    \item non passare necessariamente dei punti ma "sufficientemente vicini"
    \item non utilizzare necessariamente polinomi
\end{enumerate}
Quella che \textbf{cerchiamo} è quindi una $\Phi(x) \approx f(x)$ con $\boxed{{\Phi(x) = \sum_{j=0}^mc_j\cdot G_j(x)}}$ con $m\leq k$.
\begin{nota}
    $G_j$ è una funzione modello: funzione elementare, non necessariamente un polinomio.
\end{nota}
\subsection{Come scegliere i coefficienti Cj}
Li cerchiamo in modo tale che l'errore rispetto a $f$ nei punti $x_i$ sia minimo/piccolo.\\
Definiamo una funzione che, dati $c_0, \ldots, c_n$, misuri l'errore:
\begin{itemize}
    \item
    $
    \psi(c_0, \ldots, c_n)
    = \sum_{j=0}^{k} \left[ \Phi(x_j) - f(x_j) \right] \Rightarrow
    $
    potrebbe essere piccola anche per approssimazioni non buone,
    quindi non va bene;

    \item
    $
    \psi(c_0, \ldots, c_n)
    = \sum_{j=0}^{k} \lvert \Phi(x_j) - f(x_j) \rvert \Rightarrow
    $
    andrebbe bene ma $\psi$ non \`e pi\`u differenziabile,
    quindi \`e difficile da minimizzare, per cui non va bene.
\end{itemize}
La \textbf{scelta definitiva} \`e:
\[
\psi(c)
:= \sum_{j=0}^{k} \bigl( \Phi(x_j) - f(x_j) \bigr)^2
= \sum_{j=0}^{k}
\left[
\sum_{h=0}^{m} c_h \, G_h(x_j) - f(x_j)
\right]^2,
\]
ottenuta sommando gli scarti quadratici nei nodi.\\
\uline{Le variabili da determinare sono} $\boxed{c_0, \dots, c_m}$, e per farlo si risolve:
\[
\boxed{
\min_{c_0, \dots, c_m} \psi(c_0, ..., c_m)
= \sum_{j=0}^{k} \bigl[ \Phi(x_j) - f(x_j) \bigr]^2
}
= \sum_{j=0}^{k}
\left[
\sum_{h=0}^{m} c_h \, G_h(x_j) - f(x_j)
\right]^2
\]
e si prende il punto minimo.\\
Se definiamo il vettore
\[
r(c) :=
\begin{bmatrix}
\Phi(x_0) - f(x_0) \\[-10pt]
\Phi(x_1) - f(x_1) \\[-10pt]
\vdots \\[-10pt]
\Phi(x_k) - f(x_k)
\end{bmatrix}
\in \mathbb{R}^{k+1},
\quad \text{e} \quad
\psi(c)
= \sum_{j=0}^{k} \bigl(r_j(c)\bigr)^2
= \|r(c)\|_2^2,
\]
segue che
\[
r_j(c)
= \sum_{h=0}^{m} c_h \, G_h(x_j) - f(x_j)
= c_0 \, G_0(x_j) + c_1 \, G_1(x_j) + \cdots + c_m \, G_m(x_j) - f(x_j)
\]
In forma matriciale:
\[
r_j(c)
=
\begin{bmatrix}
\, G_0(x_j) & \, G_1(x_j) & \cdots & \, G_m(x_j)
\end{bmatrix}
\cdot
\begin{bmatrix}
c_0 \\[-10pt] \vdots \\[-10pt] c_m
\end{bmatrix}
- f(x_j)
\]
per cui possiamo ottenere la \uline{formulazione matriciale del problema}:
\[
A =
\begin{bmatrix}
G_0(x_0) & \cdots & G_m(x_0) \\
\vdots  &        & \vdots  \\
G_0(x_k) & \cdots & G_m(x_k)
\end{bmatrix}
\in \mathbb{R}^{(k+1)\times(m+1)},
\qquad
b =
\begin{bmatrix}
f(x_0) \\
\vdots \\
f(x_k)
\end{bmatrix}
\in \mathbb{R}^{k+1}.
\]

\[
r(c) = Ac - b
\;\Longrightarrow\;
\min_{c}\psi(c)
= \min_{c}\|r(c)\|_2^2
= \min_{c}\|Ac - b\|_2^2.
\]
Il problema é quindi equivalente a risolvere il sistema sovradeterminato $Ac=b$ nel senso dei minimi quadrati.\\
Quindi possiamo utilizzare il \textbf{sistema delle equazioni normali} (\autoref{def: eq normali}) oppure il \textbf{metodo QR} per trovare $c$.

\subsubsection{Esempio}
Calcolare l'equazione della \uline{retta} che approssima nel senso dei minimi quadrati i dati:
\begin{center}
\begin{tabular}{c|ccccc}
    $x$ & 1 & 2 & 3 & 4 & 5\\
    \hline
    $y$ & 1.4 & 3.1 & 4.8 & 6.8 & 8.5
\end{tabular}
\end{center}

\noindent
Quindi possiamo utilizzare le equazioni normali o il metodo QR per trovare $c$.

\subparagraph{Retta}
\[
\Phi_1(x) = c_1x+c_0, \qquad \varphi_0(x)=1,\quad \varphi_1(x)=x
\]
\[
\renewcommand{\arraystretch}{1}
A = \begin{bmatrix} 1 & 1\\ 1 & 2\\ 1 & 3\\ 1 & 4\\ 1 & 5\end{bmatrix}, \qquad
b = \begin{bmatrix} 1.4\\ 3.1\\ 4.8\\ 6.8\\ 8.5\end{bmatrix}
\]
Le equazioni normali ci portano a considerare il sistema $A^TAc=A^Tb$:
\[
A^TA = \begin{bmatrix} 5 & 15\\ 15 & 55\end{bmatrix}, \qquad
A^Tb = \begin{bmatrix} 24.6 \\ 91.7 \end{bmatrix}
\]
Risolvendo il sistema $2\times2$ si ottiene:
\[
c = \begin{bmatrix} -0.45\\ 1.79\end{bmatrix}
\]
\[
\Phi_1(x) = 1.79\cdot x - 0.45
\]

\subparagraph{Caso con Parabola}
\[
\Phi_2(x) = c_0+c_1x+c_2x^2, \qquad \varphi_0(x)=1,\quad \varphi_1(x)=x,\quad \varphi_2(x)=x^2
\]
\[
\renewcommand{\arraystretch}{1}
A = \begin{bmatrix} 1 & 1 & 1\\ 1 & 2 & 4\\ 1 & 3 & 9\\ 1 & 4 & 16\\ 1 & 5 & 25\end{bmatrix}
\]
\[
A^TA = \begin{bmatrix} 5 & 15 & 55\\ 15 & 55 & 225\\ 55 & 225 & 979\end{bmatrix}, \quad
b = \begin{bmatrix} \vdots\end{bmatrix} \; \Rightarrow \;
A^Tb = \begin{bmatrix} 24.6\\ 91.7\\ 378.3\end{bmatrix}
\Rightarrow
c=
\begin{bmatrix}
    -0.3\\
    1.66\\
    0.021
\end{bmatrix}
\]

\chapter{Metodi per l'integrazione}

\begin{lezione}
    [Lezione 16 26/04]
\end{lezione}
\noindent
\section{Integrazione Numerica}
Si vogliono fornire metodi per approssimare quantità del tipo
\[
I(\rho \cdot f) = \int^b_a f(x) \cdot \rho(x) dx
\]
con $\rho(x)$ funzione peso definita come:
$\rho(x) \geq 0$ su $[a, b]$ \textbf{funzione peso} tale che 
\[
m_k = \int_a^b |x|^k \rho(x)\,dx < +\infty \quad \forall k = 0,1,2,...
\]
dove le quantità $m_k$ sono dette \textbf{momenti}.
\\
I motivi per cui può servire una rappresentazione di questo tipo sono:
\begin{itemize}
    \item Di molte funzioni non esiste una espressione esplicita della primitiva.
    \item Anche se la primitiva esiste, è o costosa da calcolare e/o da valutare.
    \item A volte conosciamo $f(x)$ solo in qualche punto o è fornita come black box, quindi si può valutare ma non si sa come è fatta $f$.
    \item Stiamo parlando del caso $f : \mathbb{R} \to \mathbb{R}$.
\end{itemize}
\section{Formula di Quadratura}
L'idea di base è quella di approssimare un integrale tramite una somma pesata di valori della funzione in punti opportuni, riprendendo quindi i presupposti per la definizione di integrale di Rienmann, cioè la sommatoria di larghezza dell'intervallo per valore della funzione nell'intervallo.\\
Di seguito un'anteprima di come la trattazione dell'argomento si differenzi in base al grado di libertà offerto:
\begin{itemize}
    \item Nodi fissati e pesi liberi
    \item Nodi e pesi liberi
\end{itemize}
\begin{nota}
    Non trattiamo il caso di nodi e pesi fissati in quanto induce a formule non ottimizzate con gradi di precisione bassi.
\end{nota}
% \begin{table}[H]
%     \centering
%     \renewcommand{\arraystretch}{1.4}
%     \resizebox{\textwidth}{!}{%
%     \begin{tabular}{|l|l|l|}
%     \hline
%     & \textbf{Newton-Cotes} & \textbf{Gaussiana} \\
%     \hline
%     Nodi
%     & equidistanti, \textbf{fissati}
%     & zeri di polinomi ortogonali, \textbf{ottimizzati} \\
%     \hline
%     Pesi
%     & calcolati da sistema lineare
%     & calcolati da sistema non lineare \\
%     \hline
%     Grado di precisione
%     & $\geq n$ (o $n+1$ se $n$ pari)
%     & $2n+1$ \\
%     \hline
%     Limite
%     & instabilità per $n>7$ (pesi negativi)
%     & nodi non intuitivi, ma sempre stabile \\
%     \hline
%     Esempi
%     & Trapezi ($n=1$), Simpson ($n=2$)
%     & Gauss-Legendre, Gauss-Chebyshev \\
%     \hline
%     \end{tabular}}
%     \caption{Confronto tra formule di Newton-Cotes e formule gaussiane con $n$ nodi}
% \end{table}
\begin{definizione}
    [Formula di Quadratura]
    Formula del tipo
    \[
    J_n(f) = \sum_{j=0}^n a_j \cdot f(x_j)
    \]
    dove $a_j$ sono detti pesi e $x_j$ sono i nodi.
    In particolare abbiamo \uline{$n+1$ pesi e $n+1$ nodi}.\\
    Le formule di quadratura sono dunque formate dalla sommatoria di coefficienti, i pesi, e dei valori della funzione di cui si vuole approssimare l'integrale in punti detti nodi.
\end{definizione}
\begin{definizione}
    [Errore della formula di quadratura]
    \[
    E_n(f) = I(\rho \cdot f) - J_n(f)
    \]
\end{definizione}
\begin{definizione}
    [Grado di precisione di una formula di quadratura]
    Una formula di quadratura ha grado di precisione $m$ se verifica:
    \[
    E_n(1) = E_n(x) = \s ...\s =E_n(x^m) = 0 \quad \land \quad E_n(x^{m+1}) \neq 0 
    \]
    ovvero la formula di quadratura con un certo grado di precisione integra esattamente ogni monomio di grado $\leq m$.
\end{definizione}

\paragraph{Osservazione}
Per linearità se {$J_n(\cdot)$} ha grado di precisione $m$ significa che \uline{integra esattamente ogni polinomio $q(x)$ che ha grado $\le m$} ($E_n(q)=0$).\\
Questo ci è garantito dal fatto che il polinomio $q(x)$ è combinazione lineare di monomi su cui la formula di quadratura è esatta.

\subsection{Formula dei Trapezi}
\begin{definizione}
    [Formula dei Trapezi]
    \label{def: trapezi}
    \noindent
    \begin{minipage}[c]{0.25\linewidth}
        ~
    \end{minipage}
    \begin{minipage}[c]{0.4\linewidth}
    \[
        J_1(f) = \frac{b-a}{2} \cdot \left[ f(a) + f(b)\right]
    \]
    \end{minipage}
    \begin{minipage}[c]{0.3\linewidth}
        \raggedleft
        \includegraphics[width=\linewidth]{img/trapezi1.png}
    \end{minipage}
\end{definizione}
\begin{nota}
Ha grado di precisione = 1.
\end{nota}

\subsection{Formula di Simpson/Cavalieri-Simpson}
Si ottiene aggiungendo un punto di valutazione alla formula dei trapezi, in particolare si aggiunge il punto di mezzo tra $a$ e $b$.
\begin{definizione}
    [Formula di Simpson/Cavalieri-Simpson]
    \label{def: simpson}
    \[
    J_2(f) = \frac{b-a}{6}\cdot\left[ f(a) + 4f\left(\frac{a+b}{2} \right) +f(b)\right]
    \]
    \begin{nota}
        Ha grado di precisione = 3.
    \end{nota}
\end{definizione}
\begin{figure}[H]
    \centering
    \includegraphics[width=0.3\linewidth]{img/cav simp.png}
    \caption{Qui l'integrale viene approssimato tramite quello della parabola che interpola la funzione nei 3 punti.}
\end{figure}

\subsection{Massimizzare il grado di precisione di una formula di quadratura}

Se i nodi sono dati ($x_0,\dots, x_n$) allora massimizzare il grado di precisione equivale a risolvere il sistema lineare di $n+1$ equazioni:
\begin{nota}
    per semplicità in questo esempio consideriamo $\rho(x) = 1$
\end{nota}
\[
\begin{cases}
a_0 + a_1 + \cdots + a_n
= \displaystyle\int_a^b 1 \, dx
= m_0 \\[6pt]

a_0 x_0 + a_1 x_1 + \cdots + a_n x_n
= \displaystyle\int_a^b x \, dx
= m_1 \\[6pt]

a_0 x_0^{2} + a_1 x_1^{2} + \cdots + a_n x_n^{2}
= \displaystyle\int_a^b x^{2} \, dx
= m_2 \\[6pt]

\vdots \\[4pt]

a_0 x_0^{n} + a_1 x_1^{n} + \cdots + a_n x_n^{n}
= \displaystyle\int_a^b x^{n} \, dx
= m_n
\end{cases}
\]
che non è altro che il sistema $V^T\tilde{a} = b$ dove $V$ è la matrice di Vandermonde, $\tilde{a}=[a_0, \dots, a_n]$ e $b = [m_0, \dots , m_n]$.\\
Anche qui possiamo quindi fare le stesso considerazioni sul calcolo del determinante:
\[
det(V) = det(V^T) = \prod (x_i - x_j)
\]
per cui possiamo concludere che se $x_i \ne x_j$ il sistema ha un'unica soluzione $\Rightarrow$ 
\begin{teorema}
    [Unicità della formula di quadratura di grado di precisione massimo]
    \label{th: unicita formula quadratura grado massimo}
    Fissati i nodi ($n+1$) \textbf{esiste ed è unica} la formula di quadratura che ha grado di precisione almeno $n$.
\end{teorema}
\begin{nota}
    Se consideriamo le formule dei Trapezi e di C-S si nota che stiamo integrando al posto di $f$ il polinomio (di grado 1 per trapezi e 0 per Simpson) che interpola $f$ nei nodi.
\end{nota}
Se si ha la possibilità di scegliere sia i nodi che i pesi ($n$ dato) per massimizzare il grado di precisione di precisione bisogna risolvere un sistema non lineare:

\section{Formule di quadratura Gaussiana}
L'idea di queste è di scegliere in modo ottimale anche i nodi, oltre che hai pesi (a differenza di quanto si farà nelle \autoref{sec:formule di newton-cotes}), in questo modo si evitano i limiti di grado di precisione e i fenomeni instabili come quello di Runge, quest'ultimo infatti si verifica proprio in presenza di nodi equispaziati.
\begin{nota}
    Più avanti vedremo come queste sono un sottoinsieme delle formule interpolatorie.
\end{nota}
Dalla situazione descritta nel sistema precedente osserviamo che abbiamo $2n+2$ quantità/parametri da trovare ($a_0,...,a_n,x_0,...,x_n$) quindi si impongono $2n+2$ equazioni:
\[
\begin{array}{cc}
    E_n(x^k) =0 \quad k=0,...,2n+1 & m_j= \int_a^bx^jdx \\
    \begin{cases}
        a_0 + \s...\s +a_n = m_0 \\
        \vdots \\
        a_0x_0^{2n+1} +\s ...\s + a_nx_n^{2n+1} = m_{2n+1}
    \end{cases} & \leftarrow \begin{cases}
        \text{sistema non lineare di} \\
        \text{$2n+2$ eq. in $2n+2$ incognite}
    \end{cases}
\end{array}
\]
Risolvendo questo sistema otteniamo una formula di quadratura che ha grado almeno $2n+1$
\begin{teorema}
    [Esistenza e unicità della formula di quadratura Gaussiana]
    Il sistema non lineare sopra descritto ammette sempre un'unica soluzione per ogni intervallo $[a,b]$, peso $\rho(x)$ e numero di nodi $n+1$.
    In particolare c'è un'unica formula di quadratura con $n+1$ nodi che ha grado almeno $2n+1$, ovvero il massimo possibile con $n$ valutazioni, e tale formula viene detta \textbf{Formula di Quadratura Gaussiana}.
\end{teorema}
\section{Errore delle formule di quadratura}
Studiamo come si comporta l'errore per formule in generale non necessariamente Gaussiane.
\begin{teorema}
    [Teorema di Peano]
Data $f(x)$ continua e derivabile $m+1$ volte in $[a,b]$ con $J_n(f)$ formula con grado di precisione $m$ allora:
\[
\boxed{E_n(f) = \frac{1}{m!} \int_a^b f^{m+1}(t) \cdot G(t) \cdot \rho(t) dt}  = I(\rho \cdot f) - J_n(f)
\]
dove:
\[
\boxed{G(t) = E_n(s_m(x-t))} \quad \text{con} \quad s_m(x) = 
\begin{cases}
    x^m \quad x\geq 0 \\
    0 \quad x<0
\end{cases}
\]
è detto nucleo di Peano, ovvero la funzione che misura come la formula pesa la derivata che non sa integrare esattamente.
\end{teorema}
\begin{nota}
    Questa equazione per la definizione di errore è derivata direttamente dallo sviluppo di Taylor con resto integrale per $f\in\mathcal{C}^{m+1}$, in particolare è proprio l'addendo che definisce tale resto.
\end{nota}
In parole semplici il nucleo di Peano, ci fa vedere come e dove la specifica formula sta accumulando l'errore su quella parte di curva che non riesce a capire, ovvero è l'errore commesso dalla formula di quadratura quando si usa come funzione $s_m(x-t)$, cioè un polinomio traslato.\\
Per esempio, nel caso dei trapezi che ha grado di precisione 1, ignora il resto e misura quanto la derivata seconda della funzione e misura quanto sta mettendo sotto stress la funzione in ogni singolo punto (nodo).\\
Se $G(t)$ non cambia di segno per $t\in[a,b]$ allora l'espressione di $E_n(f)$ diventa:
\[
\frac{1}{m!} f^{m+1}(t) \, \cdot \int_a^b G(t) dt
\]
Quindi l'errore dipende dalla derivata di ordine $m+1$ della funzione f e dalla forma del nucleo G, il quale dipende solo dalla formula di quadratura (non da f).\\
Questo significa che se si valuta l'espressione precedente per 
\[f(x) = x^{m+1}\implies f^{(m+1)}(x) = (m + 1)!\]
si ottiene che
\[
E_n(x^{m+1}) = \frac{1}{m!} ((m + 1)!) \cdot \int_a^b G(t)\,dt = (m + 1) \int_a^b G(t)\,dt
\]
\[
\implies \int_a^b G(t)\,dt = \frac{E_n(x^{m+1})}{m + 1}
\]
\paragraph{Corollario}
Per cui si ha il corollario:\\
$f \in C^{m+1}([a, b])$, $J_n(f)$ con grado di precisione $m$ e tale che il nucleo di Peano $G(t)$ non cambia segno per $t \in [a, b]$, allora:
\[
E_n(f) = \frac{f^{(m+1)}(\varepsilon)}{(m + 1)!} E_n(x^{m+1}), \text{ per un certo } \varepsilon \in [a, b]
\]
Questa formula permette di quantificare l'errore sapendo solo:
\begin{itemize}
    \item quanto vale la derivata (m+1)-esima di $f$
    \item l'errore fatto su un polinomio noto (come $x^{m+1}$)
\end{itemize}


\begin{lezione}
    [Lezione 17 29/04]
\end{lezione}

\subsection{Recap Trapezi/Peano}
...
\section{Formule Interpolatorie}
\begin{nota}
    È una sottoclasse specifica delle formule di quadratura, l'idea è quella di approssimare l'integrale tramite l'integrale del polinomio interpolante della funzione, in quanto integrare un polinomio risulta banale. 
\end{nota}
% L'\textbf{idea} è: invece di approssimare direttamente l'integrale di f, si interpola f con un polinomio $P_n(x)$ e si integra quello, dato che integrare un polinomio è banale. I pesi, a questo punto, vengono di conseguenza, non sono più liberi.\\
% Prendendo in esempio Lagrange:
% \[
% \int_{a}^{b}l_j(x)\rho(x)dx
% \]
% abbiamo $l_j$ che sono i polinomi di Lagrange, mentre i nodi della formula coincidono esattamente con i nodi di interpolazione.
% % Ora invece di approssimare direttamente l'integrale di $f$ interpoliamo $f$ con un polinomio e integriamo il polinomio.\\
\noindent
Ripartendo dal problema originale abbiamo da calcolare $\displaystyle\int_{a}^{b} f(x)\rho(x)\,dx \approx \sum \dots$, selezioniamo quindi i punti $a \leq x_0 < x_1 < \ldots < x_n \leq b$ come $n+1$ nodi di interpolazione in $[a, b]$.\\
Si considera $P_n(x)$ il polinomio di grado al più $n$ che interpola $f$ nei nodi 
\[
x_j:\, f(x_j) = P_n(x_j)\quad \forall j = 0,\ldots,n
\]
e scriviamo $f(x)$ come somma tra il polinomio interpolante e il suo resto:
$$f(x) = P_n(x) + R_n(x)$$
Sostituendo questa espressione nell'integrale si ha:
\[
\int_{a}^{b} \bigl(P_n(x) + R_N(x)\bigr)\rho(x)\,dx
= \int_{a}^{b} P_n(x)\rho(x)\,dx + \int_{a}^{b} R_n(x) + \rho(x)\,dx
\]
Prendiamo il primo termine della somma come approssimazione dell'integrale iniziale ed il secondo come l'errore commesso.
\\
Sappiamo, dalla \autoref{def: polinomio interpolante di Lagrange}, che $P_n(x) = \displaystyle\sum_{j=0}^{n} f(x_j) \cdot l_j(x)$, con $l_j = \displaystyle\prod_{\substack{i \neq j,\, i=0}}^{n} \dfrac{(x - x_i)}{x_j - x_i}$\\
Per cui:
\[
\int_{a}^{b} P_n(x)\rho(x)\,dx
= \int_{a}^{b} \sum_{j=0}^{n} f(x_j) l_j(x)\rho(x)\,dx
= \boxed{\sum_{j=0}^{n} f(x_j) \int_{a}^{b} l_j(x)\rho(x)\,dx}
\]
\begin{nota}
    La scelta di prendere Lagrange come polinomio interpolante, piuttosto che Newton o altri, è dovuta al fatto che stiamo cercando di ottenere una sommatoria riconducibile a quella che definisce una formula di quadratura, e che quindi presenti un valore di $f(x_i)$ moltiplicato per un peso.
    La formula di Lagrange contiene proprio questo tipo di struttura e quindi è perfetta per questo scopo. 
\end{nota}
\medskip
\noindent
Abbiamo quindi ottenuto una formula di quadratura in cui i nodi di interpolazione coincidono con i nodi della formula interpolatoria e in cui i pesi $a_j$ sono:
$a_j = \displaystyle\int_{a}^{b} l_j(x)\rho(x)\,dx$, $j = 0,\ldots,n$\\
Per l'errore commesso si ha: $E_n(f) = \displaystyle\int_{a}^{b} R_n(x)\rho(x)\,dx$ e per $f \in C^{n+1}([a,b])$, dal \autoref{th: errore interpolazione polinomiale}, sappiamo che:
\[
R_n(x) = \frac{\pi(x)}{(n+1)!}\,f^{(n+1)}(\varepsilon_x),
\quad \text{con } \pi(x) = \prod_{j=0}^{n}(x - x_j)
\text{ e } \varepsilon_x \in [a, b]
\]
\paragraph{Osservazione}
Se non si può utilizzare il corollario del teorema di Peano (perché $G(t)$ cambia di segno si può utilizzare $R_n(x)$ per provare a stimare l'errore.

\begin{teorema}
    [Grado di precisione di una formula interpolatoria]
    Per costruzione, una formula interpolatoria su $n+1$ nodi è esatta su tutti i polinomi di grado $n$, quindi tale formula ha grado di precisione $\geq n$.
\end{teorema}
\begin{teorema}
    [Corrispondenza tra formula di quadratura e interpolatoria]
    Dato il \autoref{th: unicita formula quadratura grado massimo}, fissando dei nodi (sempre $n+1$) esiste un'unica formula di quadratura che ha grado di precisione $\ge n$, quindi \uline{fissati i nodi} l'unica formula di quadratura che ha grado di precisione \textbf{almeno} $n$ è la formula interpolatoria su questi nodi.
\end{teorema}
\begin{nota}
    Concetto importante per i vero/falso.
\end{nota}
\noindent
Per concludere: fissati i nodi $x_0, ..., x_n$ l'unica formula che ha grado di precisione almeno $n$ è la formula interpolatoria su questi nodi.\\
In base a come scegliamo i nodi diamo nomi differenti alle formule interpolatorie, tra le quali le più importanti sono:
\begin{itemize}
    \item le formule di Newton-Cotes, caratterizzate da nodi equispaziati nell'intervallo.
    \item le formule Gaussiane, caratterizzate da nodi scelti in modo da massimizzare il grado di precisione.
\end{itemize}

\subsection{Formule di Newton-Cotes}
\label{sec:formule di newton-cotes}
Le formule di Newton-Cotes sono le formule interpolatorie che si ottengono scegliendo gli $n+1$ nodi $x_0,\dots,x_n$ \textbf{equispaziati} nell'intervallo $[a,b]$. Per il teorema precedente, fissati i nodi in questo modo, i pesi $a_j$ risultano univocamente determinati e la formula ha grado di precisione almeno $n$.\\
Le formule di \uline{Trapezi e Cavalieri-Simpson sono infatti i casi $n=1$ e $n= 2$} delle formule di Newton-Cotes.

Ricapitolando: Nodi equidistanti di $h$ con $h= \frac{b-a}{n}$ (passo della formula di quadratura) e dove $n$ numero di nodi.

\subsubsection{Come si trovano i pesi}
I pesi si trovano valutando con $\rho(x) = 1$:
\[
a_i = I(l_i(x)) = \int_a^b l_i(x)dx
\]
oppure risolvendo il sistema lineare derivante da
\[
E_p(1)=E_n(x)=\dots = E_n(x^n) =0
\]
\begin{nota}
    [Proprietà]
    fino a $n=7$ i pesi $a_j$ risultano essere positivi, si hanno poi segni alternati.
\end{nota}
\begin{nota}
    [Proprietà]
    il nucleo di Peano delle formule di Newton-Cotes non cambia segno in $[a,b]$.
\end{nota}

\subsection{Errore nelle formule dei Trapezi e Simpson}
\paragraph{Trapezi}
\[
E_1 = - \frac{(b-a)^3}{12} \cdot f^{(2)}(\varepsilon) \quad\varepsilon\in [a,b]
\]
\paragraph{Simpson}
\[
E_2 = - \frac{(b-a)^5}{2880} \cdot f^{(4)}(\varepsilon) \quad\varepsilon\in [a,b]
\]
\paragraph{Osservazione}
Le formula di N.C. se definite escludendo $a$ e $b$ si dicono "aperte".

\subsection{Polinomi Ortogonali}
\begin{definizione}
    [Spazio vettoriale con prodotto scalare]
    Sia $\Pi = \{\text{polinomi con coefficienti reali}\}$ e consideriamo il prodotto scalare
    $p(x), q(x) \in \Pi$,
    \[
    \langle p(x), q(x) \rangle = \int_{a}^{b} p(x)q(x)\rho(x)\,dx,
    \]
    questo è uno spazio vettoriale, con prodotto scalare.
\end{definizione}
\begin{definizione}
    [Famiglia di polinomi ortogonali]
    \[
    \Pi^* = \left\{ a_i(x) \in \Pi,\, \deg(q_i) = i,\,
    \langle q_i, q_j \rangle =
    \begin{cases} 0 & \text{se } i \neq j \\ c_i & \text{se } i = j \end{cases}
    \right\}_{\!i=0,1,2\ldots}
    \]
    si dice famiglia di polinomi ortogonali su $[a, b]$ rispetto al peso $\rho(x)$.

    In parole povere: una successione di polinomi $q_0, q_1, q_2, \ldots$ dove $q_i$ ha grado esattamente $i$ e ogni coppia distinta è "perpendicolare" nel senso del prodotto scalare pesato $\langle f, g \rangle = \int_a^b f(x)g(x)\rho(x)\,dx$ - stesso concetto di ortogonalità tra vettori, ma applicato a funzioni.
\end{definizione}

\paragraph{Proprietà}
\begin{itemize}
  \item gli elementi di $\Pi^*$ sono univocamente determinati a meno di costanti
  \item I polinomi in $\Pi^*$ di grado $\leq j$, generano tutti i polinomi di grado al più $j$
  \item $\langle p(x), q_n(x) \rangle = 0$, $\quad \forall p(x)$ di grado $\leq n-1$
  \item $q_n(x)$ ha $n$ zeri semplici contenuti in $[a, b]$ $\forall n \in \mathbb{N}$.
        Questi coincidono con i nodi della formula gaussiana a $n$ nodi,
        per l'intervallo $[a, b]$ rispetto al peso $\rho(x)$
\end{itemize}

\paragraph{Lemma}
$q_n(x)$ polinomio ortogonale su $[a, b]$ con peso $\rho(x)$,
allora $q_n(x)$ ha $n$ zeri semplici in $[a, b]$.
\begin{dimostrazionelibera}
    sul quaderno.
\end{dimostrazionelibera}

\begin{lezione}
    [Lezione 18 03/05]
\end{lezione}

\subsection{Come migliorare l'accuratezza}
\begin{enumerate}
    \item aumentando i nodi nella formula di Newton-Cotes.
    \item spezzando $[a,b]$ in sottointervalli più piccoli e stimando gli integrali nei sottointervalli singolarmente.
\end{enumerate}
Poiché la proposta del punto 1 è limitata dal fatto che per $n>7$ la formula di Newton Cotes diventa instabile numericamente di solito si predilige procedere per la seconda strada.
L'idea è infatti che se $c\in[a,b]$ allora:
\[
\int_a^bf(x)\,dx = \int_a^cf(x)\,dx + \int_c^bf(x)\,dx 
\]
dove il primo integrale della somma viene stimato con una formula di Newton Cotes mentre il secondo tramite un altro metodo.\\
Formalmente: se vogliamo $L$ sottointervalli \uline{si scelgono $L+1$ punti equispaziati in $[a,b]$}:
\[
\int_a^b f(x)\,dx = \sum_{i=1}^L f(x)\,dx \approx \sum_{i=1}^L \overset{[x_{i-1}, x_i]}{J_n(f)}
\]
dove $J_n(f)$ è una formula di Newton Cotes con $n+1$ nodi.\\
Questo procedimento ci costringe a valutare $f$ nei punti $x_0, \dots, x_L$ ed inoltre in eventuali altri punti interni ai singoli intervalli, quanti nodi servono in generale?\\
Se usiamo la formula di N.C. con $n+1$ nodi ci servono $n+1-2 = n-1$ punti aggiuntivi in ogni intervallo (i due estremi saranno già parte delle $L$ valutazioni).
Quindi in totale abbiamo come \textbf{numero totale di nodi}:
\[
(n-1)L + L+1 = \boxed{nL+1}
\]
perciò ne concludiamo che: per una formula di N.C. generalizzata sono necessarie $nL+1$ valutazioni di $f(x)$ dove:
\begin{itemize}
    \item $L=$ numero di sottointervalli.
    \item $n$ ordine della formula.
\end{itemize}

Andiamo dunque a vedere le formule di Newton-Cotes generalizzate.
L'idea alla base di queste, come abbiamo appena visto, è che per rendere l'approssimazione dell'integrale più accurata si può suddividere l'intervallo in L sottointervalli di uguale ampiezza e applicare una formula di Newton-Cotes in ciascuno di essi.
\section{Formule generalizzate o composite}
Dati $L$ intervalli, definiti da un'ampliezza $h=\frac{b-a}{L}$, si hanno:
\begin{definizione}
    [Formula dei Trapezi generalizzata]
    Generalizzazione di \autoref{def: trapezi} su $L$ sottointervalli:
    \[
    J_1^{(G)} = \frac{b-a}{2L} \cdot \left[ f(x_0) +2\sum_{i=1}^{L-1}f(x_i) + f(x_L)\right]
    \]
\end{definizione}
\begin{definizione}
    [Formula di Cavalieri-Simpson generalizzata]
    Generalizzazione di \autoref{def: simpson} su $L$ sottointervalli che quindi risultano in $nL+1$ valutazioni\footnote{Ricordarsi come Cavalieri-Simpson utilizzava 3 punti rispetto a Trapezi che ne utilizzava solo 2}:
    \[
    J_2^{(G)} = \frac{b-a}{6L} \cdot \left[ f(x_0) +2\sum_{i=1}^{L-1}f(x_i) +4\sum_{i=1}^{L}f\left(\frac{x_i+x_{i-1}}{2}\right) + f(x_L)\right]
    \]
    anche se si capisce meglio ragionando con indici pari o dispari rispetto alle $nL+1$ valutazioni:
    \[
        J_2^{(G)} = \frac{b-a}{6L} \cdot \left[ f(x_0) +2\sum_{i=1}^{L-1}f(x_{2i}) +4\sum_{i=0}^{L-1}f\left(x_{2i+1}\right) + f(x_L)\right]
    \]
\end{definizione}
\subsection{Errore delle formule generalizzate}
Il \textbf{grado di precisione} delle formule generalizzate \textbf{coincide con} quello delle \textbf{formule base}.\\
L'\textbf{errore} invece \textbf{migliora} (diminuisce) \textbf{al crescere di} $L$, esso vale infatti:
\[
\dots \Longrightarrow E^{(G)}(f) = c\, h^\alpha  L  f^{(\alpha -1)} (\varepsilon)
\]
dove:
\begin{itemize}
    \item $c$: costante numerica dipendente dalla formula base usata;
    \item $h = \dfrac{b-a}{L}$: ampiezza di ogni sottointervallo (passo di discretizzazione);
    \item $\alpha$: ordine di convergenza dell'errore rispetto a $h$;
    \item $L$: numero di sottointervalli in cui è diviso $[a,b]$;
    \item $f^{(\alpha-1)}(\varepsilon)$: derivata di ordine $\alpha{-}1$ di $f$ valutata in un punto $\varepsilon \in [a,b]$ (esiste per il teorema del valor medio, assumendo $f$ sufficientemente regolare).
\end{itemize}
Da cui:
\begin{definizione}
    [Errore delle formule di Newton-Cotes generalizzate]
    \[
    E_n^{(G)} =
    \begin{cases}
    c_n \, \dfrac{(b-a)^{n+2}}{L^{\,n+1}} \, f^{(n+1)}(\varepsilon)
    & \text{se $n$ dispari}, \\[8pt]
    c_n \, \dfrac{(b-a)^{n+3}}{L^{\,n+2}} \, f^{(n+2)}(\varepsilon)
    & \text{se $n$ pari}.
    \end{cases}
    \]
    Da cui possiamo definire che se abbiamo un numero di nodi $k+1$ allora in base a $n=k-1$ possiamo definire in modo più preciso il grado di precisione: dall'espressione sopra vediamo che con \uline{$n$ pari} la formula "sente" solo la derivata di ordine $n+2$, il che significa che \uline{è esatta anche per polinomi di grado $n+1$}.
\end{definizione}
\paragraph{Trapezi}
\[
E_1^{(G)} = - \frac{(b-a)^3}{12L^2} \cdot f^{(2)}(\varepsilon) \quad\varepsilon\in [a,b]
\]
\paragraph{Simpson}
\[
E_2^{(G)} = - \frac{(b-a)^5}{2880L^4} \cdot f^{(4)}(\varepsilon) \quad\varepsilon\in [a,b]
\]

\section{Formule di quadratura che coinvolgono derivate di f}
Si possono generalizzare le formule di quadratura con espressioni che contengono derivate.
\[
\int_a^b f(x)\,dx \approx \sum_{i=0}^{n_1} f(x_i)a_i + \sum_{i=0}^{n_2} f'(y_i)b_i + \ldots + \sum_{i=0}^{n_d} f^{(n_d-1)}(z_i)c_i
\]
Dati i nodi in cui si conosce il valore della funzione e/o delle sue derivate, ci si può chiedere qual'è la scelta dei pesi che massimizza il grado di precisione della formula. Si procede nella maniera usuale determinando i pesi in modo da avere $E(1) = E(x) = \ldots = E(x^m) = 0$ per $m$ più grande possibile.\\
Il problema da risolvere è un sistema lineare dove le incognite sono i pesi della formula.\\

\chapter{Metodi per equazioni non lineari}

\begin{lezione}
    [Lezione 19 06/05]
\end{lezione}
\noindent
In generale non esiste un metodo diretto che permette di calcolare una o tutte le soluzioni di $f(x) = 0$.
\section{Intro}
\begin{definizione}
    [Radice o Zero]
    Soluzioni di $f(x) = 0$.
\end{definizione}
\noindent
Per gradi superiori al quinto di $f$ non ci sono formule che le esprimono tali zeri in termini di funzioni elementari dei coefficienti, dunque per trovare gli zeri di $f$ non lineare si usano metodi iterativi che definiamo come:


\begin{definizione}
    [Metodo iterativo per equazioni non lineari]
    Si dice tale un metodo definito come:
    \[
    \begin{cases}
    x_0,\, x_{-1},\, \ldots,\, x_{1-k} \quad \text{dati}, \\[4pt]
    x_{n+1}
    = \phi_n \bigl( x_n,\, x_{n-1},\, \ldots,\, x_{n-k+1} \bigr),
    \qquad k = 0,1,2,\ldots
    \end{cases}
    \]
\end{definizione}

\begin{definizione}
    [Metodo stazionario]
    Metodo iterativo che utilizza la stessa funzione ad ogni passo.
\end{definizione}
\begin{definizione}
    [Successione convergente]
    La successione $\{x_n\}$ è detta convergente ad $\alpha$ (radice di $f$) se:
    \[
    \lim_{n\to+\infty} x_n = \alpha
    \]
\end{definizione}
\begin{definizione}
    [Successione convergente con ordine p]
    $p\geq 1$ tale che:
    \[
    \lim_{n\to\infty} \frac{|e_{n+1}|}{|e_n|^p}=c
    \]
    dato $e_n = x_n -\alpha$ e $c\neq 0$.
\end{definizione}
\begin{nota}
    Nei casi di $p=1$ o $p = 2$ si parla di convergenza lineare e quadratica.
\end{nota}
\begin{nota}
    Nel caso $p = 1$ per la convergenza è necessario $0<c<1$
\end{nota}
\begin{nota}
    Interpretazione del significato del limite:
    \begin{itemize}
        \item per avere un $c\to 0$ dovrebbe voler dire che al denominatore abbiamo un valore che, come ordine di grandezza, tende a $\infty$ e che fa tendere il rapporto a $0$. Se $c$ rimane $\ne 0$ significa che il rapporto tra ordini di grandezza dell'errore al passo successivo e la potenza $p$ dell'errore precedente è costante, da questo la definizione di ordine di convergenza $p$.
        \item se c fosse zero allora vorrebbe dire che l'errore va a zero più velocemente di $|e_n|^p$ e quindi l'ordine, per definizione, sarebbe maggiore di p.
    \end{itemize}
\end{nota}

In generale metodi iterativi convergono solo localmente, è quindi importante individuare gli intervalli per partire sufficientemente vicino le radici.\\
Per questo si possono utilizzare strumenti di analisi matematica ad esempio:
\begin{enumerate}
    \item $f \in C^0(\mathbb{R})$ e $a, b \in \mathbb{R} : f(a)f(b) < 0$ allora $\exists$ almeno una radice $\alpha \in [a, b]$
    \item Studio del segno della derivata, \textbf{esempio}:
    Sia $f(x) = xlog(x) + \frac{1}{3}$ su $(0,+\infty)$.
    Vogliamo mostrare che $f$ ha esattamente 2 radici in $(0, +\infty)$.

    \medskip
    Si osserva innanzitutto che: $\lim_{x \to 0^+} f(x) = \frac{1}{3}$\\
    La derivata è $f'(x) = \log(x) + 1$, da cui:
    \[
        \begin{cases}
            f'(x) > 0 & x > e^{-1} \\
            f'(x) < 0 & x < e^{-1}
        \end{cases}
    \]
    Quindi $f$ è decrescente su $(0, e^{-1})$ e crescente su $(e^{-1}, +\infty)$,
    con minimo in $x = e^{-1}$. Valutiamo $f$ nel minimo:
    \[
        f(e^{-1}) = -\frac{1}{e} + \frac{1}{3} < 0
    \]
    Poiché $f(0^+) = \frac{1}{3} > 0$, $f(e^{-1}) < 0$ e $f(x) \to +\infty$
    per $x \to +\infty$, per il teorema degli zeri esistono esattamente
    2 radici $\alpha_1 < \alpha_2$ con:
    \[
        \alpha_1 \in (0,\, e^{-1}), \qquad \alpha_2 \in (e^{-1},\, +\infty)
    \]
    \item Separazione grafica: se $f(x) = g(x) - h(x)$ dove $g(x)$, $h(x)$ hanno grafici o proprietà note, potrebbe essere conveniente riscrivere $f(x) = 0$ come $g(x) = h(x)$, ovvero cercare le intersezioni fra grafico di $g$ e grafico di $h\to$  possiamo quindi identificare un intervallo $[a, b]$ in cui sappiamo esserci una radice.
    \begin{figure}[H]
        \centering
        \includegraphics[width= 0.65\textwidth]{img/Separazione grafica.png}
    \end{figure}
\end{enumerate}

\begin{metodo}
    [Metodo di Bisezione (ricerca binaria)]
    Data $f$ continua su $[a,b]$ e tale che $f(a)\cdot f(b) < 0$ e considerando $x_0, x_1$ come gli estremi mentre $x_2=\frac{b-a}{2}$ il punto medio si ricava il metodo:
    \[
    \hat{x}_n = 
    \begin{cases}
        \text{dipende dal sengo}
    \end{cases}
    \]
    che quindi è un metodo a 2 punti non stazionario in quanto la funzione può cambiare ad ogni passo.
\end{metodo}
\begin{metodo}
    [Metodo delle secanti]
    Metodo stazionario a 2 punti
    \[
    \begin{cases}
        x_{n+1} = x_n - f(x_n) \cdot \frac{x_n - x_{n-1}}{f(x_n) - f(x_{n-1})} \\
        x_0,x_1 = dati
    \end{cases}
    \]
\end{metodo}
\begin{teorema}
    [Ordine di convergenza del metodo delle secanti]
    $f \in C^2([a, b])$, e partendo sufficientemente vicino ad $\alpha$, allora il metodo converge con ordine dato da
    \[p = \frac{1 + \sqrt{5}}{2}\]
\end{teorema}

\section{Metodi stazionari a un punto (metodi di punto fisso)}
\label{sec: punto fisso}

\subsection{Definizione e convergenza locale}
Il metodo del punto fisso prevede di cercare lo zero di una funzione $f$ attraverso la ricerca del punto fisso di un'altra funzione, chiamiamola $\phi$, ottenuta manipolando la funzione di partenza $f$ .\\
In sostanza, non esistendo un metodo risolutivo per questo tipo di funzioni, ci costruiamo un metodo iterativo a partire dalla funzione stessa.\\
Se l'equazione $f (x) = 0$ è riscritta nella forma $x = \phi(x)$, è possibile utilizzare un metodo iterativo di punto fisso di tale forma:
\begin{definizione}
    [Metodo iterativo a punto fisso]
    Un metodo iterativo a punto fisso è definito dalla successione:
    \[
    \begin{cases}
        x_0 = dato\\
        x_{n+1} = \phi(x_n)
    \end{cases}
    \]
\end{definizione}
\noindent
e sperare che la successione $\{x_n\}$ converga ad una radice di $f$.\\
Il \textbf{ragionamento} è questo: costruiamo $\phi$ in modo tale che i suoi punti fissi coincidano con gli zeri di $f$ (lo facciamo, come visto sopra, partendo da $f (x) = 0$), ovvero che la funzione $\phi$ incontri la bisettrice primo-terzo quadrante in corrispondenza delle $x$ in cui $f$ presenta uno zero. Se nella regione che ci interessa c'è un solo punto fisso, allora l'iterazione
$x_{n+1} = \phi(x_n)$ 
partendo da un $x_0$ qualsiasi vicino a quel punto, converge esattamente lì.\\
% Il metodo non "usa" i punti fissi come input ma come destinazione: si parte da $x_0$ generico, si applica $\phi(x_n)$ ripetutamente, e la successione, se convergente, viene attrattaverso il punto fisso.
Di seguito un \textbf{esempio} con $f(x) = x- cos(x)$:
% Richiede nel preambolo: \usepackage{subcaption} (oltre a \usepackage{graphicx})
\begin{figure}[H]
    \centering
    \begin{subfigure}[t]{0.32\textwidth}
        \centering
        \includegraphics[width=\linewidth]{img/step1_f.png}
        \caption{Il problema $f(\alpha)=0$: $f(x)=x-\cos(x)$ non ha formula chiusa.}
    \end{subfigure}
    \hfill
    \begin{subfigure}[t]{0.32\textwidth}
        \centering
        \includegraphics[width=\linewidth]{img/step2_cobweb0.png}
        \caption{Riscrittura come punto fisso: $x=\cos(x)$, $\varphi(x)=\cos(x)$, punto di partenza $x_0=0.1$.}
    \end{subfigure}
    \hfill
    \begin{subfigure}[t]{0.32\textwidth}
        \centering
        \includegraphics[width=\linewidth]{img/step3_cobweb1.png}
        \caption{Prima iterazione: $x_1=\cos(x_0)\approx 0.995$.}
    \end{subfigure}
    % \caption{Dal problema $f(x)=0$ al metodo di punto fisso: costruzione del diagramma a ragnatela (cobweb).}
    \label{fig:cobweb-steps}
\end{figure}
% Richiede nel preambolo: \usepackage{subcaption} (oltre a \usepackage{graphicx})
\begin{figure}[H]
    \centering
    \begin{subfigure}[t]{0.32\textwidth}
        \centering
        \includegraphics[width=\linewidth]{img/step4_cobweb2.png}
        \caption{Seconda iterazione: $x_2=\cos(x_1)\approx 0.549$. La successione salta dall'altro lato di $\alpha$.}
    \end{subfigure}
    \hfill
    \begin{subfigure}[t]{0.32\textwidth}
        \centering
        \includegraphics[width=\linewidth]{img/step5_cobweb3.png}
        \caption{Terza iterazione: $x_3\approx 0.857$. Le oscillazioni si restringono poiché $|\varphi'(\alpha)|\approx 0.67<1$.}
    \end{subfigure}
    \hfill
    \begin{subfigure}[t]{0.32\textwidth}
        \centering
        \includegraphics[width=\linewidth]{img/step6_cobweb4.png}
        \caption{Quarta iterazione: $x_4\approx 0.655$. La spirale converge verso $(\alpha,\alpha)$.}
    \end{subfigure}
    \caption{Prosecuzione del metodo di punto fisso: la ragnatela si stringe progressivamente attorno ad $\alpha$ (convergenza lineare, fattore $\approx 0.67$ ad ogni passo).}
    \label{fig:cobweb-steps-2}
\end{figure}

\paragraph{Osservazione}
Partendo da una certa funzione $f$ si possono derivare infinite funzioni $\phi$ che sono la riscrittura di $f$ nella forma di punto fisso, ma non è detto che ognuna di queste funzioni porti ad una successione convergente.
\begin{definizione}
    [Punto fisso]
    Elemento di una funzione che coincide con la sua immagine.
\end{definizione}
\noindent
Notiamo che cercare le radici  di $f$ è equivalente a cercare i punti fissi di $\phi(x)$, pertanto, se la successione converge allora necessariamente converge a un punto fisso di $\phi$.\\
\noindent
Quando funziona? lo dimostra il prossimo teorema, questo spiega infatti quando e perché un metodo iterativo converge se si parte abbastanza vicino alla soluzione, e con quale velocità.
\begin{teorema}
    [Teorema di convergenza locale]
    Dati $I\subseteq \mathbb{R}$ intervallo, $\alpha \in I$, $\phi(\alpha) = \alpha$, $\phi \in \mathbb{C}^1(I)$ e supponiamo che $\exists \rho \in\mathbb{R}^+ $ e $k\in(0,1)$ tali che:
    \[
    |\phi'(x)| \leq k \quad\forall x \in [\alpha - \rho, \alpha + \rho] 
    \]
    allora :   
    \begin{enumerate}
        \item $\forall x_0 \in [\alpha - \rho, \alpha + \rho] \Rightarrow x_n \in [\alpha - \rho, \alpha + \rho] \quad\forall n\in\mathbb{N}$
        \item $\forall x_0 \in [\alpha - \rho, \alpha + \rho]$, $\lim_{n\to\infty} x_n = \alpha$
        \item $\alpha$ è l'unico punto fisso di $\phi$ in $[\alpha - \rho, \alpha + \rho]$
    \end{enumerate}
\end{teorema}
\begin{dimostrazione}
    $ $
    \begin{enumerate}
        \item Procedo per induzione su n
        \[
        n = 0
        \]
        supponiamo la tesi vera per $x_0,x_1,..., x_n$
        \[
        |x_{n+1} - \alpha| = |\phi(x_n) - \phi(\alpha)| \underbrace{=}_{Th. Lagrange} |x_n-\alpha | \cdot |\phi'(\varepsilon)|  \leq k|x_n - \alpha| \leq \underbrace{k}_{(0,1)} \cdot \rho \Rightarrow
        \]
        \[
        \Rightarrow x_{n+1} \in [\alpha - \rho, \alpha + \rho]
        \]
        $\xi \in [x_n, \alpha]$, in particolare $\xi \in [\alpha - \rho, \alpha + \rho]$
        \item $(x_{n+1} - \alpha) = |x_m - \alpha| \cdot |\phi'(\xi)| \leq k |x_n-\alpha| =$ 
        \[
        = k |\phi(x_{n-1})-\phi(x)| \leq k\phi'(\tilde{\xi})\cdot |x_{n-1}-\alpha| \leq k^2|x_{n-1} -\alpha|\leq ... \leq k^n|x_0 - \alpha | \leq k^n\rho
        \]
        da cui
        \[
        0 \le \lim_{n\to + \infty} |x_{n+1} - \alpha| 
        \le \lim_{n\to + \infty} k^n - \rho = 0 
        \implies \lim_{n\to + \infty} |x_{n+1} - \alpha| = 0
        \]
        \item supponiamo che assurdo che esista 
        $\alpha_1, \alpha_2 \in [\alpha - \rho, \alpha + \rho]$
        allora:
        \[
            \phi(\alpha_1) = \alpha_1 \, , \quad \phi(\alpha_2) = \alpha_2
        \]
        allora:
        \[
        | \alpha_1 - \alpha_2 | = |\phi(\alpha_1) - \phi(\alpha_2) |
        = | \phi'(\varepsilon) | \cdot |\alpha_1 - \alpha_2|
        \le k |\alpha_1 - \alpha_2|
        < |\alpha_1 - \alpha_2|
        \]
        che risulta essere impossibile (guardando il primo e l'ultimo termine vediamo come la disequazione è impossibile).
        \end{enumerate}
\end{dimostrazione}
\begin{lezione}
    [Lezione 20 10/05]
\end{lezione}

\begin{teorema}
    [Teorema di Convergenza Locale (semplificato)]
    Se $\phi\in \mathcal{C}^1$ e $|\phi'(\alpha)| < 1$ allora partendo sufficientemente vicini ad $\alpha$ il metodo di punto fisso converge.
\end{teorema}

\subsection{Convergenza monotona/alternata}
\begin{definizione}
    [Convergenza monotona]
    Se $\phi'>0$ in un intorno destro o sinistro (o entrambi) di $\alpha$ allora la convergenza è monotona.
\end{definizione}
\begin{definizione}
    [Convergenza alternata]
    Se $\phi'<0$ intorno ad $\alpha$
    \[
    (x_{n+1} - \alpha) = (x_n - \alpha)\cdot \underset{(<0)}{\phi'(\xi)}
    \]
    allora il segno di $x_{n+1} - \alpha$ è l'opposto di quello di $x_n - \alpha$, quindi è necessario che $|\phi'|<1$ in un intervallo simmetrico rispetto alla radice per garantire la convergenza:
    \[
    \phi'<0 \text{ in } (\alpha - \rho, \alpha + \rho) \text{ allora:}
    \]
    \[
    x_0 < \alpha, \; x_1 > \alpha, \; x_2 < \alpha, \; x_3 > \alpha, \;\dots
    \]
    dunque la convergenza è alternata.
\end{definizione}
\subsection{Quanto in fretta converge un metodo di punto fisso?}
\begin{teorema}
    [Velocità (ordine) di convergenza di un metodo di punto fisso]
    Sia $\phi(x) \in C^{p}(I)$, $\alpha \in I$ tale che $\phi(\alpha) = \alpha$.
Allora $\exists \,\rho > 0$ tale per cui, se
$x_0 \in [\alpha - \rho, \alpha + \rho]$,
la successione $\{x_n\}_{n \in \mathbb{N}}$ converge
superlinearmente con ordine $p$ se e solo se ($\Leftrightarrow$)
\[
\phi'(\alpha) = \phi''(\alpha) = \cdots = \phi^{(p-1)}(\alpha) = 0,
\qquad
\phi^{(p)}(\alpha) \neq 0.
\]
\end{teorema}

\subsubsection{Note sulle convergenze}
\begin{itemize}
    \item
    $\{x_n\}$ converge superlinearmente con ordine $p > 1$ se
    \[
    \lim_{n \to \infty}
    \frac{|x_{n+1} - \alpha|}{|x_n - \alpha|^{p}}
    = c < \infty,
    \qquad c \neq 0.
    \]

    \item
    Quando $p = 1$ si richiede $0 < c < 1$:
    convergenza lineare.

    \item
    Quando $p = 1$ e $c = 1$ (e vale comunque $\lim_{n \to \infty} x_n = \alpha$):
    convergenza sublineare.
\end{itemize}

Inoltre,
\[
|\phi'(\alpha)| = 1,
\qquad
|\phi'(x)| < 1
\quad \text{in } [\alpha - \rho, \alpha + \rho].
\]
\begin{dimostrazione}
    \underline{($\Leftarrow$)} Per dimostrare questo assumiamo l'altra proposizione:
    \[
    \varphi'(\alpha) = \cdots = \varphi^{(p-1)}(\alpha) = 0, \qquad \varphi^{(p)}(\alpha)\neq 0
    \]
    Consideriamo lo sviluppo di Taylor di $\varphi$ in $\alpha$:
    \[
    x_{n+1} = \varphi(x_n) = \varphi(\alpha) + (x_n-\alpha)\varphi'(\alpha) + \cdots + \varphi^{(p-1)}(\alpha)\frac{(x_n-\alpha)^{p-1}}{(p-1)!} + \frac{(x_n-\alpha)^p}{p!}\varphi^{(p)}(\varepsilon), \quad \varepsilon\in[x_n,\alpha]
    \]
    \[
    \implies x_{n+1} = \alpha + \frac{(x_n-\alpha)^p}{p!}\varphi^{(p)}(\varepsilon)
    \]
    \[
    \implies x_{n+1}-\alpha = \frac{(x_n-\alpha)^p}{p!}\varphi^{(p)}(\varepsilon), \qquad \varepsilon\in[x_n,\alpha] \xrightarrow{n\to+\infty} \alpha
    \]
    e quindi:
    \[
    \lim_{n\to+\infty}\frac{|x_{n+1}-\alpha|}{|x_n-\alpha|^p} = \lim_{n\to+\infty}\frac{\varphi^{(p)}(\varepsilon)}{p!} = \frac{\varphi^{(p)}(\alpha)}{p!} \neq 0
    \]

    \bigskip
    \noindent
    \underline{($\Rightarrow$)} Prendiamo $1\leq r\leq p-1$ un $r$ intero, facciamo vedere che $\varphi^{(r)}(\alpha)=0$.

    \medskip
    \noindent
    \textbf{Per $r=1$:}
    \[
    \lim_{n\to+\infty}\frac{|\varphi(x_n)-\varphi(\alpha)|}{|x_n-\alpha|} = \lim_{n\to+\infty}\frac{|x_n-\alpha|\,\varphi'(\varepsilon)}{|x_n-\alpha|} = \lim_{n\to+\infty}\varphi'(\varepsilon) = \varphi'(\alpha)
    \]

    \medskip
    \noindent
    \textbf{Per induzione} supponiamo che $\varphi'(\alpha)=\varphi''(\alpha)=\cdots=\varphi^{(r-1)}(\alpha)=0$ (ipotesi induttiva). Allora:
    \[
    \varphi(x_n)-\varphi(\alpha) = \frac{(x_n-\alpha)}{1!}\varphi'(\alpha) + \frac{(x_n-\alpha)^2}{2!}\varphi''(\alpha) + \cdots + \frac{(x_n-\alpha)^{r-1}}{(r-1)!}\varphi^{(r-1)}(\alpha) + \frac{(x_n-\alpha)^r}{r!}\varphi^{(r)}(\varepsilon)
    \]
    \[
    = \frac{(x_n-\alpha)^r}{r!}\varphi^{(r)}(\varepsilon)
    \]
    % NOTA A MARGINE (tua): "perché la succ. x_n converge con ordine p≥2" -- 
    % giustifica il passaggio successivo (probabilmente: essendo l'ordine di
    % convergenza p, per r<p il rapporto |x_{n+1}-\alpha|/|x_n-\alpha|^r tende a 0).
    \[
    0 = \lim_{n\to+\infty}\frac{|x_{n+1}-\alpha|}{|x_n-\alpha|^r} = \lim_{n\to+\infty}\frac{|\varphi(x_n)-\varphi(\alpha)|}{|x_n-\alpha|^r} = \lim_{n\to+\infty}\frac{\varphi^{(r)}(\varepsilon)}{r!} = \frac{\varphi^{(r)}(\alpha)}{r!}
    \]
    \[
    \implies \varphi^{(r)}(\alpha) = 0
    \]
\end{dimostrazione}

\begin{nota}
    Si possono avere ordini di convergenza $p$ non interi
    solo se $\phi \in C^{p}$.
    Ad esempio, la successione $x_{n+1} = \sqrt{|x_n|^3}$ ha ordine di convergenza $\frac{3}{2}$.
\end{nota}
\subsubsection{esempi}

\subsection{Metodo delle Tangenti}
\begin{metodo}
    [Metodo delle Tangenti (o di Newton)]
    Dati $f(x) = 0 \longrightarrow x = \phi(x), \quad f\in\mathcal{C}^1$ scegliendo
    \[
    \phi(x) = x - \frac{f(x)}{f'(x)} \longrightarrow
    \begin{cases}
        x_0 = dato \\
        x_{n+1} = x_n - \frac{f(x_n)}{f'(x_n)}
    \end{cases}
    \]
\end{metodo}
\begin{definizione}
    [Radice di ordine o molteplicità r]
    \[
    \lim_{x\to\alpha}\frac{f(x)}{(x-\alpha)^r} = c < +\infty \quad c\ne 0
    \]
\end{definizione}
\begin{definizione}
    [Radice semplice]
    Con radice semplice si indica una radice nella quale la funzione ha derivata non nulla, ovvero la funzione non si appoggia sull'asse x ma lo attraversa:
    \[
    \lim_{x\to\alpha}\frac{f(x)}{x-\alpha} = c < +\infty \quad c\ne 0
    \]
    ovvero è il caso con $r=1$ della precedente definizione.
\end{definizione}

\begin{teorema}
    [Convergenza locale del Metodo delle Tangenti]
    Se $\alpha$ è una radice semplice di $f(x)$ con $f\in\mathbb{C}^2$ allora il metodo di Newton converge localmente in modo superlineare con ordine $p\geq 2$. in particolare vale:
    \[
    \lim_{n\to \infty} \frac{|x_{n+1}- \alpha|}{|x_n- \alpha|^2}=\frac{1}{2}\frac{f''(\alpha)}{f'(\alpha)} \quad \Longrightarrow 
    \begin{cases}
        f''(\alpha) \neq 0 \to p=2\\
        f''(\alpha)=0 \to p>2
    \end{cases}
    \]
\end{teorema}

\noindent
Quaderno 3

\begin{dimostrazione}
    \[
    \varphi(x) = x - \frac{f(x)}{f'(x)} \qquad \implies \qquad \varphi'(x) = 1 - \frac{[f'(x)]^2 - f(x)f''(x)}{[f'(x)]^2}
    \]
    \[
    = \frac{f(x)f''(x)}{[f'(x)]^2}
    \qquad \implies \qquad
    \varphi'(\alpha) = \frac{\overbrace{f(\alpha)}^{0}f''(\alpha)}{[f'(\alpha)]^2} = 0
    \]
    % Nota a margine (tua): "\alpha rad. semplice \implies f'(\alpha)\neq 0 [...]"
    % -- giustifica che il denominatore [f'(\alpha)]^2 sia ben definito e non nullo.
    Poiché $\alpha$ è radice semplice, $f'(\alpha)\neq 0$: il rapporto è quindi ben definito, e $\varphi'(\alpha)=0$ implica che il metodo converge \textbf{superlinearmente} (ordine $\geq 2$).

    \bigskip
    \noindent
    \textbf{Per la 2\textsuperscript{a} parte del Teorema:}
    \[
    \varphi''(x) = \frac{(f'\cdot f'' + f\cdot f''')(f')^2 - f\cdot f''\cdot 2f'\cdot f''}{(f')^4}
    \]
    Valutando in $\alpha$ (dove $f(\alpha)=0$, quindi tutti i termini con $f(\alpha)$ a fattore si annullano):
    \[
    \implies \varphi''(\alpha) = \frac{[f'(\alpha)]^3\cdot f''(\alpha)}{[f'(\alpha)]^4} = \frac{f''(\alpha)}{f'(\alpha)}, \qquad f'(\alpha)\neq 0
    \]

    \bigskip
    \noindent
    A questo punto, applicando lo sviluppo di Taylor:
    \[
    \frac{|x_{n+1}-\alpha|}{|x_n-\alpha|^2} = \frac{|\varphi(x_n)-\varphi(\alpha)|}{|x_n-\alpha|^2} \underset{\text{Taylor}}{=} \frac{\left|\overbrace{\varphi'(\alpha)}^{0}(x_n-\alpha) + \dfrac{\varphi''(\varepsilon)(x_n-\alpha)^2}{2}\right|}{|x_n-\alpha|^2}
    \]
    \[
    = \frac{1}{2}|\varphi''(\varepsilon)| \xrightarrow[n\to+\infty]{} \frac{1}{2}\varphi''(\alpha) = \frac{1}{2}\left|\frac{f''(\alpha)}{f'(\alpha)}\right|
    \]
\end{dimostrazione}

\paragraph{Proprietà}
Se $f \in C^r([a, b])$, $\alpha \in [a, b]$ è radice di ordine $r$ se e solo se
\[
f(\alpha) = f'(\alpha) = f''(\alpha) = \ldots = f^{(r-1)}(\alpha) = 0, \quad f^r(\alpha) \neq 0
\]
Intuitivamente $\alpha$ radice di ordine $r$ significa che possiamo scrivere;
\[
f(x) = (x - \alpha)^r \cdot g(x) \text{ tale che } g(\alpha) \neq 0
\]
$f'(x) = r(x-\alpha)^{r-1}g(x) + (x-\alpha)^r g(x) \to f'(\alpha) = 0 + 0$\\
$\vdots$\\
$f^{r-1}(\alpha) = 0$\\
In particolare, se $f \in C^1$ e $\alpha$ radice semplice $\Longrightarrow f'(\alpha) \neq 0$

\paragraph{Esempio}
Supponiamo di voler calcolare $t^{-1}$, quindi una semplice divisione, con $t \in \mathbb{R}^+$ ma senza usare divisioni. Ciò è equivalente a cercare gli zeri di $f(x) = x^{-1}-t$.
\[
\phi(x) = x - \frac{f(x)}{f'(x)} = x - \frac{x^{-1} - t}{-x^{-2}} = 2x - tx^2 \implies x_{n+1} = 2x_n - tx_n^2
\]
Come scelgo $x_0$ in modo da convergere?
\[
\phi'(x) = 2 - 2xt = 2(1 - xt)
\]
\[
|\phi'(x)| < 1 \iff
\begin{cases}
1 - xt < \dfrac{1}{2} \\
1 - xt > -\dfrac{1}{2}
\end{cases}
\iff
\begin{cases}
x > \dfrac{t^{-1}}{2} \\
x < \dfrac{3}{2}t^{-1}
\end{cases}
\]
Quindi Newton converge se parto da una stima non troppo sbagliata, ovvero devo partire con
\[
x_0 \in \left[t^{-1} - \frac{t^{-1}}{2};\, t^{-1} + \frac{t^{-1}}{2}\right]
\]

\paragraph{Radici multiple}
Supponiamo $\alpha$ radice di molteplicità $s > 1$, ovvero $f(x) = g(x)(x-\alpha)^s$ con $g(x) \neq 0$, $f \in C^2$ ($f'(\alpha) = 0$). Definiamo Newton in questo caso come:
\[
\phi(x) = \begin{cases}
x - \dfrac{f(x)}{f'(x)} & \text{se } x \neq \alpha \\
\alpha & \text{se } x = \alpha
\end{cases}
\]
In questo caso si ha:
\[
\phi'(x) = 1 - \frac{(f')^2 - ff''}{(f')^2} = \frac{ff''}{(f')^2} \quad (*)
\]
$f = g(x)(x-\alpha)^s$, $f' = g'(x)(x-\alpha)^s + sg(x)(x-\alpha)^{s-1}$, $f'' = g''(x-\alpha)^s + sg'(x-s)^{s-1} + sg'(x-\alpha)^{s-1} + s(s-1)g(x-\alpha)^{s-2}$\\
E sostituendo queste ultime in $(*)$, otteniamo
\[
\phi'(\alpha) = 1 - \frac{sg(\alpha)^2}{s^2g(\alpha^2)} = 1 - \frac{1}{s} < 1; \quad \in (0,1) \,\forall s > 1
\]
$\implies$ Newton converge in maniera lineare per $s > 1$
\noindent
$\to$ se si conosce $s$, si può modificare Newton in maniera minimale per riacquisire la convergenza quadratica $\to x_{n+1} = x_n - s\dfrac{f(x_n)}{f'(x_n)}$

\begin{lezione}
    [Lezione 21 13/05]
\end{lezione}

\section{Convergenza di un metodo}
\begin{teorema}
    [Convergenza di un metodo]
    Se $\phi(x) \in \mathcal{C}^1([a,b])$ e 
    \begin{enumerate}
        \item $\phi(x) \in [a,b] \quad\forall x \in [a,b]$
        \item $|\phi'(x)| <1  \quad\forall x \in [a,b]$
    \end{enumerate}
    allora $\exists !$ punto fisso $\alpha$ di $\phi(x)$ in $[a,b]$ ed il metodo $x_{n+1} = \phi(x_n)$ per $n\geq 1 $ converge $\forall x_0 \in[a,b]$.
\end{teorema}
\begin{teorema}
    [Convergenza applicata al metodo di Newton]
    Sia $\alpha \in [a,b] \subset \mathbb{R}$ una radice semplice di
    $f \in C^{2}([a,b])$.
    Se $f'$ e $f''$ hanno segno costante in $[a,b]$,
    allora il metodo di Newton converge per ogni $x_0$ tale che
    \[
    f(x_0)\,f''(x_0) > 0
    \]
    da cui:
    \begin{enumerate}
        \item Se la funzione è concava (triste) allora il metodo di Newton converge partendo da un dato iniziale $x_0$ in cui la funzione è negativa. 
        \item Viceversa se la funzione è convessa (felice) si ha convergenza partendo da un punto corrispondente ad un valore positivo di $f$.
    \end{enumerate}
    \begin{nota}
        Prendendo il caso 2 come esempio: ho convergenza perché in una funzione convessa la tangente sta sempre sotto la curva e quindi la sottostima ad ogni passaggio.\\
        Per capire il problema dell'opposto: se fosse concava la tangente sovrastima la curva e il prossimo $x_k$ potrebbe finire oltre $\alpha$, dall'altro lato, portanto quindi a oscillazione invece che a convergenza.
    \end{nota}
\end{teorema}

\section{Equazioni non lineari in Rn}
La teoria vista finora può essere estesa al caso $f:\mathbb{R}^n\to\mathbb{R}^n$ per il problema non lineare, ovvero trovare $x\in\mathbb{R}^m$ tale che
\[
\begin{cases}
f_1(x_1, \ldots, x_m) = 0 \\
f_2(x_1, \ldots, x_m) = 0 \\
\ldots \\
f_m(x_1, \ldots, x_m) = 0
\end{cases}
\quad \text{con } f_j : \mathbb{R}^m \to \mathbb{R}, \, j = 1,...,m
\]
L'idea è ancora una volta quella di scrivere l'equazione $f(x)=0$ come un'equazione di punto fisso vettoriale:
\[
x=\phi(x) \quad \phi:\mathbb{R}^n\to\mathbb{R}^n
\]
Una scelta canonica consiste nel prendere $\phi(x) = x - G(x)f(x)$ dove $G : \mathbb{R}^n \to \mathbb{R}^{n \times n}$ 
\[
G(x) = \begin{bmatrix}
G_{11}(x) & \cdots & G_{1n}(x) \\
\cdots & \cdots & \cdots \\
G_{n1}(x) & \cdots & G_{nn}(x)
\end{bmatrix}
\quad G_{ij} : \mathbb{R}^m \to \mathbb{R}, \, \forall \, 1 \leq i, j \leq m.
\]
Si assume che $G(x)$ sia non singolare in un intorno della radice $\alpha$.\\
Si considera infine l'iterazione $x^{(k+1)} = \phi(x^{(k)})$.\\
Definiamo la successione convergente di ordine $p\ge 1$.\\
Nel caso in cui esistano e siano continue le derivate di $\phi$ rispetto alle variabili $x_1,...,x_n$ si considera la matrice Jacobiana:
\[
J\phi(x) = 
\begin{bmatrix}
\frac{\partial \phi_1(x)}{\partial x_1} & \frac{\partial \phi_1(x)}{\partial x_2} & \cdots & \frac{\partial \phi_1(x)}{\partial x_n} \\
\frac{\partial \phi_2(x)}{\partial x_1} & \frac{\partial \phi_2(x)}{\partial x_2} & \cdots & \frac{\partial \phi_2(x)}{\partial x_n} \\
\vdots & \vdots & \ddots & \vdots \\
\frac{\partial \phi_m(x)}{\partial x_1} & \frac{\partial \phi_m(x)}{\partial x_2} & \cdots & \frac{\partial \phi_m(x)}{\partial x_n}
\end{bmatrix}
\]
che porta alla generalizzazione del teorema di convergenza locale.
\begin{nota}
    [Recap: come si arriva alla condizione di convergenza]
    Teniamo a mente quanto visto per i metodi di punto fisso in $\mathbb{R}$.\\
    \textbf{Il salto a $\mathbb{R}^n$.} In più dimensioni non esiste più un'unica "pendenza" (in una dimensione questa era data dalla derivata): il Jacobiano $J\phi(\alpha)$ agisce su un vettore errore in modo diverso a seconda della direzione. Alcune direzioni possono essere contratte e altre dilatate. Ciò che conta alla lunga, dopo molte iterazioni, è la direzione che si contrae/dilata \textbf{di meno}, cioè l'autovalore di modulo massimo: il raggio spettrale $\rho(J\phi(\alpha))$. Se anche la direzione "peggiore" viene comunque contratta ($\rho<1$), tutte le altre lo saranno a maggior ragione, e l'intero vettore errore collassa verso $0$.\\
    (Idea già vista per Jacobi/Gauss-Seidel ($\rho(H)<1$)).
\end{nota}
\begin{teorema}
    [Teorema di Convergenza Locale Generalizzato]
    Sia $\phi(x) \in \mathcal{C}^1(\Omega)$ con $\Omega \in \mathbb{R}^n$ e sia $\alpha \in \Omega$ tale che $\phi(\alpha) = \alpha $, se 
    $$\boxed{\rho\left(J\phi(\alpha)\right)<1 }$$
    allora $\exists \delta > 0 : \forall x^{(0)}$ con $||x^{(0)} - \alpha||_2 \leq \delta$ e quindi \uline{la successione}:
    \[
    \begin{cases}
        x^{(0)} = dato\\
        x^{(k+1)} = \phi(x^{(k)}) \quad k>0
    \end{cases}
    \]
    \uline{converge ad $\alpha$}.
\end{teorema}
\begin{nota}
    Con le norme matriciali il teorema diventa "se norma del Jacobiano compresa tra 0 e 1 allora converge".
\end{nota}
\begin{metodo}
    [Metodo di Newton-Raphson in Rn (delle tangenti)]
    Se le funzioni $f_j$ sono derivabili e il Jacobiano è non singolare (e quindi invertibile) in un intorno Omega di una radice di f allora si può generalizzare il metodo di Newton come:
    \[
    \phi(x) = x - (Jf(x))^{-1}f(x)
    \]
    si genera quindi la successione
    \[
    \begin{cases}
        x^{(0)} = \text{dato}\\
        x^{(k+1)} = x^{(k)} - (Jf(x^{(k)}))^{-1}f(x^{(k)})
    \end{cases}
    \]
    \begin{nota}
        [Importante per la risoluzione degli esercizi]
        Nell'implementazione si formano elementi del tipo $f(x^{(k)})\in\mathbb{R}^n$ e $Jf(x^{(k)})\in\mathbb{R}^{nxn}$, quindi si risolve il sistema lineare:
        \[
        Jf(x^{(k)})\cdot y= f(x^{(k)})
        \]
        e infine $x^{(k+1)}=x^{(k)}-y$
    \end{nota}
\end{metodo}

Per quanto riguarda la \textbf{convergenza} vale un risultato analogo al caso scalare che assicura convergenza quadratica.

\begin{teorema}
    [Convergenza locale quadratica di Newton-Raphson]
    Sia $f \in C^2(\Omega)$, con $\Omega \subseteq \mathbb{R}^m$, e sia $\alpha \in \Omega$ tale che $f(\alpha) = 0$.\\
    Se $Jf(x)$ è non singolare, quindi invertibile, $\forall x \in \Omega$ allora $\exists \delta \subseteq \Omega$ intorno di $\alpha$, e $\beta \geq 0$ tale che $\forall x^{(0)} \in \delta$ la successione $\{x^{(k)}\}_{k \in \mathbb{N}}$ del metodo di Newton-Raphson soddisfa
    \[
    \|x^{(k+1)} - \alpha\|_2 \leq \beta \|x^{(k)} - \alpha\|_2^2, \quad k \geq 0
    \]
   quindi \textbf{il metodo converge quadraticamente a $\alpha$ se il Jacobiano è invertibile nella radice} (vedi sotto).
\end{teorema}
\subsubsection{Osservazioni}
\begin{itemize}
    \item questi risultati possono estendersi facilmente a norme diverse da $\|\bullet\|_2$
    \item se la radice $\alpha$ è semplice (vedi sotto) e $Jf(\alpha)$ è invertibile, il metodo converge in maniera superlineare (almeno quadratica).
\end{itemize}
\begin{definizione}
    [Radice semplice in Rn]
    Una radice si dice semplice in $\mathbb{R}^n$ se
    \[
    det(J(\alpha)) \neq 0
    \]
\end{definizione}
La convergenza locale quadratica di Newton-Raphson in $\mathbb{R}^n$ si verifica tramite questa condizione, che garantisce l'invertibilità del Jacobiano in quel punto.

\begin{lezione}
    [Lezione 22 17/05]
\end{lezione}

\section{Varianti del metodo di Newton-Raphson in Rn}
Ad ogni passo, Newton-Raphson richiede 2 operazioni che possono essere costose nel caso in cui $m$ sia grande:
\begin{enumerate}
    \item calcolare $Jf(x^{(k)})$ che richiede di valutare $m^2$ funzioni non lineari da $\mathbb{R}^m \to \mathbb{R}$
    \item Risolvere il sistema lineare $Jf(x^{(k)})d = f(x^{(k)})$ che poi da $x^{(k+1)} = x^{(k)} - d$, il quale come sappiamo costa $O(n^3)$ ad esempio con Gauss
\end{enumerate}
Per questo motivo ci sono contesti in cui conviene rinunciare alla garanzia di convergenza di Newton-Raphson per passare a metodi simili ma con iterazioni meno costose.
\begin{nota}
    Ricordiamo che: se la radice $\alpha$ è semplice e $Jf(\alpha)$ è invertibile il metodo converge in maniera superlineare (più veloce di lineare, almeno quadratica).
\end{nota}

\begin{metodo}
    [Metodo di Newton semplificato]
    Si calcola il Jacobiano sempre nel punto iniziale $x_0$
    \[
    x^{(k+1)} = x^{(k)} - (Jf(x^{(0)}))^{-1}f(x^{(k)})
    \]
    in questo modo no bisogna aggiornare il Jacobiano ad ogni passo e si può precalcolare la fattorizzazione LU con pivoting di $Jf(x^{(0)})=\pi LU$ e poi all'iterazione generica risolvere il sistema lineare sfruttando la fattorizzazione...%aggiungere 
    , questo porta il costo dell'operazione ad essere di $O(n^2)$ + il costo di valutare $f(x^{(k)})$.
\end{metodo}
\begin{metodo}
    [Metodo di Jacobi non lineare (Jacobi-Newton)]
    Si sostituisce il Jacobiano con una matrice diagonale composta dalla diagonale del Jacobiano.
    \[
    x_i^{(n+1)} = x_i^{(n)} - \frac{f_i(x_1^{(n)},...,x_m^{(n)})}{\frac{\partial }{\partial x_i}f_i(x_1^{(n)},...,x_m^{(n)})} 
    \quad i=0,...,m
    \]
    costa $O(m)$ funzioni non lineari per calcolare $D_J$ ed il passo costa $O(m)$ perché richiede di risolvere un sistema diagonale.
\end{metodo}
\begin{metodo}
    [Metodo di Gauss-seidel non lineare]
    Agisce come Jacobi per l'aggiornamento di $x_1^{(n+1)}$ ma dopo utilizza le componenti di $x^{(n+1)}$ già calcolate
    \[
    x_i^{(n+1)} = x_i^{(n)} - \frac{f_i(x_1^{(n+1)},...,x_{i-1}^{(n+1)},x_i^{(n)},...,x_m^{(n)})}{\frac{\partial }{\partial x_i}f_i(x_1^{(n)},...,x_{i-1}^{(n+1)},x_i^{(n)},...,x_m^{(n)})} 
    \quad i=0,...,m
    \]
\end{metodo}
\subsection{Metodi Quasi-Newton}
Sostituiscono il Jacobiano con qualcosa che richiede meno valutazioni di funzioni non lineari e/o con cui è più facile risolvere sistemi lineari.

\chapter{Metodi per il calcolo di autovalori e autovettori}

\section{Intro}
Vogliamo trovare le autocoppie $(\lambda, v)$ tali che $Av=\lambda v$, in particolare vediamo 2 metodi che approcciano 2 versioni di  questo problema:
\begin{itemize}
    \item calcolo di una specifica autocoppia, ad esempio quella associata all'autovalore di modulo massimo.
    \item calcolo simultaneo di tutte le autocoppie.
\end{itemize}
\section{Metodo delle potenze}
Metodo iterativo pensato per trovare l'autovalore di modulo massimo. 
Questo può essere utile in casi come:
\begin{itemize}
    \item trovare il raggio spettrale di $A$
    \item trovare la $||A||_2= \rho(A^TA)$
\end{itemize}
Si può poi adattare il metodo per il calcolo di altre autocoppie.

\begin{metodo}
    [Metodo delle potenze]
Supponiamo $A\in \mathbb{C}^{n\times n}$ \textbf{diagonalizzabile} con un autovalori:
\[
|\lambda_1| > |\lambda_2| \geq |\lambda_3| \geq \cdots \geq |\lambda_n|
\]
e chiamiamo $v^{(1)}, \dots, v^{(n)}$ i relativi autovettori associati.\\
Sfruttiamo ora la seguente idea: prendiamo un certo $z\in\mathbb{C}^n$ e, dato che A è diagonalizzabile e dato che l'insieme degli autovettori forma un base di $\mathbb{C}^n$, allora vale che $z = c_1v^{(1)}+\dots+c_nv^{(n)}$ per certi scalari $c_j\in\mathbb{C}$.\\
A questo punto se applichiamo $A$ a $z$ si ottiene la seguente successione, che è anche quella che definisce il metodo:
\[
A \cdot z = A \cdot \sum_{j=1}^{n} c_j v^{(j)} = \sum_{j=1}^{n} c_j A v^{(j)} = \sum_{j=1}^{n} c_j \lambda_j v^{(j)}
\]
Se continuiamo e calcoliamo:
\[
A^2 z = A(Az) = A \sum_{j=1}^{n} c_j \lambda_j v^{(j)} = \sum_{j=1}^{n} c_j \lambda_j A v^{(j)} = \sum_{j=1}^{n} c_j \lambda_j^2 v^{(j)}
\]
Per k volte:
\[
A^k z = \sum_{j=1}^{n} c_j \lambda_j^k v^{(j)} = c_1 \lambda_1^k v^{(1)} + \ldots + c_n \lambda_n^k v^{(n)} = \lambda_1^k \left( c_1 v^{(1)} + c_2 \underbrace{\left(\frac{\lambda_2}{\lambda_1}\right)^k}_{=0} v^{(2)} + \ldots + c_n \underbrace{\left(\frac{\lambda_n}{\lambda_1}\right)^k}_{=0} v^{(n)} \right)
\]
Quei termini tendono a 0 per $k \to \infty$, quindi alla fine otteniamo:
\[
\boxed{A^k z \approx \lambda_1^k c_1 v^{(1)}}
\]
ovvero un vettore con la stessa direzione di $v^{(1)}$.
\begin{nota}
    [Perché funziona il Metodo delle potenze - schema riassuntivo]
    Il ragionamento alla base del metodo si articola in 3 passaggi concatenati:
    \begin{enumerate}
        \item \textbf{Scomposizione di $z^{(0)}$ nella base degli autovettori.} $A$ diagonalizzabile $\iff \exists S$ invertibile con $S^{-1}AS=D$: per similitudine (\autoref{def: matrici simili}) le colonne di $S$ sono gli autovettori $v^{(1)},\dots,v^{(n)}$ di $A$ e formano una \textbf{base} di $\mathbb{C}^n$. Ogni vettore di partenza (anche arbitrario) si scrive quindi in modo unico come
        \[
        z^{(0)} = c_1v^{(1)}+\dots+c_nv^{(n)}.
        \]
        \item \textbf{Applicazione ripetuta di $A$.} Per linearità di $A$ e per definizione di autovettore ($Av^{(j)}=\lambda_jv^{(j)}$), ogni componente si riscala \textit{indipendentemente} dalle altre:
        \[
        A^kz^{(0)} = \sum_{j=1}^n c_j\lambda_j^k v^{(j)} = \lambda_1^k\left(c_1v^{(1)} + \sum_{j=2}^n c_j\Big(\frac{\lambda_j}{\lambda_1}\Big)^k v^{(j)}\right).
        \]
        \item \textbf{Allineamento a $v^{(1)}$.} Poiché $|\lambda_1|>|\lambda_j|\ \forall j\geq 2$, si ha $\left(\frac{\lambda_j}{\lambda_1}\right)^k \to 0$ per $k\to\infty$: tutte le componenti non dominanti si annullano esponenzialmente, e $A^kz^{(0)}$ si allinea sempre più alla direzione di $v^{(1)}$.
    \end{enumerate}
    La logica algebrica che fa funzionare questo moltiplicare a sinistra ripetutamente è: in una matrice diagonalizzabile gli autovettori formano una base per tutto lo spazio; quindi il vettore casuale di partenza è in realtà una combinazione degli autovettori della matrice, presentando nelle sue componenti una parte per ogni autovettore, pesata in base alla rilevanza dello stesso.
    Ogni volta che quindi si applica la trasformazione A, per la definizione stessa di autocoppia ($Av=\lambda v$) si sta moltiplicando ogni componente per il suo valore.
    Farlo $k$ volte significa quindi elevare alla $k$, in particolare la componente dell'autovalore dominante viene moltiplicata per $\lambda_1^k$, le altre per il loro lambda alla k.
    Siccome $\lambda_1$ è più grande in modulo la sua potenza cresce in modo dominante rispetto alle altre e quindi il vettore di partenza si allinea sempre di più verso l'autovettore dominante.
\end{nota}
\paragraph{Riassunto dei passaggi dell'iterazione del metodo delle potenze}
I passaggi possono essere riassunti come:
\[
\boxed{
\begin{aligned}
y^{(k)} &= A \, z^{(k-1)} \\
\lambda^{(k)} &= \bigl[ z^{(k-1)} \bigr]^H y^{(k)} \\
z^{(k)} &= \frac{y^{(k)}}{\|y^{(k)}\|_2}
\end{aligned}
}
\]
dove l'ultimo passaggio serve ad evitare problemi di under/over-flow.
\end{metodo}
\begin{nota}
    Per capire il secondo passaggio invece prendiamo in esame l'ultima iterazione del metodo: per quanto visto prima siamo giunti a una perfetta corrispondenza tra $z^{(k-1)}$ e il cercato $v^{(1)}$:
    \[
    z^{k-1}\approx v^{(1)}
    \]
    Sostituendo dunque $z$ con l'autovettore possiamo riscrivere $y^{(k)}$, seguendo la definizione di autovalore e autovettore come $y^{(k)} = A v^{(1)}=\lambda_1v^{(1)}$.\\
    Dato che quello che ci interessa è il valore di $\lambda_1$ e dato che la base è unitaria, allora ci basta moltiplicare il prodotto $\lambda_1v^{(1)}$ per $[v^{(1)}]^H$ per ottenere l'autovalore:
    \[
    y^{(k)} = \lambda_1\underbrace{[v^{(1)}]^Hv^{(1)}}_{=\|v^{(1)}\|_2^2=1} = \lambda_1
    \]
    infatti, per via del fatto che $A$ è hermitiana, la base che i suoi autovettori formano è \uline{unitaria} e dunque la norma 2 di qualunque di questi autovettori risulta pari a 1.
\end{nota}
\begin{nota}
    Il prossimo teorema è alla base del metodo delle potenze, si spiega quindi anche nel dettaglio il significato della successione prima definita.
\end{nota}
\begin{teorema}
    [Convergenza del Metodo delle Potenze]
    Se A è Hermitiana (condizione rilassabile: si veda \ref{sec: ipotesi metodo potenze}) con un autovalori 
    \[|\lambda_1| > |\lambda_2| \geq \cdots \geq |\lambda_n|\]
     e con $v^{(1)}$ autovettore associato a $\lambda_1$ e $h\in\{1, \dots , n\}$ tali che $v^{(1)}_h \ne 0$ e dato un punto di partenza $z^{(0)}\in \mathbb{C}^n$ tale che $\left[ v^{(1)} \right]^H z^{(0)}\ne 0$ (ovvero il punto $z^{(0)}$ non è ortogonale a $v^{(1)}$) allora la successione:
    \[
    \begin{cases}
        z^{(0)} = dato \\
        z^{(k)} = Az^{(k-1)}
    \end{cases}
    \]
    è tale che 
    \[
    \lim_{k\to\infty} \frac{z^{(k)}}{z^{(k)}_h} = \frac{v^{(1)}}{v_h^{(1)}} = \tilde{v}^{(1)}
    \]
    dove $\tilde{v}^{(1)}$ è un multiplo scalare di $v^{(1)}$ in cui la componente $h$ è uguale a 1.\\
    \begin{nota}
        Come illustrato sotto nella dimostrazione, dividere per la componente h-esima il vettore ottenuto $k$ serve a rimuovere il fattore esponenziale incontrollato $\lambda_1^k$.
    \end{nota}
    \noindent
    La successione \textbf{quindi converge} all'autovettore mentre il limite:
    \[
    \lim_{k\to\infty} \frac{[z^{(k)}]^HAz^{(k)}}{[z^{(k)}]^Hz^{(k)}}=\lambda_1
    \]
    è uguale all'autovalore.
\end{teorema}
\begin{nota}
    La convergenza quindi si trova se uno dei passi (numericamente finito) della successione presenta componente non nulla lungo l'autovettore dominante.
\end{nota}

\begin{definizione}
    [Quoziente di Rayleigh]
    È così detta la quantità:
    \[
    r(x) = \frac{x^HAx}{x^Hx}
    \]
    presente nel limite appena utilizzato sopra.\\
    Non fa altro che estrarre il valore numerico del lambda dominante.\\
\end{definizione}

\begin{dimostrazione}
    Poiché $A$ è Hermitiana, per il teorema spettrale (\autoref{th: teorema spettrale}) è diagonalizzabile e posso scegliere i vettori $v^{(1)}, \ldots, v^{(n)}$ in modo che
    \[
        V = \bigl[\,v^{(1)} \mid \cdots \mid v^{(n)}\,\bigr]
    \]
    sia \textbf{unitaria}, ovvero $V^H V = V V^H = I$. In particolare:
    \[
        \bigl[v^{(i)}\bigr]^{\!H} v^{(j)} =
        \begin{cases}
            0 & i \neq j, \\
            1 & i = j.
        \end{cases}
    \]
    Poiché $\{v^{(1)}, \ldots, v^{(n)}\}$ è una base ortonormale di $\mathbb{C}^n$,
    possiamo scrivere
    \[
        z^{(0)} = \sum_{j=1}^{n} c_j\, v^{(j)}.
    \]
    Il coefficiente $i$-esimo si ricava moltiplicando a sinistra per
    $\bigl[v^{(i)}\bigr]^{\!H}$:
    \[
        \bigl[v^{(i)}\bigr]^{\!H} z^{(0)}
        = \bigl[v^{(i)}\bigr]^{\!H} \sum_{j=1}^{n} c_j\, v^{(j)}
        = \sum_{j=1}^{n} c_j
        \underbrace{\bigl[v^{(i)}\bigr]^{\!H} v^{(j)}}_{\substack{=\,0 \; \text{se } i\neq j\\=\,1 \; \text{se } i=j}}
        = c_i.
    \]
    In particolare, l'ipotesi $\bigl[v^{(1)}\bigr]^{\!H} z^{(0)} \neq 0$ implica $c_1 \neq 0$ (questo ci serve perché garantisce che $z^{(0)}$ non ha componente nulla lungo la direzione di $v^{(1)}$), basta pensare all riscrittura di $z$ come combinazione lineare della base composta da $V$).
    \\
    Come già discusso in precedenza, iterando si ottiene:
    \[
        z^{(k)} = A\,z^{(k-1)} \to A^k\,z^{(0)}
        = \sum_{j=1}^{n} c_j\,\lambda_j^k\,v^{(j)}
        = \lambda_1^k \left(
            c_1\,v^{(1)}
            + \sum_{j=2}^{n} c_j \left(\frac{\lambda_j}{\lambda_1}\right)^{\!k} v^{(j)}
        \right).
    \]
    \medskip
    \textbf{Convergenza all'autovettore.}
    Sia $h \in \{1,\ldots,n\}$ tale che $v_h^{(1)} \neq 0$.
    Dividiamo $z^{(k)}$ per la sua componente $h$-esima $z_h^{(k)}$:
    \[
        \frac{z^{(k)}}{z_h^{(k)}}
        =
        \frac{
            \lambda_1^k\!\left(
                c_1\,v^{(1)}
                + \displaystyle\sum_{j=2}^{n} c_j
                \!\left(\tfrac{\lambda_j}{\lambda_1}\right)^{\!k}\! v^{(j)}
            \right)
        }{
            \lambda_1^k\!\left(
                c_1\,v_h^{(1)}
                + \displaystyle\sum_{j=2}^{n} c_j
                \!\left(\tfrac{\lambda_j}{\lambda_1}\right)^{\!k}\! v_h^{(j)}
            \right)
        }.
    \]
    Poiché $\left|\dfrac{\lambda_j}{\lambda_1}\right| < 1$ per $j \geq 2$,
    il termine $\displaystyle\sum_{j=2}^{n}(\cdots) \to 0$
    (come già dimostrato in precedenza), quindi:
    \[
        \lim_{k \to +\infty} \frac{z^{(k)}}{z_h^{(k)}}
        = \lim_{k \to +\infty} \frac{c_1\,v^{(1)}}{c_1\,v_h^{(1)}}
        = \frac{v^{(1)}}{v_h^{(1)}}
    \]
    che è un multiplo scalare dell'autovettore $v^{(1)}$
    (la componente $h$-esima è normalizzata a $1$).
    \medskip
    \textbf{Convergenza all'autovalore tramite quoziente di Rayleigh.}
    Il quoziente di Rayleigh è invariante per riscalamento, dunque:
    \[
        \lim_{k \to +\infty}
        \frac{\bigl[z^{(k)}\bigr]^{\!H} A\, z^{(k)}}
            {\bigl[z^{(k)}\bigr]^{\!H} z^{(k)}}
        = \lim_{k \to +\infty}
        \frac{
            \left[\dfrac{z^{(k)}}{z_h^{(k)}}\right]^{\!\!H}
            A\,
            \dfrac{z^{(k)}}{z_h^{(k)}}
        }{
            \left[\dfrac{z^{(k)}}{z_h^{(k)}}\right]^{\!\!H}
            \dfrac{z^{(k)}}{z_h^{(k)}}
        }
        =
        \frac{
            \left[\dfrac{v^{(1)}}{v_h^{(1)}}\right]^{\!\!H}
            A\,
            \dfrac{v^{(1)}}{v_h^{(1)}}
        }{
            \left[\dfrac{v^{(1)}}{v_h^{(1)}}\right]^{\!\!H}
            \dfrac{v^{(1)}}{v_h^{(1)}}
        }.
    \]
    Il fattore $1/v_h^{(1)}$ si cancella al numeratore e al denominatore:
    \[
        = \frac{\bigl[v^{(1)}\bigr]^{\!H} A\, v^{(1)}}
            {\bigl[v^{(1)}\bigr]^{\!H} v^{(1)}}
        = \mathcal{R}\!\left(v^{(1)}\right)
        = \lambda_1
    \]
\end{dimostrazione}
\paragraph{Osservazione}
La successione $z^{(k)} = A^Kz^{(k-1)}$ può soffrire facilmente di under/over-flow.\\
Dato che moltiplicare o dividere per costanti non cambia la proprietà di essere o meno un autovettore di $A$ (conta solo la direzione), il metodo delle potenze prende la successione e rende i vettori di modulo 1 ad ogni passo (divisione per la norma del vettore)

\subsection{Metodo delle potenze in codice}
\begin{lstlisting}
    z0 dato in input
    for k=1,2,...,MAXIT
        yk =Az(k-1)
        zk = yk / ||yk||_2
    end
\end{lstlisting}
Il costo è quindi dato dal costo del prodotto matrice-vettore + $O(n)$ per il calcolo della norma. Tutto moltiplicato per il numero di iterazioni.\\
Si ha quindi una complessità totale di $O(k(n^2+n))$ con $K$ iterazioni.
\subsection{Criterio d'arresto}
\begin{definizione}
    [Criterio d'arresto del Metodo delle potenze]
    \[|r(z^{(k)}) - r(z^{(k-1)})| < \text{tol}\]
\end{definizione}

\subsection{Commenti sulle ipotesi}
\label{sec: ipotesi metodo potenze}
\begin{enumerate}
    \item La condizione $A=A^H$, che comportava $A$ diagonalizzabile, si può rilassare: \textbf{è} infatti \textbf{dimostrabile} il risultato di \uline{\textbf{convergenza per sola $A$ diagonalizzabile}}.
    \item L'assunzione $|v^{(1)}|^T z^{(0)} \neq 0$ può sembrare difficile da garantire a priori. Però se si genera $z^{(0)}$ in maniera casuale vale con probabilità 1.
    \item L'assunzione $|\lambda_1| > |\lambda_2|$ non è rimuovibile, ad esempio se $A = \begin{bmatrix} 0 & 1 \\ 1 & 0 \end{bmatrix}$ $(\lambda_1 = 1, \lambda_2 = -1)$ e prendiamo $z^{(0)} \neq \frac{1}{1}$, si finisce in un loop senza convergenza:
    \[
    z^{(0)} = \frac{2}{3}, \, Az^{(0)} = \frac{3}{2}, \, A^2 z^{(0)} = \frac{2}{3}, \, A^3 z^{(0)} = \frac{3}{2} \ldots
    \]
\end{enumerate}
\subsection{Varianti del metodo delle potenze}
In questa variante supponiamo invece che un autovalore sia, in modulo, strettamente minore degli altri, per applicare poi il metodo delle potenze ad $A^{-1}$ che ha gli autovalori, guardandoli come moduli dei reciproci di essi, con uno strettamente maggiore degli altri.
Si parla quindi di:
\begin{metodo}
    [Metodo delle potenze inverse]
    Applicando il metodo delle potenze alla matrice $A^{-1}$,
    si pu\`o calcolare l'autovalore minimo $\lambda_n$ della matrice $A$.
    Supponiamo che $A$ sia non singolare, diagonalizzabile e con autovalori
    \[
    |\lambda_1| \ge |\lambda_2| \ge \cdots \ge |\lambda_{n-1}| > |\lambda_n| > 0.
    \]
    Come sappiamo, la matrice $A^{-1}$ avr\`a autovalori $\frac{1}{\lambda_i}$
    tali che
    \[
    \left| \frac{1}{\lambda_n} \right|
    >
    \left| \frac{1}{\lambda_{n-1}} \right|
    \ge \cdots \ge
    \left| \frac{1}{\lambda_1} \right|
    \]
    Inoltre $A$ e $A^{-1}$ hanno gli stessi autovettori (proprietà \ref{item: proprieta autovalori inversa}).\\
    Poich\'e in generale $A^{-1}$ non \`e nota
    (e comunque sconveniente da calcolare),
    il vettore $y^{(k)}$ nel metodo delle potenze per $A^{-1}$
    si trover\`a risolvendo il sistema lineare
    \[
    A y^{(k)} = z^{(k-1)}
    \]
    Tali sistemi si risolvono efficientemente utilizzando la fattorizzazione
    LU di $A$ (eventualmente con pivoting),
    calcolata una volta per tutte,
    prima di iniziare a calcolare la successione $\{z^{(n)}\}$.\\
    Vediamo dunque come si può tradurre in codice questo metodo:
    \begin{lstlisting}
function [lam, zvec] = potenze_inverse(A, z0, maxit)
    n = length(z0);
    zvec = zeros(n, maxit);
    [L, U] = lu(A);
    for j = 1:maxit
        y = U \ (L \ z0);
        lam = z0' * y;
        z = y / norm(y);
        zvec(:, j) = z;
        z0 = z;
    end
    lam = 1 / lam;
end
    \end{lstlisting}
\end{metodo}

\noindent
Vediamo una ulteriore variante:
\begin{metodo}
    [Metodo delle potenze shiftato]
Pseudocodice del metodo:
    \begin{lstlisting}[mathescape=true]
$z^{(0)}$ dato
Calcolo di $(A - \sigma I) = \pi LU$
for $k = 1, 2, 3, \dots, \hat{k}$
    risolvo $\pi LU v^{(k)} = z^{(k-1)}$
    $z^{(k)} = v^{(k)} / \|v^{(k)}\|_2$
end
$\mu = [z^{(k)}]^H A z^{(k)}$
return $\underbrace{\frac{1}{\mu}}_{\text{stima di } \lambda_j} + \sigma$, $z^{(k)}$
    \end{lstlisting}
\end{metodo}
    
\begin{lezione}
    [Lezione 23 20/05]
\end{lezione}

\section{Metodo QR per gli autovalori}
\begin{nota}
    (dovrebbe essere più roba da orale???)
\end{nota}
\noindent
Idea di base: se applico una trasformazione di similitudine a una matrice gli autovalori non cambiano e gli autovettori cambiano in una maniera che conosco.

\begin{nota}
    [Trasformazioni per similitudine e autovalori/autovettori]
    Data $A \in \mathbb{C}^{n\times n}$, $S \in \mathbb{C}^{n\times n}$ invertibile, $B = S^{-1}AS$.\\
    $B$ ha gli stessi autovalori di $A$ ed autovettori $S^{-1}v$, dove $v$ è autovettore di $A$.\\
    Se conosco autovalori e autovettori di $B$, e conosco $S$, allora posso ricavare facilmente autovalori e autovettori di $A$.
\end{nota}

L'idea è quindi costruire una successione di trasformazioni per similitudine che porti ad una matrice $T$ di cui sia facile calcolare autovalori e autovettori. Nel caso del metodo QR si cerca $T$ triangolare:
\[
A_1 \xrightarrow{S_1^{-1}A_1S_1} A_2 \xrightarrow{S_2^{-1}A_2S_2} A_3 \longrightarrow \cdots \longrightarrow T
\]

\paragraph{1° problema}
Le trasformazioni per similitudine coinvolgono delle inverse, ovvero la risoluzione di sistemi lineari: si presentano quindi problemi di condizionamento se $S_j$ non ha particolari proprietà.

\paragraph{Soluzione}
Si formano le trasformazioni per similitudine con matrici \textbf{unitarie} (per una matrice unitaria $Q^{-1}=Q^H$, quindi l'inversa è immediata da calcolare e non introduce errore):
\[
A \xrightarrow{Q_1^HAQ_1} A_1 \xrightarrow{Q_2^HA_1Q_2} A_2 \longrightarrow \cdots \longrightarrow T
\]

\subsection{Come è definita la successione Ak}
\begin{definizione}
    [Successione base del metodo QR]
    \[
    \begin{cases}
        A_1 = A\\
        \text{per } k=1,2,3,\dots\\
        \quad \text{calcolo } A_k = Q_kR_k \quad \text{(fattorizzazione QR di $A_k$)}\\
        \quad \text{calcolo } A_{k+1} = R_kQ_k
    \end{cases}
    \]
\end{definizione}
\paragraph{Perché è una trasformazione per similitudine}
    \[
    A_{k+1} = R_kQ_k = Q_k^H\underbrace{Q_kR_k}_{A_k}Q_k = Q_k^HA_kQ_k
    \]
    quindi $A_{k+1}$ è simile ad $A_k$, e per transitività $A_{k+1}$ è simile ad $A$. In particolare:
    \[
    \begin{split}
    A_{k+1} &= Q_k^HA_kQ_k = Q_k^HQ_{k-1}^HA_{k-1}Q_{k-1}Q_k = \cdots = Q_k^H\cdots Q_1^H\,A\,Q_1\cdots Q_k \\
    &= \tilde{Q}^HA\tilde{Q} \qquad \textit{con }\tilde{Q} = Q_1\cdots Q_k
    \end{split}
    \]

\begin{teorema}
    [Convergenza del metodo QR]
    Sia $A\in\mathbb{R}^{n\times n}$ e supponiamo $|\lambda_1|>|\lambda_2|>\cdots>|\lambda_n|$. Allora
    \[
    \lim_{k\to+\infty}A_k = T = \begin{bmatrix} t_{11} & \cdots & \cdots\\ & \ddots & \vdots \\ 0 & & t_{nn}\end{bmatrix} 
    \]
    Se gli autovalori di $A$ sono distinti ma possibilmente con lo stesso modulo (es. coppie complesse coniugate), ovvero \uline{se $A$ ha autovalori complessi semplici}, allora
    \[
    \lim_{k\to+\infty}A_k = T = \begin{bmatrix} T_{11} & & \\ & \ddots & \\ & & T_{ss}\end{bmatrix} \quad \text{(triangolare superiore a blocchi)}
    \]
    dove ogni blocco diagonale $T_{ii}$ è $1\times1$ oppure $2\times2$.
\end{teorema}
\begin{nota}
    Nel secondo caso, gli autovalori di $T$ (e quindi di $A$) sono dati dai blocchi $1\times1$ sulla diagonale, oppure dagli autovalori dei blocchi diagonali $2\times2$ (complessi coniugati).
    % NOTA DI TRASCRIZIONE: l'esempio numerico riportato negli appunti (un blocco 2x2 con voci del tipo ±√2/2, oltre a un blocco 1x1 pari a 3) non è riprodotto qui con certezza: la struttura suggerisce un blocco di tipo "rotazione" i cui autovalori si calcolano risolvendo il polinomio caratteristico 2x2. Verificare sul quaderno per i valori esatti.
\end{nota}

Il metodo QR base consiste dunque nel generare la successione $\{A_{k+1}\}=\{R_kQ_k\}$ e fermarsi quando gli elementi sotto la diagonale di $A_k$ sono sufficientemente vicini a $0$. A quel punto si calcolano autovalori e autovettori di $T$, parte triangolare superiore di $A_k$:
\[
\setlength{\arraycolsep}{2pt}
\renewcommand{\arraystretch}{1.3}
A_k = \begin{bmatrix}
    * & \cdots & * \\
    & \ddots &\vdots \\
    \varepsilon & & *
\end{bmatrix}
 \xrightarrow{\varepsilon\to 0} T =
\begin{bmatrix}
    * & \cdots & * \\
    & \ddots &\vdots \\
     & & *
\end{bmatrix}
\quad \longrightarrow \quad \text{calcolo } (\lambda_T, v_T)
\]
Si usa $\lambda_T$ come approssimazione degli autovalori di $A$, e $\tilde{Q}v_T$ (con $\tilde{Q}=Q_1\cdots Q_k$) come approssimazione degli autovettori di $A$.

\subsection{Come si calcolano gli autovettori di una matrice triangolare con autovalori distinti}
\begin{nota}
    In pratica questo schema non è direttamente utilizzabile (è troppo lento); esistono accorgimenti per renderlo efficiente (si veda il Recap più avanti).
\end{nota}
Sia $T$ triangolare superiore con autovalori distinti sulla diagonale, $t_{ii}\neq t_{jj}\ \forall i\neq j$. Se prendo $\lambda_j = t_{jj}$, so che l'autovettore associato è un vettore $v_j$ che risolve $(T-\lambda_jI)v_j = 0$. 
\begin{figure}[H]
    \centering
    \includegraphics[width=0.6\linewidth]{img/T in qr per autoval.jpg}
\end{figure}
\noindent
dove $T_1, T_2$ sono ancora triangolari superiori. 
Dobbiamo quindi risolvere un sistema della forma:
\[
\begin{bmatrix} T_1 & z & * \\ 0 & 0 & * \\ 0 & 0 & T_2 \end{bmatrix}\begin{bmatrix} x_1\\x_2\\x_3\end{bmatrix} = \begin{bmatrix}0\\0\\0\end{bmatrix}
\]
Potrei prevedere $x_2=0$.
Scegliamo poi una soluzione $\gamma=1$ (normalizzazione). 
Questo perché rimane da risolvere:
\[
T_1x_1 + \gamma z = 0 \implies T_1x_1 = -\gamma z
\]
sistema triangolare superiore, da cui $x_1 = -T_1^{-1} z$. 
In conclusione:
\[
\renewcommand{\arraystretch}{1}
x^{(j)} = \begin{bmatrix}-T_1^{-1} z\\ 1\\ 0\\ \vdots \\ 0\end{bmatrix}
\]
\begin{nota}
    [Costo]
    Calcolare $x^{(j)}$ richiede di risolvere un sistema triangolare superiore $(j-1)\times(j-1)$, quindi costa $O((j-1)^2)=O(j^2)$. Per trovare tutti gli autovettori:
    \[
    O\left(\sum_{j=1}^n j^2\right) = O(n^3)
    \]
\end{nota}

\paragraph{Problema}
Il metodo QR richiede di calcolare fattorizzazioni QR e moltiplicare matrici ad ogni passo anche assumento $O(n)$ iterazioni per arrivare a conergenza ( $O(1)$ iterazioni per autovalore) questo ci darebbe costo $O(n^4)$.

\paragraph{Soluzione}
Si deve fare \textit{pre-processing} su $A$.

\paragraph{Fatto funzionale alla risoluzione}
Se $A$ è in forma di Hessemberg (ha elementi $= 0$ solo da sotto la prima sottodiagonale) allora posso fare una iterazione/passo del metodo QR con costo $O(n^)$.\\
Questo perché:
\begin{enumerate}
    \item calcolare la fattorizzazione QR di una matrice in forma di Hessenberg costa $O(n^2)$;
    \item calcolare $R_kQ_k$ costa $O(n^2)$, e il risultato resta in forma di Hessenberg.
\end{enumerate}

\subsection{Come si sfrutta la forma di Hessenberg}
Ogni matrice $A$ è simile a una matrice $H$ in forma di Hessenberg: $\exists\,W\in\mathbb{C}^{n\times n}$ unitaria tale che 
$$A = W\cdot H\cdot W^H$$

Per trovare $H$ (e $W$) si usa un algoritmo simile a quello della fattorizzazione QR, con trasformazioni di Householder applicate però \textbf{per similitudine} (a destra e a sinistra).

\textbf{Passo 1.} Sia $a_1\in\mathbb{C}^{n-1}$ la prima colonna di $A$ sotto l'elemento diagonale.\\
Si prende $\tilde{M}_1\in\mathbb{C}^{(n-1)\times(n-1)}$ di Householder tale che 
$
\renewcommand{\arraystretch}{1.1}
\tilde{M}_1a_1 = \begin{bmatrix}*\\0\\ \vdots\\0\end{bmatrix}
$
, e si pone
\[
M_1 = \begin{bmatrix}1 & 0\\ 0 & \tilde{M}_1\end{bmatrix}\in\mathbb{C}^{n\times n}
\]
Si calcola $M_1A$ (prima colonna azzerata sotto la sottodiagonale) e poi, per similitudine, $M_1AM_1^H$.

\textbf{Passo 2.} Si ripete sulla colonna successiva $a_2\in\mathbb{C}^{n-2}$: si costruisce $\tilde{M}_2\in\mathbb{C}^{(n-2)\times(n-2)}$ di Householder con 
$
\renewcommand{\arraystretch}{1.1}
\tilde{M}_2a_2=\begin{bmatrix}*\\0\\\vdots\\0\end{bmatrix}
$
, si pone $M_2=\begin{bmatrix}I_2 & 0\\0 & \tilde{M}_2\end{bmatrix}$ e si calcola $M_2M_1AM_1^HM_2^H$.\\
Dopo $n-2$ passi si ottiene
\[
M_{n-2}\cdots M_1\,A\,M_1^H\cdots M_{n-2}^H = H
\]
in forma di Hessenberg, con $W = M_1^H\cdots M_{n-2}^H$. Il costo totale è $O(n^3)$ e va calcolato \textbf{una sola volta}, prima di iniziare le iterazioni del metodo QR.

\subsection{Recap Metodo QR}
\begin{enumerate}
    \item Calcolo la forma di Hessenberg di $A$: $A = WHW^H$ (costo $O(n^3)$, effettuato una sola volta)
    \item Genero i termini della successione:
    \[
    \begin{cases}
        A_0 = H\\
        A_k = Q_kR_k \quad \text{(costo $O(n^2)$)}\\
        A_{k+1} = R_kQ_k \quad \text{(costo $O(n^2)$)}
    \end{cases}
    \quad\Longrightarrow\quad \text{costo totale } O(\hat{k}n^2)
    \]
    \item Prendo $T$, parte triangolare superiore di $A_{\hat{k}}$
    \item Calcolo autovalori e autovettori di $T$ (costo $O(n^3)$)
    \item Restituisco $\lambda_T$ come approssimazione degli autovalori di $A$, e le quantità $W\tilde{Q}v_T$ come approssimazione degli autovettori di $A$ (dove $\tilde{Q}=Q_1\cdots Q_{\hat{k}}$ è tale che $A_{\hat{k}} = \tilde{Q}^HH\tilde{Q}$) (costo $O(n^3)$)
\end{enumerate}

\section{Ulteriori criticità}
Non le affrontiamo.

\chapter{Esercizi}

\begin{lezione}
    [Lezione 24 21/05]
\end{lezione}
 
\section{Fattorizzazione LU e risoluzione sistema lineare}
\section{Approssimazione di funzione nel senso dei minimi quadrati}
\section{Equazione non lineare}
Discussione radici, intervalli di separazione, determinazione della convergenza locale di un metodo date le varie radici.
\section{Analisi dell'errore algoritmico delle formule}
\section{Esercizi di algebra}
\subsection{Determinare per quali alpha di A il metodo di Gauss-Seidel converge}
\subsection{Esercizio 2 pag. 112 Prove scritte 2022-2023}
Calcolo della fattorizzazione spettrale di A.
\subsection{Esercizio 2 pag. 4 Prove scritte 2022-2023}
Dimostrare che A è riducibile, trovare matrice di permutazione che porti A in forma triangolare superiore a blocchi.

\chapter{Laboratorio di Calcolo Numerico}

% ============================================================
\section{Lezione 1 -- Introduzione a MATLAB}
% ============================================================

\subsection{Installazione}

Per svolgere gli esercizi è necessario disporre di MATLAB o di un'alternativa
compatibile:
\begin{itemize}
  \item \textbf{MATLAB}: disponibile tramite licenza Campus dell'Ateneo
        (\texttt{mathworks.com}) oppure nella versione online
        (\texttt{matlab.mathworks.com}).
  \item \textbf{GNU Octave}: alternativa open source, gratuita e compatibile
        per tutti gli esercizi del laboratorio (\texttt{gnu.org/software/octave}).
\end{itemize}

\subsection{Calcoli in virgola mobile}

MATLAB utilizza la \textit{doppia precisione} (8 byte per numero).
I comandi principali per esplorare i limiti della rappresentazione sono:

\begin{lstlisting}[language=Matlab]
realmin   % minimo numero positivo: 2.2251e-308
realmax   % massimo numero positivo: 1.7977e+308
eps       % precisione di macchina:  2.2204e-16
\end{lstlisting}

Il punto e virgola alla fine di una riga sopprime la stampa del risultato.
Il comando \texttt{format long} aumenta il numero di cifre visualizzate.

\subsection{Variabili}

In MATLAB non è necessario dichiarare né specificare il tipo di una variabile;
basta procedere direttamente con l'assegnamento. Le convenzioni sui nomi sono:
\begin{itemize}
  \item devono iniziare con una lettera;
  \item possono contenere lettere, numeri e underscore;
  \item lunghezza massima: 63 caratteri;
  \item il linguaggio è \textit{case sensitive} (\texttt{pAlla} $\neq$ \texttt{palla}).
\end{itemize}

\subsection{Script e Funzioni}

Uno \textbf{script} è una sequenza di comandi salvata in un file \texttt{.m}.
Per eseguirlo è sufficiente digitarne il nome nel prompt (senza estensione).\\
Una \textbf{funzione} segue la sintassi:
\begin{lstlisting}[language=Matlab]
function [y1,...,yN] = miafunzione(x1,...,xM)
    % corpo della funzione
end
\end{lstlisting}
La dichiarazione \texttt{function} deve essere la prima istruzione eseguibile
del file. La parola chiave \texttt{return} interrompe l'esecuzione anticipatamente.
Esempio:
\begin{lstlisting}[language=Matlab]
function [area] = areacerchio(raggio)
    if raggio < 0
        area = -1;
        return
    end
    area = pi * raggio^2;
end
\end{lstlisting}
Per consultare la documentazione di una funzione:
\begin{lstlisting}[language=Matlab]
help cosd
\end{lstlisting}

\subsection{Perdita di precisione}

Non tutti i numeri reali sono rappresentabili in virgola mobile; ciò provoca
\textit{errori di arrotondamento}. Un esempio classico:
\begin{lstlisting}[language=Matlab]
e = 1 - 3*(4/3 - 1)   % ci si aspetta 0, si ottiene 2.2204e-16
\end{lstlisting}
Gli \textit{errori di cancellazione} si verificano sottraendo numeri grandi
e molto vicini tra loro. Per risolverli occorre riformulare l'algoritmo.
Ad esempio, per le radici di $ax^2 + bx + c = 0$, invece di applicare
direttamente la formula quadratica per entrambe le radici, si calcola solo
quella senza cancellazione e si ricava l'altra dalla relazione $x_1 \cdot x_2 = c/a$.

\subsection{Approssimazione dell'esponenziale}

La serie di Taylor troncata approssima $e^x$:
\[
  e^x \approx \sum_{k=0}^{n} \frac{x^k}{k!}.
\]
Un'implementazione efficiente usa un accumulatore per il termine generico,
evitando il calcolo di $k!$ e $x^k$ separatamente ($O(n)$ operazioni anziché $O(n^2)$):
\begin{lstlisting}[language=Matlab]
function a = myexp2(x, n)
    t = 1;  % termine generico
    a = 1;  % somme parziali
    for k = 1:n
        t = t * x / k;
        a = a + t;
    end
end
\end{lstlisting}
Per $x < 0$ la serie soffre di errori di cancellazione; si preferisce calcolare
$1/e^{-x}$ (usando la serie con $-x > 0$) per evitare il problema.


% ============================================================
\section{Lezione 2 -- Vettori e Matrici}
% ============================================================

\subsection{Creare vettori e matrici}

\begin{lstlisting}[language=Matlab]
A = [1 2 3; 4 5 6]   % matrice 2x3 (righe separate da ;)
zeros(3,2)            % matrice di zeri
ones(3,2)             % matrice di uni
eye(3)                % identita' 3x3
randn(2,3)            % elementi gaussiani casuali
\end{lstlisting}
Il \textbf{range operator} \texttt{a:t:b} genera un vettore riga con passo $t$:
\begin{lstlisting}[language=Matlab]
1:0.5:4    % [1.0 1.5 2.0 ... 4.0]
1:10       % [1 2 3 ... 10]
\end{lstlisting}
Il comando \texttt{linspace(a,b,n)} suddivide $[a,b]$ in $n-1$ intervalli uguali.

\subsection{Operazioni su vettori e matrici}

\begin{lstlisting}[language=Matlab]
a = 1:3; b = 4:6;
a + b          % somma elemento per elemento: [5 7 9]
sin(a)         % funzioni applicate elemento per elemento
a .* b         % prodotto elemento per elemento (operatore .)
c = a'         % trasposta (vettore colonna)
C = a' * b     % prodotto esterno (matrice 3x3)
length(a)      % lunghezza di un vettore
size(C)        % dimensioni (righe, colonne)
\end{lstlisting}
Per leggere/scrivere sottomatrici si usa l'operatore "\texttt{:}":
\begin{lstlisting}[language=Matlab]
A(1:2, 2:3)      % sottomatrice righe 1-2, colonne 2-3
A(2:end, 1:end-2)
A(1, :)          % prima riga intera
\end{lstlisting}
Scrivere un elemento fuori dai limiti espande automaticamente la matrice.

\subsection{Matrici per diagonali}

Il comando \texttt{diag(v,k)} inserisce gli elementi del vettore $v$ sulla
$k$-esima diagonale (0 = principale, $k>0$ = sopradiagonale, $k<0$ = sottodiagonale):
\begin{lstlisting}[language=Matlab]
A = diag(1:10) + diag(ones(9,1), 1)  % bidiagonale superiore
\end{lstlisting}

\subsection{Grafici}

\texttt{plot(x,y)} collega con una linea i punti $(x_i, y_i)$:
\begin{lstlisting}[language=Matlab]
r = 1:10;
plot(r, r.^2)

t = linspace(0, 2*pi, 1000);
plot(cos(t), sin(t))   % cerchio unitario
\end{lstlisting}

\subsection{Tempi di calcolo}

\texttt{tic} e \texttt{toc} misurano il tempo di esecuzione di un blocco di codice:
\begin{lstlisting}[language=Matlab]
A = randn(100); B = randn(100);
tic;
C = A * B;
t = toc;
\end{lstlisting}
Il prodotto di matrici $n \times n$ ha costo $O(n^3)$: raddoppiando $n$,
il tempo scala di un fattore $\approx 8$.
Su grafico \texttt{loglog}, la pendenza della retta del costo dovrebbe essere 3.


% ============================================================
\section{Lezione 3 -- Cerchi di Gerschgorin}
% ============================================================

\subsection{Disegnare un cerchio}

Un cerchio di centro $z \in \mathbb{C}$ e raggio $r$ si traccia parametricamente:
\begin{lstlisting}[language=Matlab]
function circle(z, r)
    t = linspace(0, 2*pi, 1000);
    x = real(z) + r * cos(t);
    y = imag(z) + r * sin(t);
    plot(x, y);
end
\end{lstlisting}

\subsection{Cerchi di Gerschgorin}

Il $k$-esimo cerchio di Gerschgorin di una matrice $A \in \mathbb{C}^{n \times n}$ è
\[
  F_k(A) = \left\{ z \in \mathbb{C} : |z - A_{kk}| \leq \sum_{j \neq k} |A_{kj}| \right\}.
\]
Ogni autovalore di $A$ appartiene all'unione di tali cerchi.\\
La funzione seguente disegna autovalori e cerchi di Gerschgorin:
\begin{lstlisting}[language=Matlab]
function disegna_gersh(A)
    n = size(A, 1);
    close all; hold on; axis('equal');
    autovalori = eig(A);
    plot(real(autovalori), imag(autovalori), '*');
    for k = 1:n
        center = A(k,k);
        radius = sum(abs(A(k,:))) - abs(A(k,k));
        circle(center, radius);
    end
end
\end{lstlisting}

\begin{figure}[H]
    \centering
    \includegraphics[width=0.75\linewidth]{img/ger matlab.png}
    \caption{Esempio di plot con questa funzione}
\end{figure}
\textbf{Nota sull'efficienza}: in MATLAB, preferire operazioni su vettori
ai cicli \texttt{for} è fondamentale per le prestazioni, poiché ogni istruzione
interpretata ha un costo fisso. Ad esempio, \texttt{sum(abs(v))} è molto più
veloce del corrispondente ciclo \texttt{for}. Questo diventa cruciale per
cicli annidati come quelli delle fattorizzazioni matriciali.


% ============================================================
\section{Lezione 4 -- Sistemi Triangolari ed Eliminazione di Gauss}
% ============================================================

\subsection{Matrici di permutazione}

Una permutazione su $\{1, \ldots, n\}$ si rappresenta con un vettore $p$ di interi.
Moltiplicare $A$ a sinistra per la matrice di permutazione $\Pi$ associata a $p$
si ottiene banalmente con l'indicizzazione:
\begin{lstlisting}[language=Matlab]
p = [2 3 4 1];
A(p, :)   % equivale a Pi * A
A(p, p)   % equivale a Pi * A * Pi^T
\end{lstlisting}
La permutazione inversa $p^{-1}$ si calcola con:
\begin{lstlisting}[language=Matlab]
ip = zeros(1, n);
ip(p) = 1:n;
\end{lstlisting}

\subsection{Soluzione di sistemi triangolari}

Un sistema triangolare inferiore $Lx = b$ si risolve per sostituzione in avanti:
\[
  x_i = \frac{b_i - \sum_{j=1}^{i-1} L_{ij} x_j}{L_{ii}}, \quad i = 1, \ldots, n.
\]
\begin{lstlisting}[language=Matlab]
function x = inf_solve(L, b)
    n = size(L, 1);
    x = zeros(n, 1);
    for i = 1:n
        p = b(i);
        for j = 1:i-1
            p = p - L(i,j) * x(j);
        end
        x(i) = p / L(i,i);
    end
end
\end{lstlisting}
Il ciclo interno si può sostituire con un'operazione vettoriale
(\texttt{L(i,1:i-1) * x(1:i-1)}) per migliorare le prestazioni.\\
I comandi \texttt{tril} e \texttt{triu} estraggono la parte triangolare inferiore e superiore di una matrice, rispettivamente.

\subsection{Eliminazione di Gauss}

Al passo $k$, si sottraggono multipli della riga $k$ alle righe $k+1, \ldots, n$
per azzerare la colonna $k$ sotto il pivot $A^{(k)}_{kk}$:
\begin{lstlisting}[language=Matlab]
function [U, c] = my_gauss(A, b)
    n = size(A, 1);
    for k = 1:n-1
        for i = k+1:n
            m = A(i,k) / A(k,k);
            A(i,:) = A(i,:) - m * A(k,:);
            b(i)   = b(i)   - m * b(k);
        end
    end
    U = A; c = b;
end
\end{lstlisting}
La riga interna può essere riscritta con operazioni su vettori eliminando
il ciclo \texttt{for} in $i$.
\begin{lstlisting}
function [U, c] = my_gauss(A, b)
    n  = size(A, 1);
    Ab = [A, b];                  % matrice aumentata (n x n+1)
    for k = 1:n-1
        % vettore dei moltiplicatori per tutte le righe sotto k
        lk = Ab(k+1:n, k) / Ab(k, k);   % vettore (n-k) x 1
        % aggiorna tutte le righe da k+1 a n in un'unica operazione
        Ab(k+1:n, k:n+1) = Ab(k+1:n, k:n+1) - lk * Ab(k, k:n+1);
        % azzera esplicitamente la colonna k (sotto il pivot)
        Ab(k+1:n, k) = 0;
    end

    U = Ab(:, 1:n);
    c = Ab(:, n+1);
end
\end{lstlisting}
La soluzione completa di $Ax = b$ si ottiene
combinando \texttt{my\_gauss} con \texttt{sup\_solve}.


% ============================================================
\section{Lezione 5 -- Sistemi con Più Termini Noti e Metodi Iterativi}
% ============================================================

\subsection{Sistemi con più termini noti}

L'equazione $AX = B$, con $B \in \mathbb{C}^{n \times s}$, è equivalente a $s$
sistemi lineari con la stessa matrice $A$. Si applica l'eliminazione di Gauss alla matrice aumentata $(A \mid B)$ e si risolvono gli $s$ sistemi triangolari
risultanti (con la stessa $U$, termini noti diversi).

\subsection{Determinante con l'eliminazione di Gauss}

Il determinante di $A$ si calcola dall'eliminazione di Gauss come prodotto dei pivot (con segno, tenendo conto degli scambi di riga):
\[
  \det(A) = \prod_{k=1}^{n} U_{kk}.
\]
\begin{nota}
    Questo deriva dal fatto che l'eliminazione di Gauss produce una fattorizzazione LU che, tramite il teorema di Binet-Cauchy sul prodotto dei determinanti, ci permette di capire da dove derivi questa semplificazione.
\end{nota}
\noindent
Questo approccio ha costo $O(n^3)$, molto più efficiente dell'espansione di Laplace
($O(n!)$).

\subsection{Metodo di Jacobi}

La formula del metodo di Jacobi per $Ax = b$ è:
\[
  x_i^{(\text{new})} = \frac{1}{A_{ii}} \left( b_i - \sum_{j \neq i} A_{ij} x_j^{(\text{old})} \right), \quad i = 1, \ldots, n.
\]
Tutte le componenti di $x^{(\text{new})}$ si calcolano usando $x^{(\text{old})}$:
\begin{lstlisting}[language=Matlab]
function x = jacobi(A, b, k)
    n = size(A, 1);
    x = b;
    for iter = 1:k
        x_new = zeros(n, 1);
        for i = 1:n
            x_new(i) = (b(i) - A(i,:)*x + A(i,i)*x(i)) / A(i,i);
        end
        x = x_new;
        norm(A*x - b)
    end
end
\end{lstlisting}

\subsection{Metodo di Gauss--Seidel}

Differisce da Jacobi per l'uso immediato dei valori aggiornati:
\[
  x_i^{(\text{new})} = \frac{1}{A_{ii}} \left( b_i
    - \sum_{j=1}^{i-1} A_{ij} x_j^{(\text{new})}
    - \sum_{j=i+1}^{n} A_{ij} x_j^{(\text{old})} \right).
\]
Spesso converge più velocemente di Jacobi, ma non sempre: la convergenza di
entrambi i metodi dipende dalle proprietà spettrali della matrice $A$.\\
Esempi.

% ============================================================
\section{Lezione 6 -- Fattorizzazione QR e Minimi Quadrati}
% ============================================================

\subsection{Trasformazioni di Householder}

Una \textbf{trasformazione di Householder} è una matrice ortogonale della forma:
\[
  H_v = I - \frac{2vv^T}{\|v\|_2^2}, \quad v \in \mathbb{R}^m.
\]
Dato $x \in \mathbb{R}^m$, scegliendo $v = x \pm \|x\|_2 e_1$ si ottiene
$H_v \cdot x = \pm\|x\|_2 e_1$ (multiplo del primo vettore della base canonica).
Per evitare errori di cancellazione nel calcolo di $v_1$, si sceglie il segno
\textbf{opposto} a quello di $x_1$:
\[
  v_1 = x_1 + \operatorname{sign}(x_1)\|x\|_2.
\]
\begin{lstlisting}
function v = householder_vector(x)
    nrm = norm(x);
    if nrm == 0
        v = 0;
    else
        v = x;
        v(1) = v(1) + sign(x(1)) * nrm;
    end
end
\end{lstlisting}

\subsection{Fattorizzazione QR}

Data $A \in \mathbb{R}^{m \times n}$ con $m \geq n$, la fattorizzazione $QR$ è
$A = QR$ dove $Q \in \mathbb{R}^{m \times m}$ è ortogonale e $R \in \mathbb{R}^{m \times n}$
è triangolare superiore.\\
Si applicano sequenzialmente $n$ trasformazioni di Householder per triangolarizzare $A$:
\[
  H_n \cdots H_1 A = R, \qquad Q = H_1^T \cdots H_n^T.
\]
\begin{lstlisting}
function [Q, R] = my_qr(A)
    %MY_QR calcola la fattorizzazione QR di A, utilizzando riflettori di Householder.
    %
    % Questa implementazione funziona anche con matrici rettangolari. 
    
    [m, n] = size(A);
    
    R = A;
    Q = eye(m);
    
    % Riduciamo tutte le colonne in modo che siano triangolari superiori.
    % Dobbiamo occuparci delle prime min(n, m - 1) colonne. 
    for j = 1 : min(m - 1, n)
        % Calcolo un riflettore di Householder che mappa la parte restante
        % della colonne j-esima di R in un multiplo di e1.
        u = householder_vector(R(j:end, j));
        b = 2/norm(u)^2;
        % Lo applico ad R a sx, e a Q a dx.
        R(j:end,j:end) = R(j:end,j:end) - b * u * (u' * R(j:end,j:end));
        Q(:,j:end) = Q(:,j:end) - b * (Q(:,j:end) * u) * u';
        
        % Dato che R e' triangolare superiore, la parte triangolare inferiore si puo' impostare a zero manualmente, per evitare di vedere del round-off
        R(j+1:end,j) = 0;
    end
end
\end{lstlisting}
Verifica della correttezza:
\begin{lstlisting}[language=Matlab]
[Q, R] = my_qr(A);
norm(A - Q*R) / norm(A)          % deve essere piccolo
norm(Q'*Q - eye(size(Q,1)))      % Q deve essere ortogonale
\end{lstlisting}

\subsection{Problemi ai minimi quadrati}

Il problema $\min_{x \in \mathbb{R}^n} \|Ax - b\|_2$ si risolve con due strategie:

\begin{enumerate}
  \item \textbf{Equazioni normali}: si risolve $A^T A x = A^T b$ con
        l'eliminazione di Gauss. Semplice ma potenzialmente instabile se
        $A$ è mal condizionata.
  \item \textbf{Metodo QR}: con $A = QR$, $R = \begin{pmatrix}R_1 \\ 0\end{pmatrix}$
        e $Q^T b = \begin{pmatrix}c_1 \\ c_2\end{pmatrix}$, si risolve il
        sistema triangolare $R_1 x = c_1$. Più stabile numericamente;
        il residuo ottimo è $\|c_2\|_2$.
\end{enumerate}
Il metodo QR è preferibile quando $A$ è quasi singolare o mal condizionata
(ad esempio quando le colonne di $A$ hanno norme molto diverse).


% ============================================================
\section{Lezione 7 -- Interpolazione Polinomiale}
% ============================================================

\subsection{Matrice di Vandermonde}

Dati $k+1$ nodi $x_0, \ldots, x_k$, la matrice di Vandermonde è:
\[
  V = \begin{pmatrix}
    1 & x_0 & x_0^2 & \cdots & x_0^k \\
    1 & x_1 & x_1^2 & \cdots & x_1^k \\
    \vdots & & & & \vdots \\
    1 & x_k & x_k^2 & \cdots & x_k^k
  \end{pmatrix}.
\]
In MATLAB:
\begin{lstlisting}[language=Matlab]
V = vander(x);
V = V(:, end:-1:1);   % ordine crescente delle potenze
\end{lstlisting}
Il numero di condizionamento di $V$ cresce \textbf{esponenzialmente} con $k$
quando i nodi sono equispaziati: questo rende il sistema mal condizionato
per gradi elevati.

\subsection{Polinomio interpolante di Lagrange}

Dati $k+1$ punti $(x_j, y_j)$, il polinomio interpolante è:
\[
  L_k(x) = \sum_{j=0}^{k} y_j \ell_j(x), \qquad
  \ell_j(x) = \prod_{\substack{i=0 \\ i \neq j}}^{k} \frac{x - x_i}{x_j - x_i}.
\]
\begin{lstlisting}
% x vettore nodi dell'interpolazione
% y valori della funzione interpolata nei nodi
% v vattore punti dove si vuole valutare la f interpolata
function w = valuta_lagrange(x, y, v)
	k = length(x);
	m = length(v);
	w = zeros(m, 1);
	% rendo i tre vettori dei vettori colonna
	x = x(:);
	y = y(:);
	v = v(:);
	for j = 1:k
	 	xx = x([1:j-1, j+1:end]); % faccio un vettore temporaneo con i nodi di interpolazione meno il j-esimo
	 	tmp = ones(m, 1); % accumulatore
	 	for h = 1:k-1
	 		tmp = tmp .* (v - xx(h)); 
	 	end
	 	tmp = tmp / prod(x(j) - xx);
	 	w = w + y(j) * tmp;
	end
end
\end{lstlisting}

\subsection{Errore di interpolazione}

L'errore commesso interpolando $f$ con un polinomio di grado $k$ è:
\[
  \max_{x \in [a,b]} |E_k(x)| \leq
  \frac{\max_{x \in [a,b]} |f^{(k+1)}(x)|}{(k+1)!}
  \cdot \max_{x \in [a,b]} \prod_{j=0}^{k} |x - x_j|.
\]

\subsection{Fenomeno di Runge}

Per la funzione $f(x) = \frac{1}{1+x^2}$ sull'intervallo $[-5,5]$,
aumentando il numero di nodi equispaziati il polinomio interpolante
\textbf{non converge} uniformemente a $f$: si osservano forti oscillazioni
agli estremi dell'intervallo al crescere del grado (fenomeno di Runge).

\begin{figure}[H]
    \centering
    \includegraphics[width=0.75\linewidth]{img/runge11-2.png}
    \caption{Fenomeno di Runge per un polinomio di interpolazione di grado 11 per la funzione $100x^2-100x+26$}
    \label{fig:runge-esempio}
\end{figure}


\subsubsection{Nodi di Chebyshev}
\label{sec: lab nodi chebyshev}

Il fenomeno di Runge si elimina usando i \textbf{nodi di Chebyshev}
sull'intervallo $[a, b]$:
\[
  x_j = \frac{a+b}{2} + \frac{b-a}{2} \cos\!\left(\frac{(2j+1)\pi}{2k}\right),
  \quad j = 0, \ldots, k-1.
\]
Questi nodi minimizzano il termine $\prod_{j=0}^{k}|x - x_j|$ e garantiscono
la convergenza dell'interpolazione al crescere di $k$.

\begin{figure}[H]
    \centering
    \includegraphics[width=0.85\linewidth]{img/runge_equi_vs_cheb.png}
    \caption{Confronto tra interpolazione su nodi equispaziati e su nodi di Chebyshev per la funzione di Runge $f(x) = \frac{1}{100x^2-100x+26}$ sull'intervallo $[0,1]$, polinomio di grado 11 (stessi nodi della fig. \ref{fig:runge-esempio}).}
\end{figure}


% ============================================================
\section{Lezione 8 -- Interpolazione a Tratti e Fitting di Dati}
% ============================================================

\subsection{Fitting ai minimi quadrati}

Dato un insieme di coppie $(x_i, f(x_i))$, si cerca una combinazione lineare
di funzioni modello $\varphi_1, \ldots, \varphi_n$ che approssimi i dati nel
senso dei minimi quadrati. La matrice del sistema è:
\[
  A_{ij} = \varphi_j(x_i).
\]
% Ad esempio, per un'approssimazione con polinomio di grado $p$:
% \begin{lstlisting}[language=Matlab]
% A = zeros(m, p+1);
% for j = 1:p+1
%     A(:,j) = x.^(j-1);
% end
% [c, res] = mq_qr(A, b);
% \end{lstlisting}

\begin{figure}[H]
    \centering
    \includegraphics[width=0.45\linewidth]{img/fitting dati.png}
    \caption{Dati sperimentali della lezione da approssimare nel senso dei minimi quadrati.}
\end{figure}

\subsection{Funzioni polinomiali a tratti}

Una funzione $F: [a,b] \to \mathbb{R}$ è \textbf{polinomiale a tratti} di grado
$\leq p$ se, su ciascun sottointervallo $[x_{i-1}, x_i]$, coincide con un
polinomio $p_i(x) = c_0^{(i)} + c_1^{(i)} x + \cdots + c_p^{(i)} x^p$.
Si rappresenta con un vettore di $k+1$ nodi e una matrice $C \in \mathbb{R}^{k \times (p+1)}$
dei coefficienti. Per valutare si usa \texttt{polyval}.

\subsection{Interpolazione lineare a tratti}

Su ciascun intervallo $[x_{i-1}, x_i]$, l'interpolante lineare a tratti è:
\[
  p_i(x) = (f(x_i) - f(x_{i-1})) \frac{x - x_{i-1}}{x_i - x_{i-1}} + f(x_{i-1}).
\]
L'errore decresce come $O(h^2)$ (con $h$ ampiezza del sottointervallo).

\subsection{Spline cubiche naturali}
\label{sec: lab spline cubiche naturali}

Le \textbf{spline cubiche naturali} sono polinomi di grado 3 a tratti che si
raccordano con continuità fino alla derivata seconda. Su nodi equispaziati
con passo $h$, i coefficienti $m_0, \ldots, m_k$ si trovano risolvendo:
\[
  \begin{pmatrix}
    2 & 1 & & \\
    1 & 4 & 1 & \\
    & \ddots & \ddots & \ddots \\
    & & 1 & 4 & 1 \\
    & & & 1 & 2
  \end{pmatrix}
  \begin{pmatrix} m_0 \\ m_1 \\ \vdots \\ m_k \end{pmatrix}
  = \frac{3}{h}
  \begin{pmatrix}
    f(x_1) - f(x_0) \\
    f(x_2) - f(x_0) \\
    \vdots \\
    f(x_k) - f(x_{k-1})
  \end{pmatrix}
\]
Le condizioni al bordo naturali impongono $p_1''(x_0) = 0$ e $p_k''(x_k) = 0$.
L'errore delle spline cubiche decresce come $O(h^4)$, molto più velocemente
dell'interpolazione lineare a tratti.

\begin{figure}[H]
    \centering
    \includegraphics[width=0.90\linewidth]{img/lineare vs spline.png}
    \caption{Confronto tra interpolazione lineare a tratti e spline cubica naturale per la funzione di Runge $f(x) = \frac{1}{100x^2-100x+26}$ sull'intervallo $[0,1]$, con 11 nodi equispaziati (stessi nodi delle fig. \ref{fig:runge-esempio} e del confronto equispaziati/Chebyshev).}
\end{figure}


% ============================================================
\section{Lezione 9 -- Integrazione Numerica}
% ============================================================

\subsection{Formule di Newton-Cotes composite}

Per approssimare $\int_a^b f(x)\,dx$ con maggiore accuratezza si suddivide
$[a,b]$ in $L$ sottointervalli di ampiezza $h = (b-a)/L$ e si applica
una formula di quadratura su ciascuno.\\
\textbf{Formula dei trapezi generalizzata}:
\[
  J_1^{(G)}(f) = \frac{b-a}{2L}
  \left[ f(x_0) + 2\sum_{i=1}^{L-1} f(x_i) + f(x_L) \right].
\]
L'errore è $O(h^2)$; esatta per polinomi di grado $\leq 1$.\\
\textbf{Formula di Cavalieri-Simpson generalizzata}:
\[
  J_2^{(G)}(f) = \frac{b-a}{6L}
  \left[ f(x_0) + 2\sum_{i=1}^{L-1} f(x_i)
         + 4\sum_{i=1}^{L} f\!\left(\tfrac{x_i + x_{i-1}}{2}\right)
         + f(x_L) \right].
\]
L'errore è $O(h^4)$; esatta per polinomi di grado $\leq 3$.\\
Implementazione della formula dei trapezi:
\begin{lstlisting}[language=Matlab]
function I = trapezi_gen(f, a, b, L)
    x = linspace(a, b, L+1);
    I = (b-a)/(2*L) * (f(x(1)) + 2*sum(f(x(2:end-1))) + f(x(end)));
end
\end{lstlisting}

\subsection{Confronto tra metodi}

A parità di $L$ sottointervalli, la formula dei trapezi richiede $L+1$ valutazioni di $f$, mentre Simpson ne richiede $2L+1$.\\
Vogliamo capire quale dei due metodi ha un costo minore (in termini di valutazioni di f (x)) se l'obbiettivo è una data
accuratezza sulla stima dell'integrale.

Simpson è nettamente più efficiente:
per raggiungere un'accuratezza $\varepsilon$, Simpson richiede molte meno valutazioni
di $f$ rispetto ai trapezi (l'errore decade come $L^{-4}$ anziché $L^{-2}$).

\begin{figure}[H]
    \centering
    \includegraphics[width=0.90\linewidth]{img/valutazioni_trapezi_simpson.png}
    \caption{Numero minimo di valutazioni di $f$ necessarie ai trapezi generalizzati e a Cavalieri-Simpson generalizzata per ottenere un errore inferiore a $\varepsilon = 10^{-1}, \ldots, 10^{-7}$, nella stima di $\int_0^3 x^2 \sin(x)^3\,dx$ (scala log-log su entrambi gli assi). Per tolleranze larghe i due metodi sono comparabili, ma al diminuire di $\varepsilon$ Simpson (pendenza $\sim \varepsilon^{-1/4}$) scala molto meglio dei trapezi (pendenza $\sim \varepsilon^{-1/2}$): per $\varepsilon = 10^{-7}$ servono 1968 valutazioni con i trapezi contro sole 113 con Simpson.}
\end{figure}


% ============================================================
\section{Lezione 10 -- Ricerca di Zeri}
% ============================================================

\subsection{Metodo del punto fisso}

Se $f(x) = 0$ si riscrive come $x = \varphi(x)$, si genera la successione
\[
  x_{n+1} = \varphi(x_n), \quad n = 0, 1, 2, \ldots
\]
La convergenza dipende dalla scelta di $\varphi$: si ha convergenza se e solo se
$|\varphi'(x^*)| < 1$ nell'intorno del punto fisso $x^*$.

\begin{lstlisting}[language=Matlab]
function xvec = puntofisso(phi, x0, tol, maxit)
    xvec = x0;
    x = x0;
    for n = 1:maxit
        xn = phi(x);
        xvec = [xvec; xn];
        if abs(xn - phi(xn)) < tol
            break;
        end
        x = xn;
    end
end
\end{lstlisting}
Per $f(x) = x - x^{1/3} - 2$ (radice unica in $[3,5]$), le tre iterazioni
$\varphi_1(x) = x^{1/3} + 2$, $\varphi_2(x) = (x-2)^3$,
$\varphi_3(x) = \frac{6 + 2x^{1/3}}{3 - x^{-2/3}}$
si comportano diversamente: $\varphi_2$ diverge ($|\varphi_2'(x^*)| \gg 1$),
$\varphi_1$ e $\varphi_3$ convergono.

\subsection{Metodo di Newton}

Il \textbf{metodo di Newton} (o delle tangenti) genera la successione:
\[
  x_{n+1} = x_n - \frac{f(x_n)}{f'(x_n)}, \quad n = 0, 1, 2, \ldots
\]
Converge \textbf{quadraticamente} nell'intorno di una radice semplice
(il numero di cifre corrette raddoppia ad ogni iterazione).

\begin{lstlisting}[language=Matlab]
function [xvec, valvec] = newton(f, df, x0, tol, maxit)
    x = x0;
    xvec = x; valvec = f(x);
    for n = 1:maxit
        x = x - f(x)/df(x);
        xvec   = [xvec;   x];
        valvec = [valvec; f(x)];
        if abs(f(x)) < tol
            break;
        end
    end
end
\end{lstlisting}
Il metodo può non convergere se il punto iniziale è lontano dalla radice
o se $f'(x_n) \approx 0$ durante le iterazioni (es.\ $f(x) = x^2 - A$ con $x_0 = 0$
porta a una divisione per zero).

\begin{figure}[H]
    \centering
    \includegraphics[width=0.85\linewidth]{img/newton_puntofisso.png}
    \caption{Errore relativo $|x_n - x^*|/|x^*|$ in funzione delle iterazioni (scala semi-logaritmica) per il metodo di punto fisso ($\varphi_1(x) = x^{1/3}+2$) e per il metodo di Newton, applicati a $f(x) = x - x^{1/3} - 2$ con $x_0 = 4$. Il punto fisso converge linearmente (retta in scala log), Newton quadraticamente (il numero di cifre corrette raddoppia a ogni passo, fino a saturare alla precisione di macchina dopo sole 3 iterazioni).}
\end{figure}


% ============================================================
\section{Lezione 11 -- Calcolo di Autovalori}
% ============================================================

\subsection{Metodo delle potenze}

Il \textbf{metodo delle potenze} calcola l'autovalore di modulo massimo $\lambda_1$
di una matrice $A$ diagonalizzabile con $|\lambda_1| > |\lambda_2| \geq \cdots \geq |\lambda_n|$.

Fissato $z^{(0)} \in \mathbb{C}^n$ con $\|z^{(0)}\|_2 = 1$, si genera:
\[
  y^{(k)} = A z^{(k-1)}, \qquad
  \lambda^{(k)} = [z^{(k-1)}]^H y^{(k)}, \qquad
  z^{(k)} = \frac{y^{(k)}}{\|y^{(k)}\|_2}.
\]
La successione $\lambda^{(k)}$ converge a $\lambda_1$; $z^{(k)}$ all'autovettore
corrispondente.

\begin{lstlisting}[language=Matlab]
function [z, lamvec] = potenze(A, z0, maxit)
    z = z0 / norm(z0);
    lamvec = [];
    for k = 1:maxit
        y = A * z;
        lam = z' * y;
        lamvec = [lamvec; lam];
        z = y / norm(y);
    end
end
\end{lstlisting}

\textbf{Attenzione}: se il vettore iniziale $z^{(0)}$ è ortogonale all'autovettore
dominante $v^{(1)}$ (ovvero $[v^{(1)}]^H z^{(0)} = 0$), il metodo potrebbe non
convergere a $\lambda_1$.

\subsection{Metodo delle potenze inverse}

Applicando il metodo delle potenze ad $A^{-1}$, si calcola l'autovalore di
modulo \textbf{minimo} $\lambda_n$ (poiché $A^{-1}$ ha autovalore dominante $1/\lambda_n$).

Poiché calcolare $A^{-1}$ è sconveniente, si sostituisce il prodotto
$y^{(k)} = A^{-1} z^{(k-1)}$ con la soluzione del sistema lineare
$A y^{(k)} = z^{(k-1)}$. La fattorizzazione LU di $A$ si calcola
una volta sola e si riutilizza ad ogni iterazione:

\begin{lstlisting}[language=Matlab]
function [z, lamvec] = potenze_inverse(A, z0, maxit)
    z = z0 / norm(z0);
    lamvec = [];
    [L, U, P] = lu(A);  % O(n^3) fatto UNA SOLA VOLTA, P permutazione
    for k = 1:maxit
        y = U \ (L \ (P * z));   % risolve A*y = z (O(n^2))
        lam = 1 / (z' * y);
        lamvec = [lamvec; lam];
        z = y / norm(y);
    end
end
\end{lstlisting}

\begin{figure}[H]
    \centering
    \includegraphics[width=0.85\linewidth]{img/potenze_due_matrici.png}
    \caption{Errore relativo $|\lambda^{(k)} - \lambda_1|/|\lambda_1|$ del metodo delle potenze in funzione delle iterazioni (scala semi-logaritmica), per due matrici simmetriche $3\times 3$ con gli stessi autovettori ma diverso rapporto spettrale: matrice A con $|\lambda_2/\lambda_1| = 0.3$ (autovalori ben separati) e matrice B con $|\lambda_2/\lambda_1| = 0.9$ (autovalori vicini). La convergenza è lineare con fattore $|\lambda_2/\lambda_1|$: la matrice A converge alla precisione di macchina in circa 15 iterazioni, mentre la B è ancora lontana dalla convergenza dopo 30 iterazioni.}
\end{figure}

\section*{Riepilogo dei contenuti}
\begin{itemize}
    \item Teoremi: \arabic{teorema}
    \item Metodi: \arabic{metodo}
    \item Definizioni: \arabic{definizione}
\end{itemize}

\chapter{MatLab utils}

\begin{utilita}
    [Forzare un vettore a una particolare forma]
    \begin{lstlisting}[language=Matlab]
v = v(:);   %forma verticale
v = v(:)';   %forma orizzontale
    \end{lstlisting}
\end{utilita}

\begin{utilita}
    [Definizione di una funzione parametrica]
    \begin{lstlisting}[language=Matlab]
f = @(t) log(abs(t.^3 + cos(t)));   % esempio
    \end{lstlisting}
\end{utilita}

\begin{utilita}
    [Somma righe/colonne matrice]
    \begin{lstlisting}[language=Matlab]
r = sum(B, 1);  % somma lungo le righe, r = sum(col)
r = sum(B, 2);  % somma lungo le colonne, r = sum(righe)
    \end{lstlisting}
\end{utilita}

\begin{utilita}
    [Matrice booleana]
    \begin{lstlisting}[language=Matlab]
M = A <= R;  % esempio, ritorna M con 1 dove la condizione e' vera
    \end{lstlisting}
\end{utilita}

\begin{utilita}
    [Broadcasting esponenziale]
    \begin{lstlisting}[language=Matlab]
        x = x(:);  % vettore colonne, e.g. n x 1
        j = 0:k;   % vettore riga esponenti, 1 x (k+1)
        V = x.^j;  % broadcasting: n x (k+1)
    \end{lstlisting}
    Ogni elemento del vettore base viene elevato per ognuna delle componenti del vettore esponente.
\end{utilita}


\end{document}
